% ------------------------------------------------------------------
% Distinction, Admissibility, Repair, and Reconstruction
% Combined volume — merged from eight standalone chapter drafts.
% Compile with XeLaTeX (two passes for cross-references).
% Author byline: Flyxion. Independent Researcher.
%
% Chapter order: Distinction, Information, Constraint Populations,
% Admissibility, Entropy, Repair, Restorability Boundary, Continuation,
% Reconstruction.
% ------------------------------------------------------------------

\documentclass[12pt]{report}

\usepackage{fontspec}
\setmainfont{TeX Gyre Pagella}

\usepackage{amsmath}
\usepackage{amssymb}
\usepackage{amsthm}
\usepackage{caption}
\usepackage[margin=1.35in]{geometry}
% \usepackage{microtype}
\usepackage[colorlinks=true,linkcolor=black,citecolor=black,urlcolor=black]{hyperref}
\usepackage{cleveref}

\newtheorem{definition}{Definition}[section]
\newtheorem{proposition}[definition]{Proposition}
\newtheorem{theorem}[definition]{Theorem}
\newtheorem{corollary}[definition]{Corollary}
\newtheorem{lemma}[definition]{Lemma}
\newtheorem{claim}[definition]{Claim}
\newtheorem{example}[definition]{Example}
\newtheorem{openproblem}[definition]{Open Problem}
\newtheorem{remark}[definition]{Remark}
\newtheorem{principle}[definition]{Principle}
\crefname{claim}{Claim}{Claims}
\Crefname{claim}{Claim}{Claims}

% Macros used across chapters (each source file defined these
% identically; consolidated here to a single definition apiece).
\newcommand{\Reach}{\ensuremath{\operatorname{Reach}}}
\newcommand{\Rboundary}[1][]{\ensuremath{\partial_{\mathrm{R}#1}}}
\newcommand{\I}{\ensuremath{\mathrm{I}}}
\newcommand{\Ddelta}{\ensuremath{\Delta\delta}}
\newcommand{\Dstar}{\ensuremath{D^{*}}}
\newcommand{\DT}{\ensuremath{D_T}}
\newcommand{\Lreach}{\ensuremath{L_T}}
\newcommand{\Lopt}{\ensuremath{L^{*}}}
\newcommand{\Health}{\ensuremath{\mathsf{H}}}
\newcommand{\SR}{\ensuremath{S_R}}

\title{Restorability Boundaries}
\author{Flyxion \\ \small Independent Researcher}
\date{}

\begin{document}

\maketitle

\begin{abstract}
\noindent
This volume develops a framework linking distinction, information,
entropy, admissibility, repair, restorability, continuation, and
reconstruction. Distinctions induce the partitions on which information
and entropy are defined. Admissibility, given two independent
characterizations that are shown to agree in type but not yet in
content, governs which transitions between states are reachable at all.
Repair, and the operators that persist as ecologies rather than isolated
fixes, keep admissible structure viable under damage. Restorability
boundaries mark where that viability fails, by four formally distinct
mechanisms rather than one. Continuation resolves an apparent
contradiction between admissibility as checked and admissibility as
authored by showing each describes a different level of the same
structure. Reconstruction closes the framework with a resource-independence
result: some collapsed distinctions are not merely expensive to recover
but structurally unrecoverable by any allocation of time, space, or
energy. The volume's central open problem --- whether its two
admissibility constructions induce the same continuation family, a
sharper question than whether they are equal as quantities --- is stated
precisely, given a formal argument for skepticism, and left open.
\end{abstract}

\chapter*{Preface}
\addcontentsline{toc}{chapter}{Preface}

Persistence depends on the preservation and restoration of distinctions.
That is the central claim of this volume, and the chapters that follow
develop it through a sequence of increasingly structured constructions.
Distinguishability induces information; information induces entropy;
entropy constrains admissibility; admissibility supports repair; repair
determines restorability; restorability governs continuation; and
reconstruction studies what survives, and what does not, when
distinctions are lost.

The mathematics is the argument. Distinguishability spaces, coding
deficits, and the four operations of Chapter 1 are not preliminaries to
a later, looser discussion --- they are the objects the rest of the
volume is built from, checked at each stage by substitution rather than
assumed by analogy. When two constructions share a name or a formula,
this volume treats that as a question to be settled, not a license to
identify them: a type test first, then a test of whether one is
literally a special case of the other. Most of the time, the answer
turns out to be no. The volume's single largest open problem --- whether
two independently-motivated notions of admissibility induce the same
family of permitted continuations --- is a direct consequence of taking
that discipline seriously rather than resolving it by convenience.

Four appendices collect material that would otherwise interrupt the
chapters that need it: a notation index, the shared machinery of
quotients and fibers underlying Chapters 1, 2, and 6, the continuation
and reachability structures used throughout the second half of the
volume, and the persistence and reconstruction results that close it.

\tableofcontents

\chapter{Distinction}
% ------------------------------------------------------------------

\section{Introduction}

Three independent lines of work, developed for unrelated purposes,
converge on the same mathematical object once their concerns are stated
precisely enough to compare. A line of work proceeding from
process-primary computation establishes that an observable state is a
projection of a richer history space, and that distinct histories
routinely collapse to the same observable state under that projection. A
line of work proceeding from kernel statistics demonstrates that distinct
continuous distributions become mutually singular after an appropriate
embedding, separated in the embedding space even where they overlapped
substantially before it. A line of work proceeding from archive-based
compression identifies a notion of reachable complexity: the shortest
description an archive can currently produce, which may be strictly
longer than the shortest description that exists in principle. The first
of these is concerned with distinctions that are destroyed. The second is
concerned with distinctions that are amplified. The third is concerned
with distinctions that are, for now, simply inaccessible. None of the
three names the object they share. This chapter makes it explicit.

\section{The Distinguishability Space}
\label{sec:space}

\begin{definition}[Distinguishability space]
\label{def:dist-space}
A \emph{distinguishability space} is a pair $(X,\sim)$ where $X$ is a set
and $\sim$ is an equivalence relation on $X$, interpreted as: $x \sim y$
exactly when $x$ and $y$ cannot be told apart within the current
representational system. The equivalence classes $[x]_\sim$ are its
\emph{fibers}.
\end{definition}

A finer relation has smaller fibers and expresses more distinctions; a
coarser relation has larger fibers and expresses fewer. The two extremes
are the discrete relation, under which every element is distinguishable
from every other, and the indiscrete relation, under which nothing is
distinguishable from anything. The canonical instance recurring
throughout this book is a projection $\sigma : H \to S$ from a history
space $H$ to an observable state space $S$, where $h \sim h'$ exactly
when $\sigma(h) = \sigma(h')$: the fiber above a state $s$ is the set of
every history that would have produced it.

\begin{principle}[Distinguishability Principle]
\label{prin:dist}
A representational system is characterized not by what objects it
contains but by what distinctions it can express.
\end{principle}

An immediate consequence of \cref{prin:dist} is that no finitary
representational system can express every distinction present in what it
represents, so every such system carries some deficit relative to what a
maximally expressive system could achieve. Making that deficit precise is
this chapter's first formal task.

\section{Two Deficits}
\label{sec:deficits}

Let $T$ denote an ontology in the sense used throughout this book: a
description language, a template hierarchy, or a representational system
generating a distinguishability relation on some set of objects.

\begin{definition}[Coding deficit]
\label{def:coding-deficit}
Let $\Lreach(x)$ be the length of the shortest description of $x$
reachable within $T$, and let $\Lopt(x)$ be the shortest description of
$x$ achievable by any ontology. The \emph{coding deficit} is
$\delta_T(x) = \Lreach(x) - \Lopt(x)$.
\end{definition}

\begin{definition}[Distinguishability deficit]
\label{def:dist-deficit}
Let $\DT(x)$ be the number of distinct objects separable from $x$ by
descriptions reachable within $T$, and let $\Dstar(x)$ be the maximum
such capacity achievable by any ontology. The \emph{distinguishability
deficit} is $\Delta_T(x) = \Dstar(x) - \DT(x)$.
\end{definition}

These measure different things: $\delta_T(x)$ is a description-length
excess, how many bits are wasted; $\Delta_T(x)$ is an inaccessible
distinction capacity, how many distinctions involving $x$ cannot be
expressed in $T$ at all. They are related rather than interchangeable.

\begin{claim}[Deficit Correspondence]
\label{thm:deficit-correspondence}
Under the assumption that descriptions and distinctions are related by a
faithful encoding, $\Delta_T(x) \leq C \cdot \delta_T(x)$ for a
system-dependent constant $C > 0$, and $\delta_T(x) = 0$ if and only if
$\Delta_T(x) = 0$: local optimality in coding corresponds exactly to
exhausting every expressible distinction.
\end{claim}

Both deficits are non-negative for every ontology and every object, and
both vanish exactly when the ontology is locally optimal at that object.
Throughout the rest of this book, $\delta_T$ is used where the argument
concerns description length and $\Delta_T$ where the argument concerns
fiber structure directly; \cref{thm:deficit-correspondence} licenses
treating a claim proved for one as evidence, though not proof, for the
corresponding claim about the other.

\section{The Four Operations}
\label{sec:operations}

Representational systems do not stay fixed. \cref{prin:dist} raises an
immediate question: what are the possible ways a distinguishability
structure can change? Four operations answer it, organized by a
two-level architecture.

\[
\text{Objects} \;\longleftrightarrow\; \text{Distinctions}
\;\longleftrightarrow\; \text{Ontologies}.
\]

Projection and embedding operate at the level of \emph{distinctions}:
they modify $\sim$ within a single fixed space, moving between finer and
coarser relations on the same underlying set. Revision and transport
operate at the level of \emph{ontologies}: revision changes the
generating ontology of a space outright, and transport connects two
distinct spaces to one another.

\begin{definition}[Projection]
\label{def:projection}
A \emph{projection} on $(X,\sim)$ is a surjective map $\sigma : X \to S$
inducing a coarsening $\sim'$ of $\sim$: $\sigma(x) = \sigma(y)$ implies
$x \sim' y$.
\end{definition}

\begin{claim}[Projection Deficit Bound]
\label{thm:proj-deficit}
For a projection $\sigma : X \to S$ with fiber $F_s = \sigma^{-1}(s)$
above a state $s$ with $|F_s| > 1$, the deficit increase satisfies
$\Ddelta(\sigma) \geq \log|F_s|$: projecting to $S$ discards at least
$\log|F_s|$ bits of distinguishing information, transferred from the
realized history into the now-inaccessible remainder rather than
destroyed outright.
\end{claim}

\begin{definition}[Embedding]
\label{def:embedding}
An \emph{embedding} of $(X,\sim)$ into $(Y,\sim')$ is an injective map
$\phi : X \to Y$ such that $\phi(x) \sim' \phi(y)$ implies $x \sim y$. It
is \emph{faithful} when the converse holds as well.
\end{definition}

\begin{claim}[Embedding Refinement]
\label{thm:embed-refine}
A faithful embedding never increases the coding deficit:
$\delta_{\phi(X)}(\phi(x)) \leq \delta_X(x)$. A non-faithful embedding may
introduce spurious distinctions --- elements separated in the target space
that were equivalent in the source --- and the refinement guarantee fails
for it in general.
\end{claim}

\begin{definition}[Revision]
\label{def:revision}
An \emph{ontological revision} is a transition $(X,\sim_T) \to
(X,\sim_{T'})$ replacing an ontology $T$ with a new ontology $T'$,
changing which descriptions are reachable rather than merely reorganizing
$\sim$ within a fixed $T$. A revision is \emph{productive} when
$\delta_{T'}(x) < \delta_T(x)$.
\end{definition}

\begin{claim}[Revision Monotonicity]
\label{thm:revision-mono}
For a sequence of ontologies $T_0 \subset T_1 \subset T_2 \subset \cdots$
each extending the last, the coding deficit is non-increasing at every
step, strictly decreasing wherever the revision is productive. The
canonical failure mode is \emph{ontological lock-in}: the deficit is
positive and stable, no productive revision is available, and the system
cannot escape its current template hierarchy because the path to a
better ontology does not lie within the admissibility structure of the
current one.
\end{claim}

\begin{definition}[Transport]
\label{def:transport}
A \emph{transport map} is a function $\tau : (X,\sim_X) \to (Y,\sim_Y)$
with bounded distortion $\epsilon(\tau)$, preserving equivalence in the
forward direction and controlling the measure of pairs that collapse
together in $Y$ without having been equivalent in $X$.
\end{definition}

\begin{claim}[Transport Stability]
\label{thm:transport-stability}
For a transport map $\tau$ with distortion at most $\epsilon$ and a
fiber-constant invariant $I$, $|I(x) - I(\tau(x))| \leq C\epsilon$ for a
constant $C$ depending on the Lipschitz properties of $I$. An analogy
$\tau$ is justified for inference exactly when its distortion is
controlled and the relevant invariants have bounded Lipschitz constants
--- the precise condition under which reasoning by analogy is valid
rather than merely suggestive.
\end{claim}

\begin{claim}[Representational Reduction]
\label{prop:rep-reduction}
Every representational transformation induces a transformation on
distinguishability structure, classified by its effect on fibers: a
transformation that merges fibers is a projection, one that separates
fibers is an embedding, one that changes the generating ontology is a
revision, and one that connects distinct spaces is a transport. These
four operations are not imposed as a taxonomy; they are the exhaustive
qualitative possibilities for fiber behavior.
\end{claim}

\begin{openproblem}[Classification]
\label{conj:classification}
It remains open whether every representational transformation factors
into a finite composition of projection, embedding, revision, and
transport, and, if so, whether a single expression bounds how the coding
deficit accumulates across an arbitrary such composition.
\end{openproblem}

\section{Toward Admissibility}
\label{sec:toward-admissibility}

The deficit apparatus of \cref{sec:deficits} connects directly to the
admissibility geometry used throughout the rest of this book, though the
connection needs to be stated carefully rather than assumed.

\begin{claim}[Admissibility via deficit]
\label{claim:admissibility-deficit}
An admissibility relation $\mathcal{A}$ specifies which transitions
between states are reachable from a given state. States outside
$\mathcal{A}$ from an ontology $T$ are exactly those whose description
would require $\delta_T(x) > 0$ from within $T$, and whose distinction
structure is correspondingly inaccessible: $\Delta_T(x) > 0$.
\end{claim}

This is a graded characterization: $\delta_T(x)$ and $\Delta_T(x)$ take
real, non-negative values, and admissibility failure is a matter of
degree rather than a single threshold crossed once and for all.

\begin{remark}[An open connection, not an identity]
\label{rem:gate-connection}
Chapter 6 uses a different and cruder formalization: an architecture gate
$\chi_\pi(x_{\mathrm{obs}}) \in \{0,1\}$, testing simple set membership in
$\operatorname{Im}(\pi)$ for a fixed prediction map $\pi : \mathcal{M} \to
\Sigma$, with no graded notion of degree once outside the image. This
chapter's $\delta_T$/$\Delta_T$ machinery is a natural candidate for
refining that binary gate into something continuous, and a partial
bridge --- though not a full identification --- can now be given
explicitly.
\end{remark}

\begin{definition}[Architectural representation deficit]
\label{def:arch-deficit}
For a metric observation space $(\Sigma, d_\Sigma)$ and prediction map
$\pi : \mathcal{M} \to \Sigma$, define
\[
\delta_\pi^{\mathrm{arch}}(x) \;=\; \inf_{m \in \mathcal{M}} d_\Sigma(\pi(m), x).
\]
\end{definition}

\begin{claim}[Gate as thresholded architectural deficit]
\label{thm:gate-threshold}
If $\operatorname{Im}(\pi)$ is closed, then
\[
\chi_\pi(x) \;=\; \mathbf{1}_{\{\delta_\pi^{\mathrm{arch}}(x) > 0\}}.
\]
\end{claim}

\begin{proof}
If $x \in \operatorname{Im}(\pi)$, some $m$ satisfies $\pi(m) = x$, so the
infimum is zero. Conversely, if the infimum is zero and
$\operatorname{Im}(\pi)$ is closed, then $x$ lies in $\operatorname{Im}(\pi)$.
The stated equality follows from the definition of $\chi_\pi$.
\end{proof}

This is a real result, but a narrower one than it may first appear:
\cref{thm:gate-threshold} shows the binary gate is the zero/nonzero
threshold of a graded quantity, $\delta_\pi^{\mathrm{arch}}$, built from
the same kind of infimum-distance construction used throughout this
chapter. It does \emph{not} identify $\delta_\pi^{\mathrm{arch}}$ with
this chapter's coding deficit $\delta_T$. A full compatibility theorem
would require a map $J : T \mapsto \pi_T$ from ontologies to prediction
maps and constants $c_1, c_2 > 0$ such that
\[
c_1 \, \delta_{\pi_T}^{\mathrm{arch}}(x) \;\leq\; \delta_T(x) \;\leq\; c_2 \, \delta_{\pi_T}^{\mathrm{arch}}(x),
\]
or showing it fails, is left as an explicit open problem rather than
assumed by continuing to use both vocabularies side by side in later
chapters.

\begin{remark}[Preserving distinctions is not preserving competence]
\label{rem:preservation-not-continuity}
Nothing developed so far should suggest that a distinction-preserving
transformation leaves a system's operational capacity unchanged. A
bijective relabeling of $X$ --- not a projection, not a quotient, not a
fiber collapse, and therefore invisible to $\delta_T$, $\Delta_T$, and
every entropy measure built from them in later chapters --- can still
destroy the alignment between a learned operator repertoire and the
space that repertoire was trained to act on. Representational
preservation and operational continuity are independent properties: a
transformation can hold the first perfectly while destroying the second
completely.\footnote{A mirror keyboard --- a piano relabeled so pitch
runs high to low rather than low to high --- is the cleanest instance:
every distinction in the note space survives exactly, yet a trained
pianist's fingering competence can be rendered nearly unusable.} This
distinction reappears, worked through in full, in Chapter 5.
\end{remark}

\section*{Summary}

This chapter develops distinctions and the two deficits they support,
along with the four operations that act on distinguishability structure
and a first, partial connection to admissibility. It does not develop
the further step of showing how distinctions induce the partitions that
information theory operates over, nor state an entropy functional on the
resulting distinguishability space. The connection is real --- a
distinguishability relation $\sim$ induces a partition of $X$, and a
partition supports an entropy measure in the ordinary way --- and
connecting the resulting entropy back to the coding deficit of
\cref{def:coding-deficit}, rather than treating them as two unrelated
quantities that happen to share units, is developed in Chapter 2.

\chapter{Information}
% ------------------------------------------------------------------

\section{Introduction}

Chapter 1 closed by naming two things it had not yet done. It had not
shown how the distinguishability relation $\sim$ on a space $(X,\sim)$
induces the partition structure that information theory presupposes, and
it had not connected the resulting information measures back to the
coding deficit $\delta_T$ of \cref{def:coding-deficit-recall-info} rather than
treating them as two quantities that happen to be measured in the same
units. This chapter does both, and is checked against Chapter 1's actual
theorems rather than merely gestured at them.

\begin{definition}[Coding deficit]
\label{def:coding-deficit-recall-info}
For an ontology $T$ and object $x$, the coding deficit is
$\delta_T(x) = L_T(x) - L^*(x)$, the excess of the shortest description
of $x$ reachable within $T$ over the shortest description achievable by
any ontology.
\end{definition}

\section{Partitions and Surprisal}
\label{sec:surprisal}

A distinguishability relation $\sim$ on $X$ is, by
\cref{def:dist-space-recall}, equivalently a partition of $X$ into fibers.

\begin{definition}[Distinguishability space]
\label{def:dist-space-recall}
A distinguishability space is a pair $(X,\sim)$ where $\sim$ is an
equivalence relation on $X$, interpreted as observational
indistinguishability; its equivalence classes are fibers.
\end{definition}

Suppose an observer knows only that a system lies somewhere within a
partition $\mathcal{P} = \{P_1, \ldots, P_n\}$ induced by $\sim$, and an
observation identifies one cell $P_i$ with probability $p_i$. The
informational content of that event is the Shannon surprisal
$I(P_i) = -\log p_i$. Read distinction-theoretically rather than
purely statistically, this is not an imported foreign quantity: it
measures how strongly the observation refines the distinction structure
already given by $\mathcal{P}$. An observation that identifies a rare
cell refines more than an observation that identifies a common one,
exactly as a finer $\sim$ expresses more distinctions than a coarser one.
The Shannon entropy of the partition, $H(\mathcal{P}) =
-\sum_i p_i \log p_i$, is the expected surprisal --- the expected
refinement an observation of this system delivers.

\section{The Compression Theorem}
\label{sec:compression}

\begin{claim}[Compression Theorem]
\label{thm:compression}
For a surjective projection $\pi : X \to Y$, $H(Y) \leq H(X)$.
\end{claim}

The argument is immediate given Chapter 1's own Fiber Monotonicity result:
a surjective map merges fibers, the induced partition on $Y$ is coarser
than the one on $X$, and a coarser partition can express no more
distinctions than a finer one. Entropy inherits the same monotonicity
that fiber containment already established. Every representation
compresses; every model compresses; every theory compresses. The
question this chapter inherits from Chapter 1 is never whether
compression occurs but which distinctions survive it.

\section{Representational Entropy}
\label{sec:rep-entropy}

Compression's effect on a single point, rather than on a whole partition,
is worth isolating as its own quantity.

\begin{definition}[Representational entropy]
\label{def:rep-entropy}
Let $\pi : X \to M$ be a representation. The representational entropy at
$m \in M$ is $S_\pi(m) = \log|\pi^{-1}(m)|$, the log-size of the fiber
collapsed beneath $m$.
\end{definition}

Large $S_\pi(m)$ indicates many underlying states rendered
indistinguishable by the representation: high compression at $m$, with a
correspondingly large blind spot.

\begin{claim}[Representational entropy bounds the coding-deficit increase]
\label{claim:bridge}
For a projection $\sigma : X \to S$ with fiber $F_s = \sigma^{-1}(s)$,
Chapter 1's Projection Deficit Bound states $\Ddelta(\sigma) \geq
\log|F_s|$. Since $\log|F_s| = S_\sigma(s)$ by \cref{def:rep-entropy}
applied to the same fiber, this is exactly
\[
\Ddelta(\sigma) \;\geq\; S_\sigma(s).
\]
\end{claim}

This is the bridge Chapter 1 was owed, and it is a checked identity of
quantities rather than an assumed correspondence: both
$\log|F_s|$ in Chapter 1's bound and $S_\pi(m)$ in
\cref{def:rep-entropy} are defined as the log-cardinality of the same
fiber, so \cref{claim:bridge} follows by substitution, not by analogy.
It licenses treating representational entropy as a lower bound on how
much coding deficit a projection is forced to introduce, and treating a
large deficit increase under a specific projection as evidence, though
not proof, of large representational entropy at the corresponding point.

\begin{remark}[What this bridge does not yet establish]
\label{rem:bridge-limits}
\cref{claim:bridge} relates $\delta_T$'s behavior under one specific
operation, projection, to $S_\pi$ at a single point. It does not relate
the two deficits of Chapter 1 --- $\delta_T$ and $\Delta_T$ --- to
$H(\mathcal{P})$ in general, and it says nothing yet about revision or
transport, the two ontology-level operations of Chapter 1. Those
connections, if they exist, are not established here.
\end{remark}

\section{The Representation Limit Theorem}
\label{sec:rep-limit}

\begin{claim}[Representation Limit Theorem]
\label{thm:rep-limit}
No finite representation can preserve every distinction of an infinite
domain: for $|X| = \infty$ and $|M| < \infty$, any $\pi : X \to M$ assigns
infinitely many elements of $X$ to at least one fiber, by the pigeonhole
principle, and so collapses infinitely many distinctions.
\end{claim}

This is the first formal appearance of what recurs throughout this book
under the name projection failure. Loss under a bounded representation of
an unbounded domain is not an implementation defect to be engineered
away; it is a mathematical necessity following directly from cardinality
alone, independent of how cleverly the representation is designed.

\section{The Fundamental Information Theorem}
\label{sec:fund-info}

\begin{claim}[Fundamental Information Theorem]
\label{thm:fund-info}
Every informational quantity can be expressed as a property of a
distinction structure: information requires distinguishable alternatives,
which define a partition; probability distributions are defined over
partition cells; entropy is computed from those probabilities; mutual
information measures shared refinement between partitions; compression
measures distinction collapse; channel capacity measures transmissible
distinctions. Each depends solely on partition structure.
\end{claim}

\begin{remark}[The inverted hierarchy]
\label{rem:inverted}
The traditional hierarchy places objects first and derives information
from them. \cref{thm:fund-info} inverts this: distinctions precede
information, and information precedes objects. An object, on this
reading, is an informational artifact generated by a
distinction-preserving compression --- not a primitive that information
is subsequently extracted from.
\end{remark}

\section*{Summary}

This chapter develops information as a Shannon-style and representational
quantity: surprisal, partition entropy, and representational entropy,
all grounded directly in the distinguishability apparatus of Chapter 1
and checked against it rather than merely motivated by analogy to it. It
does not touch the entropy used to define admissibility: the monotonic,
thermodynamic-style field $S$ satisfying $\dot{S} \geq 0$ along admissible
trajectories in the RSVP formalization. Whether that $S$ coincides with,
is derived from, or is genuinely independent of the $H(\mathcal{P})$ and
$S_\pi(m)$ developed here is addressed in Chapter 4, which also checks
whether ``entropy'' is being used consistently across every context it
appears in, rather than assuming any two uses coincide because they
share a word.

This chapter has nothing to say about a chapter titled ``Constraint'':
no independent technical substrate for such a chapter, distinct from
admissibility, has been found. Admissibility is the working technical
term throughout the relevant material; ``constraint'' is never
operationalized there as a separate object. That chapter is accordingly
absent from this volume.

\chapter{Admissibility}
% ------------------------------------------------------------------

\section{Two Definitions}
\label{sec:two-definitions}

An admissibility relation specifies which transitions between states
are reachable from a given state. Two independent constructions give
this relation formal content.

\begin{definition}[Coding deficit]
\label{def:coding-deficit-recall-adm}
For an ontology $T$ and object $x$, the coding deficit is
$\delta_T(x) = L_T(x) - L^*(x)$, the excess of the shortest description
of $x$ reachable within $T$ over the shortest description achievable by
any ontology.
\end{definition}

\begin{definition}[Deficit-based admissibility]
\label{def:deficit-admissibility}
For an ontology $T$, a state $x$ is inadmissible from $T$ exactly when
$\delta_T(x) > 0$.
\end{definition}

\begin{definition}[Entropy-monotonicity admissibility]
\label{def:entropy-admissibility}
For a field triple $(\Phi, \mathbf{v}, S)$ over configuration space
$\mathcal{F}$, a configuration $\gamma_2$ is admissible from $\gamma_1$
when there is a continuous path in $\mathcal{F}$ connecting them along
which the entropy density $S$ is non-decreasing, $\dot{S} \geq 0$, and
regularity holds. This is a directed reachability relation, a partial
order compatible with entropy monotonicity.
\end{definition}

Both definitions describe the same underlying idea: which states can be
reached from which. That shared ambition does not by itself make them
the same relation.

\section{Testing the Identification}
\label{sec:testing}

\begin{claim}[Type compatibility]
\label{claim:type-test}
\cref{def:deficit-admissibility,def:entropy-admissibility} pass the type
test: both are relations on pairs of states, or equivalently on paths
between them, not objects of different kinds.
\end{claim}

Passing the type test is necessary but not sufficient. The stronger
requirement is a definition-by-substitution test: can one definition be
obtained from the other by specializing variables, the way historical
entropy is representational entropy evaluated at a particular
projection.

A schema general enough to compare the two precisely treats admissibility
as membership in a continuation family.

\begin{definition}[Continuation family and trajectory-relative admissibility]
\label{def:continuation-family}
A continuation family assigns to each state $x$ a set $C(x)$ of
trajectories beginning at $x$ considered permitted. A state $x_i$
occurring within a trajectory $\tau$ is admissible relative to $\tau$
and $C$ when the suffix of $\tau$ beginning at $x_i$ is a member of
$C(x_i)$.
\end{definition}

\begin{proposition}[Both definitions are continuation families]
\label{rem:schema-unification}
$C_S(x) = \{\tau : \dot{S}(\tau) \geq 0\}$ is the continuation family
induced by \cref{def:entropy-admissibility}: checking membership requires
nothing beyond evaluating a static condition at each instant.
$C_\delta(x) = \{\tau : \delta_T(\tau) \leq \epsilon\}$, for a threshold
$\epsilon$, is the continuation family induced by
\cref{def:deficit-admissibility}, built from a distortion bound rather
than an entropy condition. Recognizing both as instances of
\cref{def:continuation-family} does not identify them --- different
constructions of $C(x)$ remain different constructions --- but replaces
a comparison between two differently-typed-looking objects with a
question about two set-valued constructions of the same type.
\end{proposition}

\begin{openproblem}[Central open problem]
\label{prob:admissibility-identity}
The relevant question is not whether $\delta_T(x) = 0$ is equivalent to
the existence of an entropy-nondecreasing path to $x$, stated bare. It
is whether the induced continuation families coincide, nest, or merely
intersect:
\[
C_\delta(x) \overset{?}{=} C_S(x), \qquad C_\delta(x) \overset{?}{\subseteq} C_S(x), \qquad C_\delta(x) \cap C_S(x) \overset{?}{\neq} \varnothing.
\]
This is sharper and more tractable than a bare quantity comparison,
since it asks about set relations between two well-typed objects rather
than an equality between two differently-motivated scalars. Every later
use of admissibility in this volume --- the admissible future region
$\Gamma_t(x)$, the architecture gate $\chi_\pi$, and any treatment of
continuation or coordination --- inherits whichever answer this
eventually receives. One structural correspondence is worth noting
without being upgraded to a proof: the Second Law's $\dot{D} \leq 0$ and
the entropy condition's $\dot{S} \geq 0$ are the same monotonic
inequality under a sign flip, suggestive of $S$ playing the role of
consumed distinction capacity, which would suggest $C_\delta(x) \subseteq
C_S(x)$ if it held. Suggestive is the correct word; nothing here
upgrades it to a proof.
\end{openproblem}

\begin{remark}[On resolving this prematurely]
\label{rem:no-rush}
A framework this dependent on admissibility has an obvious incentive to
declare \cref{prob:admissibility-identity} resolved as soon as a
plausible bridge appears. That incentive is a reason for caution, not
license. Resolving it requires an explicit map with checked
functoriality between $C_\delta$ and $C_S$, not the convenience of
having many prior results click into a single coherent structure once
it is assumed.
\end{remark}

\subsection{A Structural Reason for Skepticism}
\label{subsec:structural-skepticism}

A general fact about conserved quantities gives this skepticism a
formal basis, together with a caveat about exactly how far it reaches.

\begin{theorem}[Constants of motion are generically non-injective on trajectory space]
\label{claim:conservation-noninjective}
Let $F$ be a constant of motion along a trajectory ($\dot{F} = 0$
identically) and let $\mathcal{T}$ be the space of dynamically
realizable trajectories. The level sets of $F$ generically contain more
than one distinct trajectory: except in fully integrable systems with as
many independent constants of motion as degrees of freedom, knowing the
value of $F$ alone does not determine which trajectory, among those
sharing that value, is realized.
\end{theorem}

This is close to a restatement of ordinary partial integrability, but
its consequence here is not incidental. A history channel earns its
keep, in the sense $\delta_T$ and the distinguishability apparatus of
Chapter 1 both depend on, by being \emph{injective} over a fiber ---
separating predecessors that would otherwise be conflated. A constant of
motion is, by construction, the opposite: it is valuable to physics
precisely because it is shared by many different trajectories, which is
what makes it a useful invariant rather than a complete description of
the motion. Distinguishing and conserving are not two routes to the
same destination; they are built to do opposite jobs.

\begin{remark}[Scope of the theorem]
\label{rem:monotonic-not-conserved}
\cref{claim:conservation-noninjective} is stated for a strict constant of
motion, $\dot{F} = 0$. The entropy condition entering
\cref{def:entropy-admissibility} is monotonic, $\dot{S} \geq 0$, not
conserved, and the theorem as proved does not directly cover it. The
structural point plausibly extends --- fixing ``$S$ is non-decreasing''
as a constraint on a trajectory is a coarse condition compatible with a
large multiplicity of distinct trajectories, in the same spirit as a
fixed level set --- but that extension is not proved here. What can be
said without overreach is that $S$ belongs to the family of quantities
\cref{claim:conservation-noninjective} concerns --- built to be shared
across trajectories, not to separate them --- and that this family has
been shown, at least in the conserved case, to be structurally unsuited
to the job $\delta_T$ does.
\end{remark}

A related question in classical mechanics sharpens this further. A
first pass might claim that no treatment of classical mechanics asks
anything resembling a distinguishability question; constrained
Hamiltonian systems are the counterexample, since a singular Legendre
transform there produces genuinely non-injective structure, formalized
in Dirac--Bergmann constraint theory. Physics's own unresolved
\emph{problem of time} in general relativity is, in substance, the
question of whether a distinction collapsed by a gauge constraint
carries physical content, and the literature has not settled it after
decades of direct engagement. The one identified formal difference is
that the standard physics criterion (the Dirac-observable condition) is
context-free and algebraic, applied once and for all, while
admissibility-relevance here is explicitly relative to a continuation
family, so that the same fiber can matter under one continuation family
and not another.

\begin{remark}[A speculative reframing]
\label{rem:speculative-reframing}
It is tempting, offered here only as speculation, to wonder whether
\cref{prob:admissibility-identity} has resisted resolution for a
structurally similar reason. If $\delta_T$'s admissibility is inherently
relative to an ontology and a continuation family, while the entropy
condition is posed, at least on its surface, as a context-free field
condition, then demanding an equivalence between them may be attempting
the same kind of category error the problem of time keeps encountering:
seeking a single, context-free answer to a question that may only admit
context-relative ones. This reframes why \cref{prob:admissibility-identity}
is hard without resolving it, and carries no more authority here than
the underlying dispute in the foundations of general relativity carries
elsewhere. A reframing of a difficulty is useful even when it settles
nothing.
\end{remark}

\section{The Architecture Gate as a Special Case}
\label{sec:gate-recall}

A partial, weaker bridge survives independently of
\cref{prob:admissibility-identity}'s resolution.

\begin{proposition}[Gate as thresholded architectural deficit]
\label{claim:gate-recall}
For a prediction map $\pi : \mathcal{M} \to \Sigma$ with closed image, the
binary architecture gate satisfies $\chi_\pi(x) =
\mathbf{1}_{\{\delta_\pi^{\mathrm{arch}}(x) > 0\}}$ for the architectural
representation deficit $\delta_\pi^{\mathrm{arch}}(x) = \inf_{m \in
\mathcal{M}} d_\Sigma(\pi(m),x)$.
\end{proposition}

This shows $\chi_\pi$ is a coarsening of a graded quantity built the same
way $\delta_T$ is built --- an infimum-distance construction --- without
showing $\delta_\pi^{\mathrm{arch}}$ and $\delta_T$ are the same
quantity. \cref{claim:gate-recall} and \cref{def:deficit-admissibility}
are siblings under the same construction pattern, not proven identical.
Admissibility notions throughout this volume are repeatedly built as
either a deficit (an infimum distance from a target set) or a
monotonicity condition (a non-decreasing quantity along a path), and it
remains an open question whether these are two families or one family
with two descriptions.

\section{Continuation Needs No New Machinery}
\label{sec:continuation-reduction}

The clearest positive result available is not an identification of
\cref{def:deficit-admissibility,def:entropy-admissibility} with each
other, but a reduction internal to \cref{def:entropy-admissibility}
alone.

\begin{theorem}[Trajectory-admissibility reduces to static admissibility, checked pointwise]
\label{thm:pointwise-reduction}
For a trajectory $x(t)$ and a coordination geometry $\mathcal{G}(t)$ with
CLIO projection $\pi_\mathcal{G}$, ``$x(t)$ is continuation-admissible''
is equivalent to $S_{\mathrm{RSVP}}(\pi_\mathcal{G}(x(t))) \leq 0$ holding
at every $t$: the same static condition of \cref{def:entropy-admissibility},
evaluated at each moment along the path, introduces no further
structure. When this condition is violated, the coordination geometry
reorganizes --- geometric repair --- until it is restored.
\end{theorem}

\Cref{thm:pointwise-reduction} is proved by direct construction, not
asserted by resemblance. Its consequence for continuation is specific.
A theory of continuation does not need to introduce a second notion of
admissibility, or a new axiom governing how admissibility changes over
time; it needs only to define trajectories properly, use
\cref{thm:pointwise-reduction}, and examine what its violation condition
--- geometric repair --- has in common with operator ecologies (Chapter
5), since both describe a system reorganizing itself to restore a
condition that has been violated. That comparison, not new foundational
machinery, is what a theory of continuation is actually for.

\begin{remark}[This also explains the admissible future region]
\label{rem:gamma-explained}
The admissible future region $\Gamma_t(x) = \overline{\mathcal{R}}_t(x)$
of Chapter 5 is not a third notion of admissibility. It is
\cref{def:entropy-admissibility}'s reachability relation, restricted to
the effective (quotiented) state space rather than the full state space,
for the specific case where entropy production is realized as
distinction coarsening. This confirms that the inventory of
admissibility notions in this volume --- deficit-based, entropy-based,
and the architecture gate --- is complete, with $\Gamma_t(x)$ falling
under the second rather than constituting a fourth.
\end{remark}

\section*{Summary}

Two independent constructions of admissibility --- a deficit-based
relation and an entropy-monotonicity condition --- pass a type test but
not a substitution test; whether they induce the same continuation
family is the volume's central open problem, sharpened here from a bare
quantity comparison into a question about set relations between two
well-typed objects, and left open. The architecture gate is shown to be
a coarsening of a deficit-style construction without being identified
with the specific deficit of Chapter 1. Continuation requires no
independent foundational treatment beyond the material developed here,
applied pointwise.

\chapter{Constraint Populations}
% ------------------------------------------------------------------

\section{From a Single Region to a Population}

The admissibility relations developed earlier in this book fix a single
region --- a viability manifold $A \subseteq X$ --- and ask whether a
trajectory remains inside it. This chapter considers instead a
\emph{population} of admissibility constraints $\chi_1, \ldots, \chi_n$,
each $\chi_i$ specifying its own admissible region $A_i \subseteq X$,
with the population's joint admissible region given by
\[
A \;=\; \bigcap_{i=1}^n A_i.
\]
Nothing about this construction is new on its own; intersecting
constraints is the ordinary way multiple requirements combine, and each
$\chi_i$ individually is admissibility-typed in exactly the sense already
developed --- it fails the type test to call it a distinct primitive.
What is new is treating the $\chi_i$ themselves as objects with their own
histories and dynamics rather than as fixed, given sets: a building code
that is amended, a type system that is extended, a norm that shifts. Each
$\chi_i$ is a compressed operator $\chi_i = C(H_i)$ arising from some
history of prior cases, and each is subject to its own recursive
modification, $\chi_i(t+1) = G_i(\chi_i(t), \ldots)$. The question this
chapter asks is what happens once that update rule is allowed to depend
on the state of the \emph{other} constraints in the population, not only
on $\chi_i$'s own trajectory.

\section{Constraint Populations Are Sympoietic Systems}

\begin{definition}[Coupled recursive update]
\label{def:coupled-recursive}
A collection of objects $o_1, \ldots, o_n$ exhibits coupled recursive
update when each $o_i$'s own update rule at time $t+1$ depends not only
on $o_i$'s own history but on the joint state of the whole collection:
$o_i(t+1) = J_i(o_i(t), E_{t+1})$, where $E_{t+1} = H(E_t, o_1(t),
\ldots, o_n(t))$ aggregates the current state of every object in the
collection. This is a Level-3 phenomenon in the sense that the rule
governing $o_i$'s future behavior, not merely $o_i$'s state, is the
object being modified --- and it is a genuinely collective one: the
population as a whole can exhibit coupled recursive update even where no
individual $o_i$, considered in isolation from the others, does.
\end{definition}

\begin{proposition}[Constraint populations are sympoietic]
\label{prop:sympoietic}
Let $\chi_1, \ldots, \chi_n$ be a constraint population sharing a common
state space $X$, with each $\chi_i$'s update depending on the joint
state of the population in the sense of \cref{def:coupled-recursive}:
$\chi_i(t+1) = J_i(\chi_i(t), E_{t+1})$, $E_{t+1} = H(E_t, \chi_1(t),
\ldots, \chi_n(t))$. Then the population exhibits coupled recursive
update at the collective level directly, by \cref{def:coupled-recursive}
applied with $o_i = \chi_i$.
\end{proposition}

\begin{proof}
Immediate from \cref{def:coupled-recursive}: the population's update
structure is stated in exactly that definition's coupled form, with
$\chi_i$ substituted for the general object $o_i$.
\end{proof}

\begin{remark}[The same pattern recurs for operator ecologies]
\label{rem:same-pattern-ch6}
Chapter 6 develops this same coupling structure again, independently,
for a different population: the operators of an operator ecology, each
compressed from its own history and each capable of having its update
depend on the state of the ecology around it. Nothing in that later
development is assumed here, and nothing here is assumed there; the two
are separate instances of \cref{def:coupled-recursive}; a reader who
meets Chapter 6's version after this one will recognize the same shape
appearing at a different scale, with operators playing the role
constraints play here.
\end{remark}

\begin{remark}[Why this costs nothing new to state]
A body of statutory law that shifts only through the accumulated,
individually modest effect of many separate judicial rulings, no one of
which represents or intends to modify the law's overall shape, is a
direct instance: each ruling is a local, low-level update, and the shape
of the law that emerges is a collective Level-3 phenomenon, with case law
standing in for the aggregate of individually unremarkable local
changes. What was informally described as one constraint's violation
being ``absorbed'' by neighboring constraints tightening or loosening
elsewhere is simply \cref{prop:sympoietic} in operation: $\chi_i$'s
update rule $J_i$ is permitted to respond to the joint state $E_{t+1}$,
which encodes the current state of every other constraint in the
population, not merely $\chi_i$'s own history. No new mechanism is
required beyond the coupling already built into
\cref{def:coupled-recursive}.
\end{remark}

\section{Two Independent Pathologies}

\begin{definition}[Feasibility]
A constraint population is \emph{feasible} if $A = \bigcap_i A_i \neq
\varnothing$.
\end{definition}

\begin{definition}[Constraint diversity]
Let $\bar\rho \in [0,1]$ denote the average pairwise correlation among
the $\chi_i$'s failure events --- the frequency with which they are
jointly violated by the same states --- and define
\[
D_\chi \;=\; \frac{n}{1 + (n-1)\bar\rho}.
\]
\end{definition}

\begin{corollary}[Constraint monoculture]
As $\bar\rho \to 1$, $D_\chi \to 1$ regardless of $n$: a single failure
mode that defeats one $\chi_i$ defeats the population as a whole, and
many nominally distinct constraints provide no more protection against
inadmissible drift than a single constraint would.
\end{corollary}

\begin{remark}[Against a folk intuition]
Ordinary usage suggests that constraints cooperating --- pulling in the
same direction, reinforcing one another --- is the healthy condition,
and constraints competing is the pathological one. This is not quite
right, and in one respect is backwards. Constraints always violated by
the same states and always satisfied by the same states are exactly the
$\bar\rho \to 1$ case: however many there are, they function as one
constraint wearing many names, and the population inherits the fragility
of a single point of failure. What ordinary usage calls competition ---
constraints binding on different states, ruling out different failure
modes --- is closer to the $\bar\rho \to 0$ case identified as healthy
diversity. Genuine pathology is not disagreement between constraints; it
is either infeasibility or redundancy dressed up as multiplicity.
\end{remark}

\section{Feasibility Constrains Measured Diversity --- Sometimes Perversely}

\begin{lemma}[Two-constraint infeasibility bound]
\label{lem:twoconstraint}
Let $\chi_1, \chi_2$ be a population of two constraints on a finite
state space, with $\bar\rho$ their pairwise correlation computed as the
standard $\phi$ coefficient. If the population is infeasible, $\bar\rho
\leq 0$.
\end{lemma}

\begin{proof}
Infeasibility means the contingency-table cell $n_{00} = 0$. The $\phi$
coefficient is $\phi = (n_{11}n_{00} - n_{10}n_{01}) /
\sqrt{n_{1\cdot}n_{0\cdot}n_{\cdot1}n_{\cdot0}}$; with $n_{00} = 0$ the
numerator reduces to $-n_{10}n_{01} \leq 0$, and the denominator is
non-negative.
\end{proof}

Two genuinely incompatible constraints can never register as redundant
under this measure: infeasibility, at $n = 2$, structurally pushes
measured diversity up, not down. Whatever pathology infeasibility
represents, it is not the pathology of redundancy, and the two cannot be
conflated even in principle at this population size. The next result
shows this protection does not survive a third constraint.

\begin{proposition}[Diversity can mask redundancy in the presence of conflict]
\label{prop:masking}
There exists a population of four constraints, jointly infeasible, in
which three of the four are pairwise identical, and yet the population's
average pairwise correlation is exactly $\bar\rho = 0$, giving $D_\chi =
n = 4$: the score a population of four genuinely independent constraints
would report.
\end{proposition}

\begin{proof}
Let $X = \{1,\ldots,8\}$. Let $\chi_1 = \chi_2 = \chi_3$ each admit
$\{1,2,3,4\}$ (violated exactly on $\{5,6,7,8\}$), and let $\chi_4$
admit $\{5,6,7,8\}$ (violated exactly on $\{1,2,3,4\}$). The joint
admissible set is $\{1,2,3,4\} \cap \{5,6,7,8\} = \varnothing$: the
population is infeasible. Among the $\binom{4}{2} = 6$ pairs, the three
pairs within $\{\chi_1,\chi_2,\chi_3\}$ have violation indicators
agreeing everywhere, giving $\phi = +1$ each. The three pairs crossing
to $\chi_4$ have violation indicators that are exact complements,
giving $\phi = -1$ each. The average over all six pairs is $\bar\rho =
(3\cdot(+1) + 3\cdot(-1))/6 = 0$ exactly, giving $D_\chi = 4/(1 + 3
\cdot 0) = 4$.
\end{proof}

This is the chapter's central finding. The positive correlation
contributed by real redundancy and the negative correlation contributed
by genuine conflict cancel in a simple average, and the resulting scalar
$D_\chi$ reports the population as maximally diverse despite
three-quarters of it being a single constraint wearing three names.
Population health cannot be read off $D_\chi$ alone whenever redundancy
and conflict are simultaneously present: \cref{lem:twoconstraint,prop:masking}
together show the measure requires a feasibility check alongside it, not
merely as an independent second axis, but because feasibility failures
can actively corrupt what the diversity score appears to say once three
or more constraints are in play.

\begin{proposition}[Feasible populations span both high and low diversity]
\label{prop:feasiblespan}
Feasibility does not determine $D_\chi$: there exist feasible
populations with $D_\chi$ well above $n$ and feasible populations with
$D_\chi$ at its minimum of $1$.
\end{proposition}

\begin{proof}
On $X = \{1,2,3,4\}$: let $\chi_1$ admit $\{1,2,3\}$ and $\chi_2$ admit
$\{2,3,4\}$. Then $A = \{2,3\} \neq \varnothing$, and the violation
indicators give $\phi = -1/3$, so $D_\chi = 2/(1 - 1/3) = 3$, above $n =
2$. By contrast, let $\chi_1 = \chi_2$ both admit $\{1,2,3\}$: $A =
\{1,2,3\} \neq \varnothing$, but the two constraints are violated by
exactly the same state in every case, giving $\phi = 1$ and $D_\chi = 1$
despite $n = 2$.
\end{proof}

\begin{remark}[$D_\chi$ has no finite upper bound]
\label{rem:no-upper-bound}
Unlike the impression the equal-and-opposite cancellation of
\cref{prop:masking} might give, $D_\chi = n/(1 + (n-1)\bar\rho)$ is not
bounded above by $n$, or by any finite value, once $\bar\rho$ is allowed
to go negative: as $\bar\rho \to -1/(n-1)$ from above, the denominator
approaches $0^+$ and $D_\chi \to \infty$. \cref{prop:feasiblespan}'s own
example already exceeds $n$. The value $D_\chi = n$ obtained in
\cref{prop:masking} is not a ceiling being reached; it is the specific
value $\bar\rho = 0$ happens to produce, arising there from an exact
cancellation between the redundant trio's positive correlation and the
conflicting pair's negative correlation. Nothing about the measure
forbids scores above $n$, and a reader computing examples independently
should expect to encounter them.
\end{remark}

Feasibility must be checked directly, $A \neq \varnothing$, and
diversity must be checked directly, via $\bar\rho$ computed with
attention to whether conflicting and redundant subgroups might be
present simultaneously. Neither can be inferred from the other, and in
populations of three or more constraints, a favorable diversity score
cannot be trusted at face value without also checking whether it is the
product of exactly this kind of cancellation.

\section{Marginal Diversity Contribution}
\label{sec:marginal}

\cref{prop:masking} shows a population-level score can be misleading.
It does not yet say which member of the population is responsible. This
section supplies that.

\begin{definition}[Marginal diversity contribution]
\label{def:marginal-diversity}
For a constraint population $C = \{\chi_1,\ldots,\chi_n\}$ and $c \in
C$, the marginal diversity contribution of $c$ is
\[
\Delta_D(c) \;=\; D_\chi(C) \;-\; D_\chi(C \setminus \{c\}),
\]
where $D_\chi(C \setminus \{c\})$ is computed over the remaining $n-1$
constraints, with $\bar\rho(C \setminus \{c\})$ recomputed as the
average pairwise correlation among only those $n-1$ constraints ---
excluding every pair that involved $c$, not merely subtracting $c$'s own
contribution from the original average.
\end{definition}

\begin{definition}[Classification]
\label{def:diversity-classification}
$c$ is \emph{diversity-positive} if $\Delta_D(c) > 0$ (removing $c$
would decrease measured diversity: $c$ is contributing genuine
coverage), \emph{diversity-neutral} if $\Delta_D(c) = 0$, and
\emph{diversity-negative}, or \emph{invasive}, if $\Delta_D(c) < 0$
(removing $c$ would increase measured diversity: $c$'s presence is
actively suppressing the population's diversity score, typically
because it duplicates coverage another constraint already provides).
\end{definition}

\begin{corollary}[The masking population, decomposed]
\label{cor:masking-decomposed}
For the population of \cref{prop:masking}, $\Delta_D(\chi_1) = -5$ and
$\Delta_D(\chi_4) = +3$.
\end{corollary}

\begin{proof}
$D_\chi(C) = 4$ by \cref{prop:masking}. Removing $\chi_1$ leaves
$\{\chi_2,\chi_3,\chi_4\}$: the surviving pair $(\chi_2,\chi_3)$ has
$\phi = +1$, and the two pairs crossing to $\chi_4$ each have $\phi =
-1$, giving $\bar\rho = (1 - 1 - 1)/3 = -1/3$ and $D_\chi(C \setminus
\{\chi_1\}) = 3/(1 - 2/3) = 9$, so $\Delta_D(\chi_1) = 4 - 9 = -5$.
Removing $\chi_4$ leaves $\{\chi_1,\chi_2,\chi_3\}$, all three pairs at
$\phi = +1$, giving $\bar\rho = 1$ and $D_\chi(C \setminus \{\chi_4\}) =
3/(1+2) = 1$, so $\Delta_D(\chi_4) = 4 - 1 = 3$.
\end{proof}

The population-level score $D_\chi(C) = 4$ gives no reason to suspect
anything is wrong. The marginal decomposition immediately does: three
of the four constraints are invasive by \cref{def:diversity-classification},
each with $\Delta_D = -5$ by the symmetry of the construction, and the
one constraint actually responsible for the population's apparent
diversity, $\chi_4$, is identified directly, with a marginal
contribution larger than its share of the headcount would suggest. This
is the masking pathology, resolved rather than merely diagnosed: reading
$D_\chi(C)$ alone cannot distinguish this population from four genuinely
independent constraints, but reading every $\Delta_D(c)$ can.

\begin{remark}[The same construction predicts, rather than only diagnoses]
\label{rem:invasive-prediction}
\cref{def:marginal-diversity} is stated for $c$ already a member of
$C$, answering which existing members are worth keeping. The same
computation, applied to a candidate $c$ not yet in $C$ by evaluating
$\Delta_D(c)$ within $C' = C \cup \{c\}$, answers a different question
with the same arithmetic: whether introducing $c$ would help, do
nothing, or hurt. $\Delta_D(c)$ computed in $C \cup \{c\}$ is exactly
$D_\chi(C \cup \{c\}) - D_\chi(C)$, since removing $c$ from $C \cup
\{c\}$ recovers $C$ --- the pruning question and the introduction
question are the same construction viewed from opposite directions,
not two separate tools.
\end{remark}

\section{Worked Examples}

\textbf{Constitutional and statutory law.} A constitution changes
slowly and sets the outer admissible region within which statutes,
which change quickly, further narrow the admissible space. This is a
two-timescale sympoietic pair in the sense of \cref{prop:sympoietic}:
statutory revision is the fast, individually modest update; constitutional
interpretation shifts more slowly and in response to the accumulated
pattern of statutory and judicial activity.

\textbf{Building codes and inspection regimes.} A code is a constraint;
an inspection regime is a verification channel. Multiple inspectors
trained identically, using the same checklist and assumptions, add
little effective diversity however many of them there are, while
inspectors drawing on independent training or independent professional
traditions raise $D_\chi$ meaningfully even at the same headcount.

\textbf{Type systems and runtime checks.} A static type system and a
runtime assertion checking the same property are a low-$D_\chi$ pair:
both fail together on the same category of error. A static type system
paired with a runtime check for a genuinely different class of failure
is a high-$D_\chi$ pair, catching different failure modes at different
times.

\textbf{Gene regulatory networks.} Transcription factor networks
routinely exhibit precisely this structure: feedback loops in which one
regulatory constraint's activation state governs another's threshold, a
pattern recurring across organisms because sympoietic constraint
coupling of this kind is a repeatedly re-discovered solution rather than
a special case.

\section{Connections to the Rest of This Volume}

This chapter sits at the intersection of three threads, two of which lie
ahead rather than behind it. Admissibility, developed in the preceding
chapter, supplies the single-constraint apparatus --- viability
manifolds and the boundary within which structure becomes possible ---
that this chapter extends from one constraint to a population of them.
The coupled-recursive-update structure of \cref{def:coupled-recursive}
anticipates Chapter 6's own treatment of operator ecologies, where the
same pattern recurs with operators as the coupled objects rather than
constraints, and where the compression apparatus underlying $\chi_i =
C(H_i)$ here reappears as Chapter 6's treatment of operators as
compressed histories. Chapter 8's coordination geometry, concerned with
how a single geometry reorganizes to restore a violated admissibility
condition, will turn out to be a special case of what this chapter
develops: a coordination geometry is a constraint population of size
one, and geometric repair is what coupled recursive update looks like
when $n = 1$. None of these connections is claimed as a theorem
belonging to the later chapters specifically; they are recorded here as
the shape a fuller unification would need to take, to be confirmed or
revised once those chapters exist.

\section*{Open Problems}

\begin{openproblem}[Ecological stability and extinction]
Under what conditions does a constraint in a population become
effectively vacuous --- never binding, because other constraints
already exclude every state it would have ruled out --- and should a
vacuous constraint be understood as having gone extinct in the sense
Chapter 6's treatment of operator ecologies will later give that word?
\end{openproblem}

\begin{remark}[Invasive constraints, resolved]
\label{rem:invasive-resolved}
The criterion this population lacked when first raised is
\cref{def:marginal-diversity} applied via \cref{rem:invasive-prediction}:
whether a candidate $\chi_{n+1}$ will be absorbed, be diversity-neutral,
or destabilize the population's diversity score is $\Delta_D(\chi_{n+1})$
computed within $C \cup \{\chi_{n+1}\}$, positive, zero, or negative
respectively --- computable before introducing the constraint from its
pairwise correlations with each existing member alone. What remains
genuinely open is narrower: rejection (infeasibility) is a separate
failure mode from destabilization (a negative but still feasible
$\Delta_D$), and nothing here characterizes which candidates produce
which of the two, only that diversity impact and feasibility impact must
both be checked, neither implying the other, in keeping with
\cref{prop:feasiblespan}.
\end{remark}

\begin{openproblem}[Dynamics of $J_i$]
\cref{prop:sympoietic} assumes each constraint's update rule $J_i$
exists without specifying its form. A general theory of how
aggressively a constraint should tighten or loosen in response to
violations elsewhere is left to future work.
\end{openproblem}

\section*{Summary}

A population of admissibility regions is not adequately described by
its size, nor by a single scalar diversity score taken at face value.
Feasibility and diversity are independent axes that must each be
checked directly; for populations of three or more, they can actively
corrupt each other's apparent reading, with genuine redundancy and
genuine conflict cancelling to report a maximally diverse population
that is mostly one constraint under several names. This is not a
restatement of single-region admissibility; it is what single-region
admissibility looks like once more than one region is held
simultaneously, coupled through exactly the sympoietic machinery
Chapter 6 already built for a different purpose.

\chapter{Entropy}
% ------------------------------------------------------------------

\section{Introduction: Five Candidates, Not One}
\label{sec:candidates}

Chapter 2 developed two quantities under headings that both, in ordinary
usage, fall under ``entropy'': the Shannon entropy of a partition and the
representational entropy of a projection fiber. Whether either coincides
with the entropy used to define admissibility is not assumed here.
Widening the search turns up at least three further candidates, for five
in total.

\begin{itemize}
\item[A.] \textbf{Shannon/distinction entropy}, $H(\mathcal{P}) =
      -\sum_i p_i \log p_i$, developed in \cref{sec:shannon-recall}.
\item[B.] \textbf{Representational entropy}, $S_\pi(m) =
      \log|\pi^{-1}(m)|$, developed in \cref{sec:rep-recall}.
\item[C.] \textbf{Historical entropy}, the multiplicity of histories
      collapsing to a single state under a history-to-state projection,
      examined in \cref{sec:historical}.
\item[D.] \textbf{Repair entropy}, $\SR(F \mid T)$, the diagnostic
      uncertainty over fault locations given a retained trace, already
      used without full justification in Chapter 6, examined in
      \cref{sec:repair-entropy}.
\item[E.] \textbf{The RSVP entropy field}, $S(x,t)$, entering the
      admissibility relation of the RSVP formalization via the
      monotonicity constraint $\dot{S} \geq 0$, examined in
      \cref{sec:rsvp-entropy}.
\end{itemize}

The goal of this chapter is to determine, for each pair, whether the
relationship is identity, derivation, structural analogy, or genuine
independence --- and to say plainly which of these has actually been
checked against a proof and which remains a resemblance.

\section{Entropy A: Shannon/Distinction Entropy}
\label{sec:shannon-recall}

For a partition $\mathcal{P} = \{P_1,\ldots,P_n\}$
with cell probabilities $p_i$, $H(\mathcal{P}) = -\sum_i p_i \log p_i$ is
the expected surprisal of an observation identifying a cell, interpreted
distinction-theoretically as the expected refinement of the distinction
structure that an observation delivers.

\section{Entropy B: Representational Entropy}
\label{sec:rep-recall}

For a representation $\pi : X \to M$, the
representational entropy at $m$ is $S_\pi(m) = \log|\pi^{-1}(m)|$, the
log-size of the fiber collapsed beneath $m$, already shown in Chapter 2
to bound Chapter 1's coding-deficit increase from below under projection.

\section{The Distinction--Entropy Duality: A and B Are Not Independent}
\label{sec:duality}

\begin{claim}[Distinction--Entropy Duality]
\label{law:duality}
Every distinction simultaneously produces information and entropy.
Information measures separation between partition cells; entropy
measures multiplicity within partition cells. For a partition
$\mathcal{P}$ with $N_D = |\mathcal{P}|$ cells, $H = \log N_D$; for the
multiplicity $\Omega(P_i) = |P_i|$ within a cell, $S(P_i) = k \log
\Omega(P_i)$ for a constant $k$. Both arise from the same partition,
viewed from opposite directions: refinement increases $N_D$ while
reducing average $\Omega(P_i)$, and coarsening does the reverse.
\end{claim}

\begin{remark}[This is a reduction, not a new object]
\label{rem:duality-reduction}
$S(P_i) = k \log \Omega(P_i)$ in \cref{law:duality} has exactly the form
of \cref{sec:rep-recall}'s representational entropy: $\Omega(P_i) =
|P_i|$ is the cardinality of a partition cell, and a partition cell is a
fiber of the quotient map $q : X \to X/{\sim}$ in the sense of Chapter 1.
Up to the constant $k$, $S(P_i)$ \emph{is} $S_q(P_i)$. The
Distinction--Entropy Duality is therefore not introducing a third
independent quantity alongside $A$ and $B$; it is the statement that $A$
and $B$ are dual descriptions of one partition structure, counting across
cells in one case and within a cell in the other. This chapter treats $A$
and $B$ as related rather than as two of the five genuinely separate
candidates from here on.
\end{remark}

\section{Entropy C: Historical Entropy Reduces to B}
\label{sec:historical}

\begin{claim}[Irreversibility Theorem]
\label{thm:irrev}
Irreversibility emerges whenever a history projection $\pi_H : \mathcal{H}
\to X$ is many-to-one: for $h_1 \neq h_2$ with $\pi_H(h_1) = \pi_H(h_2) =
x$, knowledge of the present state $x$ does not determine which history
occurred, since $\pi_H^{-1}$ is not unique.
\end{claim}

\begin{claim}[Historical entropy is representational entropy]
\label{claim:historical-reduces}
Historical entropy, informally the multiplicity of histories collapsing
to a state, is $\log|\pi_H^{-1}(x)|$ for the history projection
$\pi_H$ of \cref{thm:irrev}. This is exactly \cref{sec:rep-recall}'s
representational entropy $S_{\pi_H}(x)$, applied to the specific case
where the representation in question is a history-to-state projection.
Historical entropy is not a fourth candidate; it is Entropy B evaluated
at a particular kind of $\pi$.
\end{claim}

This reduction is a genuine derivation rather than an analogy: both
quantities are defined as the log-cardinality of a fiber under a
projection, and $\pi_H$ is a projection in exactly Chapter 1's sense.
Nothing about the reduction required inventing new machinery.

\section{Entropy Production, the Second Law, and Reachability}
\label{sec:production}

Three further results build on \cref{law:duality} without introducing
new entropy candidates.

\begin{claim}[Entropy Production Theorem]
\label{thm:production}
Coarsening of distinction structure --- merging cells $P_i$ and $P_j$
into $P_i \cup P_j$ --- increases $\log\Omega$ in the merged cell relative
to either constituent, since $|P_i| + |P_j| > \max(|P_i|,|P_j|)$. Entropy
increases wherever distinction structure coarsens.
\end{claim}

\begin{claim}[Second Law, reinterpreted]
\label{law:second}
In the absence of repair or distinction-generating dynamics, recoverable
distinction capacity $D(t)$ tends to decrease: $\dot{D} \leq 0$. Entropy
increase is the shadow cast by distinction loss.
\end{claim}

The relation between entropy and reachability deserves more care than
stating the two are monotonically linked. Coarsening observational
distinctions does not, by itself, shrink the set of states an underlying
system can \emph{physically} reach; it shrinks the set of reachable
states an observer can \emph{tell apart}. The quotient connecting the two
must be built explicitly rather than asserted.

\begin{definition}[Effective reachable set]
\label{def:effective-reach}
Let $\mathcal{R}(x)$ be the physically reachable set from $x$, and let
$q_t : X \to X/{\sim_t}$ be the quotient map induced by the
distinguishability relation at time $t$. The \emph{effective reachable
set} is $\overline{\mathcal{R}}_t(x) = q_t(\mathcal{R}(x))$, with
effective reachability volume $V_{\mathrm{eff}}(x,t) =
\bar\mu_t(\overline{\mathcal{R}}_t(x))$ for a measure $\bar\mu_t$ on the
quotient (cardinality, in the finite case).
\end{definition}

\begin{claim}[Effective Reachability Contraction]
\label{thm:effective-reach-contraction}
Suppose the distinction relation coarsens from $t_1$ to $t_2$:
$x \sim_{t_1} y \Rightarrow x \sim_{t_2} y$. Then there is a canonical
surjection $c_{21} : X/{\sim_{t_1}} \to X/{\sim_{t_2}}$ with $q_{t_2} =
c_{21} \circ q_{t_1}$, so $\overline{\mathcal{R}}_{t_2}(x) =
c_{21}(\overline{\mathcal{R}}_{t_1}(x))$, and consequently
$|\overline{\mathcal{R}}_{t_2}(x)| \leq |\overline{\mathcal{R}}_{t_1}(x)|$
in the finite case.
\end{claim}

\begin{proof}
Since $\sim_{t_2}$ is coarser, each $\sim_{t_1}$-class lies within a
unique $\sim_{t_2}$-class; sending the former to the latter defines
$c_{21}$. Then $\overline{\mathcal{R}}_{t_2}(x) = q_{t_2}(\mathcal{R}(x))
= c_{21}(q_{t_1}(\mathcal{R}(x))) = c_{21}(\overline{\mathcal{R}}_{t_1}(x))$,
and a surjective image cannot contain more distinct elements than its
domain.
\end{proof}

\Cref{thm:effective-reach-contraction} proves distinction loss implies
effective future contraction. It does not prove distinction loss implies
\emph{physical} future contraction --- that stronger claim requires
dynamics coupling representational loss back into executable action, and
holds only conditionally.

\begin{definition}[Distinction-dependent controllability]
\label{def:action-dependent}
A controlled system is \emph{distinction-dependent} when the executable
action set at $x$ depends only on its currently distinguished class:
$\mathcal{A}_t(x) = \mathcal{A}_t([x]_{\sim_t})$. Coarsening is
\emph{action-monotone} when $\sim_{t_1} \subseteq \sim_{t_2} \Rightarrow
\mathcal{A}_{t_2}(x) \subseteq \mathcal{A}_{t_1}(x)$ for every $x$.
\end{definition}

\begin{claim}[Physical Reachability Contraction under Action Dependence]
\label{thm:physical-contraction}
If distinction coarsening is action-monotone, the physically reachable
set generated by executable actions satisfies $\mathcal{R}_{t_2}(x)
\subseteq \mathcal{R}_{t_1}(x)$.
\end{claim}

\begin{proof}
Every trajectory available at $t_2$ is composed from actions in
$\mathcal{A}_{t_2}$, each of which, by action-monotonicity, was already
available at $t_1$. Hence every $t_2$-trajectory is also a
$t_1$-trajectory.
\end{proof}

\Cref{thm:effective-reach-contraction,thm:physical-contraction} together
give this section's actual result: entropy production provably contracts
\emph{effective} reachability unconditionally, and contracts
\emph{physical} reachability under the further assumption that action
availability is itself distinction-dependent. The two should not be
conflated; the second is a substantive cybernetic hypothesis about how
representation couples to action, not a free consequence of the first.

\begin{definition}[Reachability admissibility]
\label{def:reach-admissibility}
A state $y$ is admissible from $x$ at time $t$ when its effective class
is reachable: $x \rightsquigarrow_t y \iff q_t(y) \in
\overline{\mathcal{R}}_t(x)$. The admissible future region is
$\Gamma_t(x) = \overline{\mathcal{R}}_t(x)$.
\end{definition}

\begin{claim}[Entropy--Admissibility Corollary]
\label{cor:entropy-admissibility}
If entropy production is realized as monotone coarsening of the
distinction relation, the effective admissible future region is
non-increasing: $\sim_{t_1} \subseteq \sim_{t_2} \Rightarrow \Gamma_{t_2}(x)
= c_{21}(\Gamma_{t_1}(x))$, and in the finite case $|\Gamma_{t_2}(x)| \leq
|\Gamma_{t_1}(x)|$.
\end{claim}

This is immediate from \cref{thm:effective-reach-contraction} under
\cref{def:reach-admissibility}'s identification of the admissible future
region with the effective reachable set. It gives the chapter sequence
this section was actually able to establish: Entropy $\to$ effective
reachability $\to$ admissibility, deliberately avoiding an unsupported
direct arrow from entropy to repair. What repair does with a contracted
admissible region is Chapter 6's business, not this chapter's.

\begin{remark}[Decorrelated dimensionality, not operator count]
\label{rem:posani-participation}
A useful comparison, and a useful caution, is provided by Posani et al.\
(2026), who study the dimensionality and separability of neural
population representations across the cortical hierarchy. Their
dimensionality measure is the participation ratio
\[
\operatorname{PR}(C) \;=\; \frac{\left(\sum_i \lambda_i\right)^2}{\sum_i \lambda_i^2},
\]
where the $\lambda_i$ are eigenvalues of a response-covariance matrix
$C$. The participation ratio is not identical to any quantity defined in
this book, and it is worth stating precisely why, rather than leaving
the resemblance to do unearned work. It is not representational entropy:
$\operatorname{PR}(C) \neq \log|\pi^{-1}(m)|$. It is not reachability
volume: $\operatorname{PR}(C) \neq \mu(\mathcal{R}_\mathcal{E}(x))$. No
definition-by-substitution identifies these constructions, and they take
different inputs --- a covariance operator over observed responses,
against a set of attainable future states.

The comparison is nevertheless structurally informative, because of what
the participation ratio discounts. If covariance is concentrated in a
small number of eigen-directions, many recorded units still yield a low
effective dimensionality; a high value requires variance distributed
across several independently contributing directions. The lesson for an
operator ecology is that repertoire cardinality alone need not measure
generative diversity: $|\Reach(\mathcal{E})|$ can be large while the
operators it contains are nearly redundant, so that the induced future
region $\mu(\mathcal{R}_\mathcal{E}(x))$ remains small. This is not a new
observation so much as a sharpening of one already present elsewhere in
this corpus: a resilience discount for correlated checks, on the
principle that a hundred people relying on one hidden assumption do not
constitute a hundred independent checks. The participation ratio and
that resilience discount are not the same quantity, but both implement
the same correction --- correlated contributions count for less than
independent ones --- from two unrelated domains arriving at it
separately. Chapter 6 develops a repertoire-level analogue of this
correction directly.
\end{remark}

\section{Entropy D: Repair Entropy}
\label{sec:repair-entropy}

Chapter 6 already used $\SR(F \mid T)$, the repair entropy of an operator
$F$ given a retained trace $T$, to state the Anti-Alignment Theorem,
without independently justifying why this quantity deserves the name
entropy rather than merely borrowing it by analogy. That justification
can now be given in full, upgrading the reduction from a structural
resemblance to an ordinary conditional-entropy construction.

\begin{definition}[Fault hypothesis space]
\label{def:fault-space}
For an operator $F$ and observed symptom $s$, let $\Omega_F(s)$ be the
finite set of fault hypotheses capable of producing $s$. A retained trace
$T$ determines an evidence map $e_T : \Omega_F(s) \to E_T$, where
$e_T(f)$ records the facts about hypothesis $f$ distinguishable using
$T$. This induces a partition $\mathcal{P}_T = \{e_T^{-1}(a) : a \in
E_T\}$ of the fault hypothesis space, and the live hypotheses after
observing the actual trace value $a_T$ are
$\mathcal{H}_{\mathrm{fault}}(F,T) = e_T^{-1}(a_T)$.
\end{definition}

\begin{definition}[Repair entropy]
\label{def:repair-entropy-shannon}
For $p(f \mid s,T)$ a probability distribution over
$\mathcal{H}_{\mathrm{fault}}(F,T)$,
\[
\SR(F \mid T) \;=\; -\sum_{f \in \mathcal{H}_{\mathrm{fault}}(F,T)} p(f\mid s,T) \log p(f \mid s,T),
\]
which under a uniform prior becomes $\SR(F\mid T) = \log|\mathcal{H}_{\mathrm{fault}}(F,T)|$.
\end{definition}

\begin{claim}[Trace Refinement Theorem]
\label{thm:trace-refinement}
Suppose $T_2$ retains every diagnostically relevant fact retained by
$T_1$, and possibly more. Then $\mathcal{P}_{T_2}$ refines $\mathcal{P}_{T_1}$
and $H(F_{\mathrm{fault}} \mid s,T_2) \leq H(F_{\mathrm{fault}} \mid s,T_1)$,
hence under the uniform-prior convention $\SR(F\mid T_2) \leq \SR(F\mid T_1)$,
strictly whenever the additional information separates two hypotheses
that remained equivalent under $T_1$ and both had positive probability.
\end{claim}

\begin{proof}
Since $T_2$ contains all diagnostically relevant information in $T_1$,
there is a forgetful map $q : E_{T_2} \to E_{T_1}$ with $e_{T_1} = q
\circ e_{T_2}$. If $e_{T_2}(f) = e_{T_2}(g)$ then $e_{T_1}(f) = e_{T_1}(g)$,
so every cell of $\mathcal{P}_{T_2}$ lies inside a cell of $\mathcal{P}_{T_1}$:
$\mathcal{P}_{T_2}$ refines $\mathcal{P}_{T_1}$. Equivalently, the
variables form a Markov chain $F_{\mathrm{fault}} \to T_2 \to T_1$, with
$T_1$ obtained from $T_2$ by discarding information. Conditional entropy
is monotone under additional conditioning, giving the stated inequality;
under a uniform distribution, conditional entropy is the log-cardinality
of the surviving hypotheses, giving the second inequality. Strictness
follows whenever refinement splits a positive-probability cell.
\end{proof}

\begin{claim}[Anti-Alignment, as a corollary]
\label{cor:anti-alignment}
Let $T_{\min}$ be a minimal trace sufficient to enact an operator $F$,
and $T \supseteq T_{\min}$ retain additional diagnostic information.
Then $\SR(F\mid T) \leq \SR(F\mid T_{\min})$, strictly whenever the
discarded information excludes at least one otherwise live hypothesis.
\end{claim}

\begin{proof}
Apply \cref{thm:trace-refinement} with $T_1 = T_{\min}$, $T_2 = T$.
\end{proof}

\Cref{cor:anti-alignment} makes Chapter 6's Anti-Alignment result a
direct consequence of the partition structure developed in this chapter
and Chapter 2, rather than an imported standalone claim resting on
analogy. Repair entropy is accordingly reclassified: it is not a fifth
candidate alongside A through C, but Entropy A itself, applied to the
specific partition a retained trace induces on a space of fault
hypotheses.

\begin{remark}[Anti-Alignment is a monotonicity lemma, not an optimization result]
\label{rem:anti-alignment-scope}
\cref{cor:anti-alignment} shows $\SR(F \mid T)$ is weakly decreasing as
$T$ grows, and reads naturally as license to retain as much trace as
possible. That reading silently assumes retention is free. Making this
assumption explicit requires a retention cost functional $K(T)$, with
the actual objective
\[
J(T) \;=\; \SR(F \mid T) \;+\; \lambda K(T)
\]
for some trade-off parameter $\lambda > 0$. Under this objective the
optimum is generally not $T = T_{\max}$ but a minimal
\emph{admissibility-sufficient} trace $T^*$: enough retained history to
support every fault-diagnosis judgment that will actually be needed,
and no more. \cref{cor:anti-alignment} remains correct as a
monotonicity statement about $\SR$ alone; it does not by itself
determine $T^*$, since it says nothing about $K$. $K(T)$ and $T^*$ are
developed in Chapter 6, where retention cost is native; Anti-Alignment
should be read as one term of a larger optimization, not as a
standalone argument for maximal retention.
\end{remark}

\section{Entropy E: The RSVP Field, Left Open}
\label{sec:rsvp-entropy}

The RSVP formalization defines an admissibility relation directly in
terms of a field triple $(\Phi, \mathbf{v}, S)$, where $S : M \to
\mathbb{R}_{\geq 0}$ is required to satisfy $\dot{S} \geq 0$ along
admissible trajectories --- entropy monotonicity is part of what
admissibility \emph{means} in that formalization, not a consequence
derived after the fact.

\begin{openproblem}
\label{prob:rsvp-entropy}
No result here derives $S(x,t)$ from the partition- and fiber-based
entropies $A$ through $D$ above. The RSVP fields are motivated by the
distinction-repair-reachability chain rather than derived from it, and
the specific functional form relating reachability volume to $S$ is a
heuristic argument, not a uniqueness theorem. There is a plausible
structural correspondence worth flagging rather than asserting: the
Second Law of \cref{law:second} gives $\dot{D} \leq 0$ for a
distinction-capacity quantity $D$, and $\dot{S} \geq 0$ in the RSVP
formalization is the same monotonic inequality under a sign flip,
suggestive of $S$ playing the role of consumed distinction capacity
rather than distinction capacity itself. This is a resemblance in the
direction and character of two monotonicity constraints, not a proof
that $S$ and $D$ are the same quantity, functionally related, or
defined over the same underlying space. Establishing or refuting this
correspondence is left open.
\end{openproblem}

\begin{remark}[Reconstructive and thermodynamic irreversibility are distinct]
\label{rem:hydra-confirms}
Reconstructive irreversibility --- the fiber non-injectivity of a
transition function, a purely combinatorial property --- is distinct
from thermodynamic irreversibility: the former concerns neither entropy,
phase-space volume, nor coarse-graining directly. This distinction, drawn
from fiber and distinguishability apparatus essentially identical to
entropies $A$ through $D$, is independent evidence for the caution urged
in \cref{prob:rsvp-entropy}. It does not resolve the open problem. It
strengthens the case that leaving it open is the right call.
\end{remark}

\section{Summary}
\label{sec:summary}

\begin{center}
\begin{tabular}{lll}
\hline
Candidate & Status & Relation \\
\hline
A: Shannon/distinction & Checked & Dual to B (\cref{law:duality}) \\
B: Representational & Checked & Dual to A; C is an instance of B \\
C: Historical & Checked & Reduces to B (\cref{claim:historical-reduces}) \\
D: Repair & Checked & Is an instance of A (\cref{cor:anti-alignment}) \\
E: RSVP field & Open & Motivated, not derived; see \cref{prob:rsvp-entropy} \\
\hline
\end{tabular}
\end{center}

Of five candidates, four are reduced above to a single underlying
partition/fiber structure with genuine derivations, and the fifth ---
the one actually load-bearing for admissibility --- remains explicitly
open. That asymmetry matters: whatever is eventually said about
admissibility inherits an unresolved dependency on Entropy E, however
cleanly A through D have been tied together.

\section*{Summary}
\label{sec:bridge-forward}

The distinction/information/entropy sequence developed across Chapters
1--2 and this chapter does not extend cleanly into a further stage
called ``History.'' History (Entropy C) is not a separate node; it is
Entropy B evaluated at a particular projection. What connects entropy to
the second half of this volume is the Entropy--Admissibility Corollary
of \cref{cor:entropy-admissibility}, not any direct entropy-to-repair
theorem. Reachability, not history or recovery as separate stages, is
the actual bridge --- but a bridge with two distinctly typed ends rather
than one shared object. The executable repertoire $\Reach(\mathcal{E})$
of Chapter 6 is a set of \emph{operators}; the admissible future region
$\Gamma_t(x)$ and reachability volume $V_{\mathrm{eff}}(x,t)$ of this
chapter are sets and measures of \emph{states}. They fail a basic type
test and cannot be identical. What connects them is a generator
relation: an executable repertoire supplies the operators from which
state-to-state trajectories are composed, and the reachable-state set is
a derived closure of that repertoire, not the repertoire itself. That
closure map, and its monotonicity, is constructed in Chapter 6.

\chapter{Repair}
% ------------------------------------------------------------------

\section{Introduction}

A structure that persists over time is easy to mistake for a static
fact, something that simply continues to be the case once it has been
established. This chapter argues that persistence is never a static
fact. A structure persists because something is continually acting to
keep it in existence against the ordinary pressure of decay, damage, and
drift, and the moment that action stops, the structure does not
immediately vanish, but it has already stopped being maintained, and what
follows is only a matter of time. This chapter develops the vocabulary
needed to say precisely what that ``something'' is, what it means for it
to persist in turn, and what quantities govern how much repair capacity a
system actually has available to it.

\section{Three Levels of Modification}
\label{sec:levels}

Before asking what maintains a structure, it is worth being precise about
what kind of change is even available to a system that modifies itself,
since discussions of adaptation and repair routinely conflate phenomena
that differ not in how much behavior changes but in what kind of object
is doing the changing.

\begin{definition}[Levels of adaptation]
\label{def:levels}
Let $X$ be a state space, $\Theta$ a parameter space, and $\mathcal{F}$ a
space of maps $X \to X$. A system exhibits \emph{Level 1} (state update)
behavior when $x_{t+1} = F(x_t)$ for fixed $F$: the system moves through
$X$ under an unchanging rule. It exhibits \emph{Level 2} (parameter
update) behavior when $x_{t+1} = F_{\theta_t}(x_t)$ for fixed
$F : \Theta \times X \to X$ with $\theta_t \in \Theta$ varying: behavior
changes, but only within the range expressible by the fixed parametric
family $F$. It exhibits \emph{Level 3} (rule update) behavior when
$F_{t+1} = G(F_t, x_t, E_t)$ for some environment $E_t$: the mapping
governing future behavior is itself the object being modified, not merely
a parameter feeding into a fixed mapping.
\end{definition}

Repair, in the sense this chapter develops, is a Level 3 phenomenon. A
system that merely moves through its state space (Level 1) or adjusts
parameters within a fixed family (Level 2) is not repairing anything; it
is only running. Repair requires that the rule governing future behavior
--- the operator itself --- be a target of modification, which is why the
operator ecologies developed in this chapter are ecologies of Level 3
objects, not of states or parameters.

\section{From Structures to Repairs to Ecologies}
\label{sec:reframe}

The argument of this section is developed at length, with a worked
example of two initially identical cities that diverge over a decade
purely in whether ordinary maintenance --- road resurfacing, apprentice
training, records-keeping, governance continuity --- continues. Its
conclusion is stated here rather than its argument reproduced: state
equivalence between two systems is insufficient to predict their
futures, because two systems can agree in every recorded particular
while already diverging in
what each is still capable of doing about that state. Structural
equivalence fails for the same reason, one level down. The natural first
correction --- structures persist, and repair is something that
occasionally happens to them --- is shown there to have the explanatory
priority backwards. The corrected view is that repair occurs, and
structures persist because of it; and one level further, that repair
occurs because an operator ecology persists, and structures are, in the
resulting phrase, frozen operator ecologies: visible residues of a
network of constraints that, for as long as it holds, keeps producing
them.

\section{A Hierarchy of Equivalence}
\label{sec:equivalence}

\emph{Operator Ecology} develops a hierarchy of increasingly demanding
notions of sameness between two systems under damage, each an improvement
on the last, each still shown insufficient in turn until the final one.

\begin{definition}[Equivalence hierarchy]
\label{def:equiv-hierarchy}
Two systems are \emph{state equivalent} at time $t$ if they occupy the
same point in state space at $t$. They are \emph{structurally equivalent}
if their underlying architectures agree, independent of current state.
They are \emph{behaviorally equivalent} if they are bisimilar in the
coalgebraic sense: indistinguishable by any observation of their
input-output behavior, including future behavior, from the current state
onward. They are \emph{regeneratively equivalent} if they possess
sufficiently similar operator ecologies to recreate the same future
capacities after damage, independent of whether their current state or
current behavior agree at all.
\end{definition}

State and structural equivalence fail for the reason \cref{sec:reframe}
already gave: two systems, or two moments of one system, can agree on
both while their capacity to persist has already diverged. Behavioral
equivalence is a genuine improvement, since it tracks not a static fact
but an ongoing capacity, and bisimulation gives that capacity a precise
mathematical name. But behavioral equivalence is defined over a system's
\emph{current} behavior, and has nothing left to say once that behavior
has already been damaged or lost: two systems can be behaviorally
equivalent right up until the moment of damage and then diverge sharply
in what they can do afterward, a divergence behavioral equivalence alone
cannot see coming because it was never asking about recovery. Regenerative
equivalence is the notion built to see it coming: it is not a claim about
what a system currently does, but about what it retains the capacity to
do again.

\section{Operator Ecologies as Constraint Networks}
\label{sec:constraints}

The natural first attempt to formalize an operator ecology is as an
inventory: the set of operators a system currently has available. This
undercounts what actually explains regenerative equivalence, since an
inventory says nothing about why any particular operator remains
executable, or what would happen to the rest of the inventory if one
entry were lost.

\begin{definition}[Operator ecology]
\label{def:operator-ecology}
Following Juarrero's account of causation as constraint rather than
efficient push, an \emph{operator ecology} $\mathcal{E}$ is not the set of
operators it contains but the network of enabling constraints among them:
the structure that determines which operators remain executable, and
which become executable or cease to be executable as other parts of the
network change.
\end{definition}

This definition is what allows the ecological vocabulary of the next
section --- keystone, extinction, health --- to be more than metaphor.
An inventory has no notion of one entry depending on another; a constraint
network does, and the dependency structure is precisely what the
following quantities measure.

\section{Keystone Operators, Extinction, and Health}
\label{sec:keystone}

\begin{definition}[Reachable repertoire and keystone impact]
\label{def:keystone}
For an operator ecology $\mathcal{E}$, let $\Reach(\mathcal{E})$ denote
the reachable executable repertoire: the operators actually executable
given $\mathcal{E}$'s current constraint graph. The \emph{impact} of an
operator $o$ is
\[
K(o) \;=\; \bigl|\Reach(\mathcal{E})\bigr| \;-\; \bigl|\Reach(\mathcal{E}\setminus\{o\})\bigr|,
\]
the size of the repertoire lost by removing $o$ and everything its
removal disconnects. A \emph{keystone operator} is one with unusually
large $K(o)$ relative to how often it is itself invoked.
\end{definition}

Apprenticeship training in the two-cities example of \cref{sec:reframe}
is a keystone operator in exactly this sense: it is invoked rarely, but
its removal disconnects every operator that depended on the next
generation of practitioners it would otherwise have supplied.

\Cref{def:keystone}'s impact measure $K(o)$ counts lost operators. It
does not, by itself, say anything about the states or futures those
operators would have made reachable, which is a different and better
question to ask of a keystone loss. Answering it requires building the
map from an executable repertoire to a reachable-state set explicitly,
since Chapter 5's admissible future region $\Gamma_t(x)$ is a set of
states, not of operators, and the two constructions fail a basic type
test: $\Reach(\mathcal{E})$ is not the same kind of object as
$\Gamma_t(x)$, and no theorem in this book identifies them.

\begin{definition}[Operator-induced reachable set]
\label{def:operator-induced-reach}
Let $X$ be a state space and $\mathcal{O}_\mathcal{E} = \Reach(\mathcal{E})$
the currently executable repertoire of an ecology $\mathcal{E}$. For
$x \in X$ and a budget $n \in \mathbb{N}$, define
\[
\mathcal{R}_\mathcal{E}^{(n)}(x) = \bigl\{ o_k \circ \cdots \circ o_1(x) :
0 \leq k \leq n,\ o_i \in \mathcal{O}_\mathcal{E},\ \text{every intermediate
composition executable} \bigr\},
\]
with unrestricted reachable set $\mathcal{R}_\mathcal{E}(x) =
\bigcup_{n \geq 0} \mathcal{R}_\mathcal{E}^{(n)}(x)$ and, where $X$ carries
a measure $\mu$, induced reachability volume $V_\mathcal{E}(x) =
\mu(\mathcal{R}_\mathcal{E}(x))$.
\end{definition}

This is the closure map Chapter 5 asked for: it takes a repertoire of
operators and returns the set of states composable from them, which is
the right kind of object to compare against $\Gamma_t(x)$. A basic
monotonicity result requires ruling out operators that disable one
another's compositions.

\begin{definition}[Extension-monotone ecology]
\label{def:extension-monotone}
An operator ecology is \emph{extension-monotone} when $\mathcal{O}_{\mathcal{E}_1}
\subseteq \mathcal{O}_{\mathcal{E}_2}$ implies every composition
executable in $\mathcal{E}_1$ remains executable in $\mathcal{E}_2$.
\end{definition}

\begin{claim}[Repertoire--Reachability Monotonicity]
\label{thm:repertoire-reachability}
For extension-monotone ecologies, $\Reach(\mathcal{E}_1) \subseteq
\Reach(\mathcal{E}_2) \Rightarrow \mathcal{R}_{\mathcal{E}_1}(x) \subseteq
\mathcal{R}_{\mathcal{E}_2}(x)$ for every $x$, and consequently
$V_{\mathcal{E}_1}(x) \leq V_{\mathcal{E}_2}(x)$ whenever $\mu$ is
monotone under inclusion.
\end{claim}

\begin{proof}
Let $y \in \mathcal{R}_{\mathcal{E}_1}(x)$, reached by an executable
composition $y = o_k \circ \cdots \circ o_1(x)$ with every $o_i \in
\Reach(\mathcal{E}_1)$. Repertoire inclusion places every $o_i$ in
$\Reach(\mathcal{E}_2)$, and extension-monotonicity ensures the same
composition remains executable there, so $y \in \mathcal{R}_{\mathcal{E}_2}(x)$.
The volume inequality follows from monotonicity of $\mu$.
\end{proof}

Without extension-monotonicity the inclusion can fail outright: an
ecology's operators can interact so that adding one disables a
composition that a smaller repertoire supported, and \cref{thm:repertoire-reachability}
is not claimed to hold in that regime.

Chapter 5 flags, via a comparison to Posani et al.\ (2026)'s work on
cortical population dimensionality, that $|\Reach(\mathcal{E})|$ is a
poor proxy for generative diversity whenever the repertoire's operators
are redundant with one another. That comparison is illustrative there,
not formal; a repertoire-native version of the same correction can be
built directly from this section's own constructions.

\begin{definition}[Effective repertoire dimension]
\label{def:effective-repertoire-dimension}
Suppose the local effect of each executable operator $o \in \Reach(\mathcal{E})$
at $x$ is represented by a vector $v_o(x) \in T_x X$ in a linearized
tangent space. For non-negative invocation weights $w_o$, define
\[
C_\mathcal{E}(x) \;=\; \sum_{o \in \Reach(\mathcal{E})} w_o\, v_o(x) v_o(x)^\top.
\]
If $\lambda_1,\ldots,\lambda_n$ are the eigenvalues of $C_\mathcal{E}(x)$,
the \emph{effective repertoire dimension} is
\[
D_{\mathrm{eff}}(\mathcal{E},x) \;=\; \frac{\left(\sum_i \lambda_i\right)^2}{\sum_i \lambda_i^2}.
\]
\end{definition}

\begin{claim}[Bounds on effective repertoire dimension]
\label{thm:effective-repertoire-bounds}
For any nonzero positive semidefinite $C_\mathcal{E}(x)$,
\[
1 \;\leq\; D_{\mathrm{eff}}(\mathcal{E},x) \;\leq\; \operatorname{rank} C_\mathcal{E}(x).
\]
Equality on the left holds when all operator effects lie in one
effective direction; equality on the right holds when the nonzero
eigenvalues are equal.
\end{claim}

\begin{proof}
Let $r = \operatorname{rank} C_\mathcal{E}(x)$ with nonzero eigenvalues
$\lambda_1,\ldots,\lambda_r > 0$. Since $\left(\sum_i \lambda_i\right)^2
\geq \sum_i \lambda_i^2$, the lower bound follows. By Cauchy--Schwarz,
$\left(\sum_{i=1}^r \lambda_i\right)^2 \leq r \sum_{i=1}^r \lambda_i^2$,
giving the upper bound. The equality conditions are the standard ones
for the two inequalities.
\end{proof}

One hundred collinear operator effects give $D_{\mathrm{eff}} = 1$
regardless of $|\Reach(\mathcal{E})|$; $k$ orthogonal, equally weighted
effects give $D_{\mathrm{eff}} = k$. This gives the correlated-operators
concern a precise, checkable form rather than leaving it at the level of
intuition.

\begin{remark}[A local proxy, not a replacement]
\label{rem:effective-repertoire-caveat}
$D_{\mathrm{eff}}(\mathcal{E},x)$ measures local effective action
dimension at a single state $x$, under a linearization. It is not
$\mu(\mathcal{R}_\mathcal{E}(x))$, and \cref{thm:repertoire-reachability}'s
monotonicity is not claimed to transfer to it automatically.
Noncommutativity among operators, state-dependent executability, and the
compounding effects of long compositions can all make ecologies with
identical $D_{\mathrm{eff}}(\mathcal{E},x)$ generate different global
reachable regions. This quantity is offered as a cheaper, locally
computable predictor of generative diversity, not as a substitute for
\cref{def:operator-induced-reach}'s actual reachable set when the two
disagree.
\end{remark}

\Cref{def:operator-induced-reach} also gives keystone impact a second,
more consequential measure than $K(o)$.

\begin{definition}[Generative keystone impact]
\label{def:generative-keystone}
The generative impact of $o$ at $x$ is $K_V(o;x) = \mu(\mathcal{R}_\mathcal{E}(x))
- \mu(\mathcal{R}_{\mathcal{E}\setminus\{o\}}(x))$.
\end{definition}

\begin{remark}[Inventory impact and generative impact can diverge]
\label{rem:keystone-divergence}
There is no reason for $K(o)$ and $K_V(o;x)$ to rank operators
identically. Repertoire impact counts lost executable operators;
generative impact measures the future region their loss disconnects. A
single operator may account for only one lost repertoire element under
$K$ while controlling access to a vast region of state space under
$K_V$, or the reverse. The former is architectural; the latter is
dynamical. This is why \cref{def:health}'s ecological health is defined
below in terms of recoverable futures rather than repertoire size: health
was already, implicitly, tracking something closer to $K_V$ than to $K$,
without the intermediate construction that makes the relationship
precise.
\end{remark}

\begin{definition}[Extinction]
\label{def:extinction}
An operator is \emph{extinct} within an ecology $\mathcal{E}$ when no
executable realization of it remains, regardless of whether a
description of it survives.
\end{definition}

Extinction in this sense is deliberately independent of documentation.
City B's apprenticeship program, in the example of \cref{sec:reframe},
can be extinct while a charter document describing it still sits,
accurately and uselessly, in an archive: the description survives, the
executable operator does not. This is the same distinction Chapter 7
makes between representation and executability in its architecture gate,
and the same distinction that recurs, independently, in the
enactment/representation asymmetry of \cref{sec:history} below.

\begin{definition}[Ecological health]
\label{def:health}
Writing $\Gamma(\mathcal{E})$ for the set of futures recoverable from
$\mathcal{E}$ under plausible further damage, the \emph{health} of an
operator ecology is $\Health(\mathcal{E}) = \mu(\Gamma(\mathcal{E}))$ for
some measure $\mu$. Extinction is what decreases $\Health$; restoration is
what increases it; keystone loss is what causes it to drop sharply rather
than gradually.
\end{definition}

Health, so defined, does not measure how much of an ecology is currently
running. It measures how much could still be regenerated if some further
part of what is currently running were lost --- a forward-looking
quantity, in keeping with regenerative equivalence being a forward-looking
notion. That phrasing has so far described what health tracks without
saying precisely what an act of repair does to it; the reachable-set
construction of \cref{def:operator-induced-reach} now makes that
precise, and identifies $\Gamma(\mathcal{E})$ above with Chapter 5's
admissible future region under the natural reading $\Gamma(\mathcal{E})
= \mathcal{R}_\mathcal{E}(x)$ for the ecology's current state $x$.

\begin{definition}[Reachability-restoring repair]
\label{def:reach-restoring-repair}
Let $\Gamma^*(x)$ be a target admissible future region and $\Gamma_d(x)$
the region available after damage $d$. An operator $r$ is a
\emph{reachability-restoring repair} when $\mu(\Gamma_{r(d)}(x)) >
\mu(\Gamma_d(x))$, and \emph{complete} relative to the target when
$\Gamma_{r(d)}(x) = \Gamma^*(x)$, partial otherwise.
\end{definition}

\begin{claim}[Repair Existence as Reachability]
\label{thm:repair-existence-reach}
A damaged system admits repair relative to target $\Gamma^*(x)$ if and
only if there exists $r \in \mathcal{O}_\mathcal{E}$ such that
$\mu(\Gamma_{r(d)}(x)) > \mu(\Gamma_d(x))$.
\end{claim}

\begin{proof}
Immediate from \cref{def:reach-restoring-repair}: the forward direction
supplies such an operator by definition; the reverse direction
identifies any executable operator satisfying the inequality as a
repair.
\end{proof}

\Cref{thm:repair-existence-reach} is close to definitional, but it fixes
the vocabulary the rest of this book needs: admissibility describes the
future region; damage contracts it; repair restores some of it. A further
distinction, between fixing a structure and restoring the ecology that
fixes structures, is worth making formal rather than leaving implicit.

\begin{definition}[Regenerative operator]
\label{def:regenerative-operator}
An operator $g$ is \emph{regenerative} for $\mathcal{E}$ when it
increases not merely the current admissible future region but the
ecology's future capacity to produce reachability-restoring repairs:
$\Reach(g(\mathcal{E})) \supsetneq \Reach(\mathcal{E})$, or more
generally $\mu(\mathcal{R}_{g(\mathcal{E})}(x)) >
\mu(\mathcal{R}_\mathcal{E}(x))$ under plausible subsequent damage.
\end{definition}

An operator satisfying \cref{def:reach-restoring-repair} fixes a
structure. An operator satisfying \cref{def:regenerative-operator}
restores the ecology that fixes structures --- retraining an apprentice
rather than repaving a road. This is the formal difference behind
\cref{sec:reframe}'s claim that structures are frozen operator ecologies:
a repair changes $\Gamma_d(x)$; a regeneration changes what $\Reach(\mathcal{E})$
itself can produce, which is a strictly higher-leverage act whenever it
succeeds, and exactly the kind of act a keystone-operator loss under
\cref{rem:keystone-divergence} disables first.

\section{History, Compression, and Repairability}
\label{sec:history}

The preceding sections describe an operator ecology as it stands at a
moment: which operators are executable, how they depend on one another,
how much future capacity they jointly support. This section adds a
dimension the preceding ones leave implicit: an operator does not appear
from nowhere. In every case, it is the compressed residue of a history of
attempts, corrections, and successes.

\begin{definition}[Compression]
\label{def:compression}
Let $H$ denote a history: an ordered trajectory of states, attempts,
corrections, and outcomes. Let $C$ be a compression map taking a history
to an operator, $C(H) = F$, where $F$ is characterized independently of
$C$ by its behavior --- what it does when invoked --- irrespective of how
$F$ came to exist.
\end{definition}

Once $F$ exists, two different things can be done with it. It can be
\emph{enacted}: invoked and executed on some occasion. And it can, where
defined, be \emph{represented}: given a symbolic, inspectable description
sufficient to reconstruct its behavior without executing it. These are
not the same capacity, and the domain of enactment is not in general
contained in the domain of representation --- an operator can be fully
executable while admitting no complete symbolic account of what it does,
in the sense Polanyi's tacit knowledge already names. This is the same
gap \cref{def:extinction} exploited in the other direction: there,
representation survived without enactment; here, enactment can exceed
what representation is available to capture at all.

The central result of \emph{History Before Function} bears directly on
repair capacity, and it is worth stating the dependency honestly rather
than re-deriving it twice: Chapter 5 develops the same result from this
book's own partition apparatus, rather than importing it by analogy, and
this section defers to that proof instead of duplicating it.

\begin{definition}[Repair entropy]
\label{def:repair-entropy}
For an operator $F$ and a retained trace $T$, $\SR(F \mid T)$ measures
the number of distinct fault explanations consistent with a symptom $s$
and the information retained in $T$. Chapter 5 gives this a full
partition-based construction (Fault hypothesis space, repair entropy as
conditional Shannon entropy) rather than the informal version stated
here.
\end{definition}

\begin{claim}[Anti-Alignment, proved in Chapter 5]
\label{claim:antialign}
Let $T_{\min}$ be the minimal trace sufficient to enact $F$ at all. Then
$\SR(F \mid T) \leq \SR(F \mid T_{\min})$ for every retained trace $T$
that contains $T_{\min}$, strictly whenever some fact in $T \setminus
T_{\min}$ excludes an otherwise-live fault location. This is Chapter 5's
Anti-Alignment result, proved there as a corollary of its Trace
Refinement Theorem; this chapter uses the result but does not reprove
it.
\end{claim}

The consequence for operator ecologies is direct. An ecology under
pressure to compress --- to retain only what is needed to keep operators
running, discarding whatever does not affect current executability ---
is, by \cref{claim:antialign}, systematically discarding exactly the
information that would have been needed to repair those operators once
they break. Efficient enactment and future repairability are not merely
different goals; they sit, by construction, in tension with each other,
and an ecology optimized purely for the former is an ecology that has
made itself harder to fix.

\begin{remark}[Why this matters for Chapter 7]
\label{rem:bridge-to-ch6}
Chapter 7 introduces a history-indexed restorability boundary,
$\Rboundary[(X,\mathcal{E},\mathcal{R},H_t)]$, on the grounds that
diagnosing recurrent, worsening, or structurally stubborn conditions
requires a remembered repair history rather than a single snapshot. This
chapter's Anti-Alignment result is the reason that requirement is not
merely convenient but load-bearing: an ecology that retains only
$T_{\min}$ for each of its operators has, in effect, already discarded
the $H_t$ that Chapter 7 needs, not through neglect but as an
unavoidable consequence of having compressed efficiently. A repair theory
that only asks whether an operator exists, without asking how much of its
generating history survives alongside it, will systematically overstate
the repair capacity of a maximally compressed ecology.
\end{remark}

\begin{remark}[The bridge from repertoire to admissible futures, in one place]
\label{rem:bridge-diagram}
Three constructions have now been built across this chapter and the
next two, and it is worth naming the chain that connects them once,
rather than leaving a reader to reassemble it from separate remarks.
\[
\Reach(\mathcal{E}) \;\longrightarrow\; \mathcal{R}_\mathcal{E}(x) \;\longrightarrow\; \Gamma_t(x).
\]
$\Reach(\mathcal{E})$ is this chapter's executable repertoire: a set of
operators, nothing more. $\mathcal{R}_\mathcal{E}(x)$, this chapter's
operator-induced reachable set, is the closure of that repertoire under
composition: the states actually reachable from $x$ by chaining
executable operators together. $\Gamma_t(x)$, Chapter 5's admissible
future region, is a further restriction of a reachable set --- there,
$\overline{\mathcal{R}}_t(x)$, the effective reachable set under a
distinguishability quotient --- to the states an observer's current
admissibility structure actually licenses. Repertoire, then generated
futures, then admissible futures: each stage is a genuine restriction of
the one before it, not a synonym for it, and the type-test discipline
used throughout this book applies to all three pairwise: none are shown
identical, each is shown to generate or restrict its neighbor.
\end{remark}

\section{A Worked Example: Misalignment, Extinction, Regeneration}
\label{sec:wittgenstein}

Chapter 1 closed with a remark that a distinction-preserving
transformation need not preserve operational competence, and deferred
the full illustration to this chapter. A single narrative, developed in
three stages, separates three phenomena this chapter's apparatus keeps
distinct but that are easy to confuse in practice: an intact repertoire
that no longer aligns with its target space, a repertoire with pieces
genuinely and permanently removed, and the appearance of a new ecology
that restores generative capacity rather than merely patching what was
lost.

\subsection*{Stage one: misalignment}

Consider a keyboard instrument relabeled so that pitch runs high to low
across the keys rather than low to high --- every note distinguishable
from every other exactly as before, no fiber of the original
distinguishability space collapsed, nothing lost in the sense Chapter 1
measures. A trained performer's fingering repertoire, $\Reach(\mathcal{E})$,
is entirely intact: every motion the performer knows how to execute is
still executable. What fails is the correspondence between that
repertoire and the relabeled target space. A motion trained to produce
a rising line now produces a falling one; a leap trained to reach a
particular interval now reaches a different one. No operator has been
lost, and no operator has become inexecutable in the sense
\cref{def:extinction} requires. The repertoire and the space it acts on
have simply come apart.

\begin{remark}[Misalignment is neither extinction nor repair deficit]
\label{rem:misalignment}
Misalignment is not covered by any failure mode developed so far. It is
not extinction, since every operator remains executable. It is not
repair deficit in Chapter 7's sense, since $R(d,x)$ for the relabeling
``deviation'' is not the right question --- there is no damage to
repair, only a representation to relearn. It is a third phenomenon:
$\Reach(\mathcal{E})$ intact, $\mathcal{R}_\mathcal{E}(x)$ well defined,
but no longer usefully predictive of what a given operator, applied by
a performer trained on the old correspondence, will actually produce.
\end{remark}

\subsection*{Stage two: extinction}

The pianist Paul Wittgenstein lost his right arm in the First World
War. This is extinction in the precise sense of \cref{def:extinction}:
an entire sub-repertoire of his trained operators, everything requiring
two hands, became permanently inexecutable. Nothing was relabeled;
nothing needed relearning in the sense of stage one. A substantial
fraction of $\Reach(\mathcal{E})$ simply ceased to exist,
$\Reach(\mathcal{E}') \subsetneq \Reach(\mathcal{E})$, and by
\cref{def:generative-keystone} the loss carried generative impact far
beyond the raw count of operators removed: an enormous fraction of the
existing piano repertoire, and the compositional conventions behind it,
assumed two independent hands, so the loss disconnected far more of
$\mathcal{R}_\mathcal{E}(x)$ than the operator count alone would
suggest.

\subsection*{Stage three: regeneration}

Wittgenstein's response was not to search for a repair operator that
would restore two-handed capacity --- none existed --- nor to make do
with the existing one-handed repertoire, which was thin and largely
pedagogical. He commissioned an entirely new body of concert repertoire
for the left hand alone, from composers including Maurice Ravel, Sergei
Prokofiev, Benjamin Britten, and Paul Hindemith. The result was not a
set of individual pieces compensating for specific missing
capabilities. It was a new ecology: new techniques for implying two
independent voices with one hand, new pedagogical traditions, a new
compositional vocabulary specifically suited to the constraint, each
piece extending what the next commission could attempt.

\begin{claim}[The commissioned repertoire is a regenerative operator]
\label{claim:wittgenstein-regenerative}
The commissions satisfy \cref{def:regenerative-operator} rather than
\cref{def:reach-restoring-repair} alone: each addition to the
repertoire increased not merely the immediately available performable
material but the ecology's own capacity to generate further
left-hand-idiomatic material, in the sense that later commissions could
draw on techniques the earlier ones had established. This is the
distinction between fixing a structure and restoring the ecology that
fixes structures, realized in a single, real, well-documented case
rather than illustrated hypothetically.
\end{claim}

\begin{remark}[Why the progression matters]
The three stages are easy to conflate --- all three involve ``a pianist
whose usual technique stops working'' --- and this chapter's vocabulary
is precisely what keeps them apart. Misalignment leaves the repertoire
untouched and breaks only its correspondence to a relabeled space.
Extinction removes part of the repertoire outright, with no
relabeling involved. Regeneration is neither a fix nor a loss but the
appearance of new generative capacity in response to one. A theory of
repair that could not distinguish these three would not yet have
earned the right to a general vocabulary of operators and ecologies;
that it can distinguish them, using nothing beyond the definitions
already given, is this example's actual claim.
\end{remark}

\section{Toward the Restorability Boundary}

This chapter has argued that structures persist because operator
ecologies persist, that the operative notion of sameness under damage is
regenerative rather than behavioral equivalence, and that an ecology's
repair capacity depends not only on which operators remain executable but
on how much of their generating history has been retained alongside
them. None of this guarantees that repair capacity is unlimited. An
ecology can lose a keystone operator outright; it can retain operators
whose generating history has been compressed past the point of
repairability; it can face a rate of damage that outpaces its capacity to
respond regardless of how well-resourced any individual repair is. The
next chapter is about what happens at that limit: not whether repair
occurs, but where and why it stops.

\begin{openproblem}
This chapter has not formally unified the Level 3 rule-update dynamics of
\cref{def:levels}, $F_{t+1} = G(F_t, x_t, E_t)$, with the executable
reachability apparatus of \cref{def:keystone} and Chapter 7's
$\Reach(\mathcal{O},x)$. Both describe an operator changing in response to
conditions, but one is a dynamical update rule and the other is a
set-valued reachability relation, and the relationship between $G$ and
the repair operators $o \in \mathcal{O}$ that realize it has not been made
precise.
\end{openproblem}

\chapter{The Restorability Boundary}
% ------------------------------------------------------------------

\section{Introduction}

A system under observation is rarely either fully known or fully unknown.
More often it occupies an intermediate condition in which some deviations
from its expected trajectory can be detected, diagnosed, and corrected,
while others cannot. The boundary between these two conditions is the
subject of this chapter. It is not a single mechanism but a coordinate at
which several distinct mechanisms can independently place a system beyond
the reach of its surrounding corrective ecology. Distinguishing those
mechanisms from one another is the chapter's primary task, since each
implies a different diagnosis of failure and a different remedy, and
conflating them risks recommending a remedy that addresses the wrong
mechanism entirely.

Restorability should not be understood as the return of a system to some
prior, passively held state. In many of the systems this chapter draws on,
the maintained state is itself an active process rather than a default
condition. A neuron's resting potential is not simply what the membrane
does when left alone; it is continuously re-established against a
thermodynamic gradient by the sodium-potassium pump, which exchanges three
sodium ions outward for two potassium ions inward at continuous metabolic
cost. The resting potential is therefore not a passive baseline but a
continuously repaired state, and if the pump's energy supply fails, the
gradient does not remain stable while some other fault develops elsewhere:
the distinction the gradient exists to maintain simply dissolves. The
absence of repair, in such systems, does not merely block recovery from a
prior failure. It can dissolve the very state that repair was maintaining.
This is one reason the four failure modes developed below are given
concrete non-computational instances wherever possible, rather than being
treated only through the stellar-reconstruction case study: restorability
is a property of maintained physical and biological systems as much as it
is a property of inferential ones.

This should not be overstated into a claim that all persistence is
maintenance. A crystal lattice, a fossil, or a written inscription
persists largely because it occupies a low-energy or kinetically trapped
configuration, not because some operator continuously re-establishes it
against an active gradient. It is more accurate to say that persistence
generally decomposes into two components, $P = S \cup M$: persistence by
stored energetic stability, $S$, and persistence by active maintenance,
$M$. Systems can sit close to either pole --- a crystal lives almost
entirely in $S$, a membrane potential almost entirely in $M$ --- but many
systems of later interest to this book, including institutional and
civilizational ones, are mixtures rather than pure cases. A DNA sequence
is stored information in the sense of $S$, yet its persistence across
cell divisions depends on active proofreading and repair enzymes in the
sense of $M$; an archive is similarly a stored record that persists only
under continuous physical and institutional maintenance. The interesting
empirical claim, developed further in later chapters, is not that
maintenance replaces storage as the general form of persistence, but that
civilizational and biological structures cluster far more heavily toward
$M$ than intuition, which is calibrated on physically inert objects,
tends to assume.

\section{The Signal--Diagnosis--Repair--Continuation Chain}
\label{sec:chain}

Every corrective process examined in this corpus can be decomposed into
four stages through which a deviation must successfully pass before the
system that produced it is considered restored. A deviation must first be
registered as signal: some observable trace of the deviation must reach
the surrounding ecology at all. That signal must then be diagnosed: the
ecology must be able to distinguish the presence and character of the
deviation from ordinary background variation. A diagnosed deviation must
then be repaired: some operator capable of acting on the diagnosis must
exist and be executable under current conditions. And a locally executed
repair must finally be continued: the corrected local structure must
compose with the surrounding structure in a way that does not reintroduce
the deviation elsewhere or fail to persist once attention moves on.

\begin{definition}[Restorability boundary]
\label{def:boundary}
Given a system $X$ and a corrective ecology $\mathcal{E}$ acting on $X$
through the chain of \cref{sec:chain}, the restorability boundary
$\Rboundary(X,\mathcal{E})$ is the set of states of $X$ beyond which
$\mathcal{E}$ can no longer complete the chain: signal that is generated
but not diagnosable, diagnosis that is accurate but not actionable, or
repair that is executable locally but does not continue globally.
\end{definition}

The restorability boundary is deliberately defined without reference to
which stage of the chain fails. This is the sense in which it is the
general object: a system found beyond $\Rboundary$ is uncorrectable, but
the reason it is uncorrectable is left open until one of the four failure
modes below is identified as the operative cause.

\section{Two Kinds of Restorability Boundary}
\label{sec:two-kinds}

\Cref{def:boundary} as stated is under-specified in two ways that matter
before any of the four failure modes can be applied responsibly. Both
corrections are adopted here from \emph{Persistent Anomalies and the
Geometry of Ontology Revision} (Flyxion, 2026), a standalone monograph
developing a general theory of repair algebras and discrepancy hierarchies.
This chapter cites that work rather than absorbing it: what follows imports
exactly the two primitives needed to make \cref{def:boundary} well posed,
and leaves the monograph's repair-algebra formalism, its discrepancy-type
hierarchy, and its Persistence--Obstruction Schema as further reading
rather than reproducing them here.

\subsection{The architecture gate}

The first under-specification is that \cref{def:boundary} does not
distinguish a target that remains representable in principle from a
target that cannot be represented by the current architecture at all.
Let $\pi : \mathcal{M} \to \Sigma$ be the prediction map of the
representational architecture in use, so that $\operatorname{Im}\pi
\subseteq \Sigma$ is the set of observations that architecture can in
principle produce.

\begin{definition}[Architecture gate]
\label{def:gate}
For an observation $x_{\mathrm{obs}} \in \Sigma$, the architecture gate is
\[
\chi_\pi(x_{\mathrm{obs}}) \;=\;
\begin{cases}
0, & x_{\mathrm{obs}} \in \operatorname{Im}\pi, \\
1, & x_{\mathrm{obs}} \notin \operatorname{Im}\pi.
\end{cases}
\]
A restorability boundary encountered at $\chi_\pi = 0$ is an
\emph{inferential boundary}: the target remains representable, but the
available observer or repair ecology cannot presently recover it. A
restorability boundary encountered at $\chi_\pi = 1$ is an
\emph{ontological boundary}: no admissible trajectory in the current
representational architecture fits the observation, and no amount of
additional repair within that architecture will resolve it.
\end{definition}

This distinction matters because the four failure modes developed later
in this chapter --- more throughput, more diagnostic capacity, more
operators, better global composition --- are all remedies internal to a
fixed architecture. None of them can close an ontological boundary. A
system that persistently treats $\chi_\pi = 1$ conditions as if they were
$\chi_\pi = 0$ conditions will apply an unbounded amount of the wrong kind
of repair, which is precisely the trap the monograph's central historical
example is chosen to illustrate.

\begin{example}[Mercury and Banik: an ontological/inferential contrast]
\label{ex:mercury}
The nineteenth-century anomaly in Mercury's perihelion precession survived
every repair available within Newtonian celestial mechanics: hypothesized
planets, atmospheric drag corrections, refined solar oblateness estimates.
The residual was not merely undiagnosed; it was $x_{\mathrm{obs}} \notin
\operatorname{Im}\pi_{\mathrm{N}}$, where $\pi_{\mathrm{N}}$ is the
Newtonian prediction map --- no parameter choice within Newtonian
mechanics could produce the observed precession. This is an ontological
boundary, and it was resolved not by a stronger repair within Newtonian
mechanics but by general relativity expanding the representational
architecture itself. By contrast, the Banik reconstruction of
\cref{ex:banik} is an inferential boundary only: the age-likelihood
architecture is fixed throughout, and the crossover to the noise floor
marks a point at which that specific reconstruction ceases to carry
reliable information, not a demonstrated inadequacy of the underlying
representational framework. It should be stated carefully that the noise
floor is evidence for an inferential boundary; it is not evidence against
an ontological one, since this chapter has not attempted to show that no
enlarged inference within a similar general representation could recover
more. The two cases are worth holding side by side precisely because they
present a superficially similar symptom --- a residual that will not go
away --- for two structurally unrelated reasons.
\end{example}

\subsection{History indexing}

The second under-specification is that \cref{def:boundary} treats
unrestorability as an instantaneous condition. But whether a deviation is
a first encounter, a repeated failure, a temporarily exhausted capacity,
or a structurally stubborn residual cannot be determined from a single
snapshot; it requires a remembered record of what has already been tried.

\begin{definition}[History-indexed restorability boundary]
\label{def:boundary-history}
Let $H_t = (\varepsilon_0, R_1, \varepsilon_1, \ldots, R_n, \varepsilon_n)$
denote the remembered sequence of residuals and attempted repairs up to
time $t$, for a repair algebra $\mathcal{R}$ available to the ecology
$\mathcal{E}$. The history-indexed restorability boundary is
\[
\Rboundary[_{t}] \;=\; \Rboundary(X, \mathcal{E}, \mathcal{R}, H_t).
\]
\end{definition}

Without $H_t$, an ecology can detect that a deviation is currently
unrestorable, but it cannot distinguish a first encounter from a repeated
failure, a temporary capacity exhaustion from a persistent obstruction, or
an untried repair from an invariant residual --- distinctions the four
failure modes below depend on whenever they are used to diagnose
\emph{recurrent, worsening,} or \emph{structurally stubborn} conditions
rather than momentary ones. The four modes are accordingly understood
throughout the rest of this chapter as evaluated relative to $H_t$: the
Expiatory Gap becomes $G_t = D_{\mathrm{sys}}(H_t) - D_{\max}(H_t)$;
diagnostic collapse becomes the history-conditioned mutual information
$\I_t(F;E \mid H_t) \to 0$; repair deficit records which operators have
already been attempted, found unavailable, or rendered non-executable
within $H_t$; and continuation failure records whether repeated local
sections fail to glue across changing covers as the history unfolds. This
chapter does not revise the four definitions given below to make this
history-dependence explicit in every equation, since doing so before the
underlying repair-history formalism is developed further would overstate
what has actually been established; it is stated here as the governing
convention under which those definitions should be read.

A history-indexed boundary is not yet a diagnostic procedure. Nothing so
far says how, from $H_t$ alone, an ecology should actually conclude that
a residual is persistent rather than merely unresolved so far. That
inferential step needs its own construction.

\begin{definition}[Residual orbit and persistent residual]
\label{def:residual-orbit}
Let $\mathcal{R}$ be a repair algebra acting on a discrepancy space
$\mathcal{D}$. For a repair word $w = R_n \circ \cdots \circ R_1 \in
\langle \mathcal{R} \rangle$, write $\varepsilon_w = w(\varepsilon_0)$.
The repair orbit of $\varepsilon_0$ is $\operatorname{Orb}_\mathcal{R}(\varepsilon_0)
= \{w(\varepsilon_0) : w \in \langle\mathcal{R}\rangle\}$. For a
discrepancy norm $\|\cdot\|$, define $P(\varepsilon_0,\mathcal{R}) =
\inf_{w \in \langle\mathcal{R}\rangle} \|w(\varepsilon_0)\|$. The
discrepancy is \emph{repairable to tolerance $\eta$} when
$P(\varepsilon_0,\mathcal{R}) \leq \eta$, and \emph{persistent above
tolerance $\eta$} when $P(\varepsilon_0,\mathcal{R}) > \eta$.
\end{definition}

$H_t$ records only attempted words, not every word in $\langle\mathcal{R}\rangle$,
so it can only approximate $P(\varepsilon_0,\mathcal{R})$ from above.

\begin{definition}[Empirical persistence bound]
\label{def:empirical-persistence}
For a history $H_t$ containing attempted repair words $w_1,\ldots,w_n$,
define $\widehat{P}_t = \min_{1 \leq k \leq n} \|w_k(\varepsilon_0)\|$.
\end{definition}

\begin{claim}[Monotonicity of the empirical persistence bound]
\label{thm:empirical-persistence}
As repair history grows, $\widehat{P}_{t+1} \leq \widehat{P}_t$, and
$P(\varepsilon_0,\mathcal{R}) \leq \widehat{P}_t$ for every finite
history.
\end{claim}

\begin{proof}
The attempted set at $t+1$ contains the set attempted by $t$, so a
minimum over the larger set cannot exceed the minimum over the smaller
one. Since the full repair monoid contains every attempted word, the
infimum over the full monoid cannot exceed the minimum over any finite
attempted subset.
\end{proof}

\Cref{thm:empirical-persistence} exposes a limitation worth stating
plainly: \emph{repeated failure does not by itself prove structural
persistence unless the attempted repair family is sufficiently
exhaustive}. ``Every repair tried so far has failed'' and ``every repair
available in the architecture must fail'' are different claims, and only
a completeness condition closes the gap between them.

\begin{claim}[Persistence Identification under Exhaustive Search]
\label{thm:exhaustive-persistence}
Suppose the attempted repair words become dense in $\langle\mathcal{R}\rangle$
under a topology for which $w \mapsto \|w(\varepsilon_0)\|$ is
continuous. Then $\widehat{P}_t \to P(\varepsilon_0,\mathcal{R})$.
\end{claim}

\begin{remark}[Sharpening Mercury and Banik]
\label{rem:mercury-banik-sharpened}
This distinction refines \cref{ex:mercury}'s contrast. The Newtonian
astronomers who exhausted planetary hypotheses, drag corrections, and
oblateness estimates against Mercury's perihelion were not merely
accumulating $\widehat{P}_t$; the historical record suggests their
attempted repair family came close to the density condition of
\cref{thm:exhaustive-persistence} within the Newtonian algebra
$\mathcal{R}_{\mathrm{N}}$, which is what licenses reading
$P_{\mathcal{R}_{\mathrm{N}}}(\varepsilon_{\mathrm{Mercury}}) > 0$ as a
genuine ontological boundary rather than an artifact of insufficient
search. Nothing comparable has been shown for the Banik reconstruction:
its noise floor reflects $\widehat{P}_t$ for one specific reconstruction
procedure, with no argument that the attempted procedures are dense in
any relevant sense, which is precisely why \cref{ex:mercury} was careful
to call it an inferential boundary only.
\end{remark}

\begin{remark}[Why $H_t$ is load-bearing, not merely convenient]
\label{rem:antialign-bridge}
Chapter 6 gives a reason $H_t$ cannot be treated as an optional refinement
that a well-run ecology could reasonably do without. Its Anti-Alignment
result (Chapter 6, Claim 5.7, imported there from \emph{History Before
Function}) shows that the minimal trace $T_{\min}$ needed merely to
\emph{enact} an operator is generically the worst-positioned trace for
\emph{repairing} that operator later, strictly worse whenever the
discarded surplus carried diagnostic information. An ecology under
ordinary pressure to compress --- to retain only what current
executability requires --- is therefore not neutral with respect to
$H_t$: it is systematically discarding exactly the record this chapter's
history-indexed boundary depends on, as an unavoidable consequence of
efficient operation rather than through neglect. This sharpens what
``without $H_t$'' means in the paragraph above. It is not merely that an
ecology which happens not to keep records cannot make the finer
diagnostic distinctions listed there; it is that keeping the enactment of
its operators efficient and keeping $H_t$ available are, by Chapter 6's
result, in tension by construction. A history-indexed restorability
boundary is therefore not a bookkeeping addition to an otherwise complete
theory of repair. It names a resource --- retained, undiscarded history
--- that ordinary operation continuously depletes.
\end{remark}

\begin{remark}[What remains open]
\label{rem:open-relation}
This section does not claim that the four failure modes of
\cref{sec:four-modes} reduce to, or are systematically organized by, the
architecture gate or the discrepancy-type hierarchy of
\emph{Persistent Anomalies}. The relation between the two classifications
remains open. An in-architecture numerical error can still produce an
Expiatory Gap if observations simply arrive faster than they can be
processed; an out-of-architecture categorical mismatch can remain highly
discriminable and therefore need not involve diagnostic collapse at all; a
transient discrepancy can encounter a temporary repair deficit because the
relevant operator exists but happens not to be executable at the moment
it is needed. None of these cases is naturally derived from a fixed cell
of $\{\chi_\pi\} \times \{\text{discrepancy type}\} \times \{\text{persistence
character}\}$, which suggests the four modes track operational coordinates
---demand versus capacity, discriminability, operator availability,
local-to-global composability--- that may be orthogonal to, rather than
subordinate to, the architecture-level classification. Establishing
whether the two classifications are orthogonal, partially reducible, or
connected only under additional assumptions is left as a separate
comparison problem.
\end{remark}

The four modes' mutual independence, established above at the level of
definitions, does not mean they cannot cause one another dynamically. One
such coupling admits a direct formal treatment.

\begin{claim}[Throughput-Induced Information Bound]
\label{thm:throughput-info}
Let signal arrive at rate $D_{\mathrm{sys}}$ while the diagnostic ecology
processes at most $D_{\max}$, and let $\rho = \min(1, D_{\max}/D_{\mathrm{sys}})$
be the retained fraction. Under independent random thinning at retention
probability $\rho$ --- discarded items uncorrelated with which
observations are diagnostically valuable --- and conditional independence
of feedback items, $\I(\widetilde{F};E) \leq \rho \, \I(F;E)$ for the
thinned feedback $\widetilde{F}$. Consequently, as $D_{\max}/D_{\mathrm{sys}}
\to 0$, $\I(\widetilde{F};E) \to 0$.
\end{claim}

\Cref{thm:throughput-info} gives a first formal coupling arrow, Expiatory
Gap $\Rightarrow$ Diagnostic Collapse under lossy thinning, while
preserving their analytical distinction: the reverse implication does not
hold, since $\I(F;E) \to 0$ does not imply $D_{\mathrm{sys}} > D_{\max}$
--- a low-volume but systematically misleading channel can collapse
diagnostically with no throughput problem at all. This resolves one
direction of one coupling among the four modes; the remaining couplings,
and the relation to the architecture gate raised above, are not addressed
by this result and remain open.

\section{Four Failure Modes}
\label{sec:four-modes}

\subsection{Expiatory Gap: throughput failure}
\label{subsec:expiatory-gap}

The Expiatory Gap names the condition in which a system generates signal,
complexity, or consequence faster than the surrounding ecology can absorb
and process it, independent of whether that ecology could in principle
diagnose the signal correctly given sufficient time.

\begin{definition}[Expiatory Gap]
Let $D_{\mathrm{sys}}$ denote the rate at which a system generates
diagnosable content and $D_{\max}$ denote the maximum rate at which the
receiving ecology can process such content while retaining diagnostic
reliability. The Expiatory Gap is
\[
G \;=\; D_{\mathrm{sys}} - D_{\max}.
\]
\end{definition}

\begin{claim}
\label{claim:expiatory-remedy}
The Expiatory Gap is, in principle, closable by rate control alone: any
combination of reduced emission rate, increased receiver bandwidth,
staging, buffering, or temporal dilation that brings effective throughput
within $D_{\max}$ restores the diagnostic reliability of the channel.
\end{claim}

\Cref{claim:expiatory-remedy} is the load-bearing assumption behind the
``natural damping'' argument developed for this failure mode, and it is
exactly the assumption that fails for the next failure mode.

\subsection{Diagnostic Collapse: discriminability failure}
\label{subsec:diagnostic-collapse}

Diagnostic collapse names the condition in which a feedback channel
remains active, and may even remain high in volume, while ceasing to
carry information about the error state it is meant to track.

\begin{definition}[Diagnostic collapse]
Let $F$ denote feedback supplied by a corrective ecology and $E$ denote
the true error state of the system under correction. The channel is
diagnostically active when $\I(F;E) > \varepsilon$ for some operationally
meaningful threshold $\varepsilon$. Diagnostic collapse is the condition
\[
\I(F;E) \;\longrightarrow\; 0
\]
independent of the volume, frequency, or apparent sophistication of $F$.
\end{definition}

\begin{claim}
\label{claim:diagnostic-remedy}
Diagnostic collapse is not closable by rate control. Increasing the
volume of feedback, or the time allotted to produce it, leaves
$\I(F;E)$ unchanged when the feedback source lacks the capacity to
discriminate correct from incorrect states. Closing the gap requires a
qualitatively new diagnostic channel: additional expertise, a new class of
observation, or a restructured overlap graph between the system's active
frontier and the population capable of checking it.
\end{claim}

\begin{example}[The oldest star as a diagnostic-collapse boundary]
\label{ex:banik}
Banik et al.\ (2026) reconstruct the latent age distribution $P(A)$ of the
oldest stars in a sample of 155{,}600 Milky Way subgiants, using the
individual age likelihoods $\mathcal{L}_i(A)$ obtained from LAMOST DR7
spectroscopy and Gaia astrometry. The reconstructed $P(A)$ never reaches
zero at any finite age, since positivity constraints and finite sampling
guarantee a nonzero posterior at every grid point. The age of the oldest
star, $A_\star$, is instead located by examining the ratio $P(A)/\sigma(A)$
across Markov Chain Monte Carlo samples: below some age this ratio
declines approximately linearly, and beyond it the ratio flattens into a
noise floor that does not further decline no matter how the reconstruction
is refined. The crossover between these two regimes, found at
$A_\star = 13.73^{+0.18}_{-0.15}$ Gyr, is a directly measured instance of
\cref{def:boundary}: not a point at which signal ceases, but a point at
which structured variation in the reconstruction gives way to variation
indistinguishable from the reconstruction process's own uncertainty.
\end{example}

The Banik example is worth stating as a general principle in its own
right, since it is the clearest available instance of the distinction this
chapter depends on.

\begin{claim}[The boundary is discriminability, not disappearance]
\label{claim:boundary-is-discriminability}
The restorability boundary is generally not the point at which a signal
disappears, but the point at which the distinction between signal and
reconstruction uncertainty ceases to be recoverable. A magnitude threshold
asks whether a quantity has fallen to zero; a discriminability threshold
asks whether variation in that quantity can still be distinguished from
the noise generated by the process used to recover it. The two coincide
only by coincidence.
\end{claim}

The four failure modes of this chapter are analytically distinct but not
claimed to be dynamically independent, and a taxonomy that implied
otherwise would misdescribe how these failures actually arise in practice.
Sustained throughput overload can itself induce diagnostic collapse, since
an ecology forced to process signal faster than $D_{\max}$ for long enough
may lose the capacity to sort useful content from noise even after the
overload subsides. Conversely, diagnostic collapse can present as apparent
throughput overload, since an ecology that can no longer discriminate
relevant from irrelevant signal will tend to treat a constant volume of
input as increasingly unmanageable, not because more of it is arriving but
because less of what arrives can be filtered. The four modes should
therefore be read as four independent quantities that can each run out,
any of which may trigger or mask the depletion of another, rather than as
four mutually exclusive diagnoses of a single failure.

\subsection{Repair Deficit: operator failure}
\label{subsec:repair-deficit}

Repair deficit names the condition in which diagnosis succeeds --- the
deviation is correctly identified and localized --- but no operator
capable of acting on that diagnosis exists or remains executable.

\begin{definition}[Repair deficit]
A diagnosed deviation $d$ at coordinate $x$ suffers a repair deficit when
the operator ecology $\mathcal{O}$ contains no operator $o \in \mathcal{O}$
such that $o$ is executable at $x$ and $o$ corrects $d$.
\end{definition}

This definition is formally correct but comparatively thin next to the
quantities available for the other three failure modes: a rate for the
Expiatory Gap, a mutual information for diagnostic collapse, a gluing
kernel for continuation failure. A more comparable formulation defines an
executable reachability relation over the operator ecology directly.

\begin{definition}[Executable reachability]
\label{def:reach}
Let $\operatorname{Reach}(\mathcal{O}, x)$ denote the set of deviations
correctable from coordinate $x$ by some composition of operators in
$\mathcal{O}$ that remain executable at $x$ under current conditions.
Repair deficit for a diagnosed deviation $d$ at $x$ is then the condition
\[
d \;\notin\; \operatorname{Reach}(\mathcal{O}, x).
\]
\end{definition}

\Cref{def:reach} makes repair deficit directly comparable to the other
three failure modes and, more importantly, makes explicit that repair
deficit is a property of the pair $(\mathcal{O}, x)$ rather than of $d$
alone: the same deviation can be a repair deficit at one coordinate and
fully correctable at another, depending on which operators remain
executable there.

Repair deficit is the point at which this chapter's chain interfaces
directly with the operator-ecology framework of Chapter 6. A correctly
diagnosed deviation that outpaces available operators
is, in ecological terms, the loss of a keystone operator: the niche is
identified with precision, but nothing remains to occupy it. This
suggests that repair deficit should ultimately be treated jointly with
that material rather than developed independently here.

\Cref{def:reach} treats ``executable at $x$'' as a primitive, but repair
deficit can arise for at least two structurally different reasons that
this primitive conflates. Let $\operatorname{Exec}(x) \subseteq
\mathcal{O}$ denote the operators that can actually run at coordinate $x$
under current conditions.

\begin{definition}[Structural and energetic repair deficit]
\label{def:deficit-types}
A deviation $d$ at $x$ suffers a \emph{structural} repair deficit when no
operator capable of correcting $d$ exists in the ecology at all, that is,
there is no $o \in \mathcal{O}$ that corrects $d$ regardless of
executability. It suffers an \emph{energetic} repair deficit when a
corrective operator exists but cannot run at $x$: there is some $o \in
\mathcal{O}$ that corrects $d$, yet $o \notin \operatorname{Exec}(x)$.
\end{definition}

\begin{example}[Energetic repair deficit: the sodium-potassium pump]
\label{ex:pump}
The neuronal membrane potential is not a passive baseline but a
continuously repaired state, maintained by the sodium-potassium pump
exchanging three sodium ions outward for two potassium ions inward at
continuous metabolic cost. The corrective operator here is never
structurally absent: the pump exists as long as the cell exists. What can
fail is executability. If ATP production fails, the pump satisfies
$o_{\mathrm{pump}} \in \mathcal{O}$ but $o_{\mathrm{pump}} \notin
\operatorname{Exec}(x)$, and the deviation --- the gradual collapse of the
ion gradient toward equilibrium --- falls out of
$\operatorname{Reach}(\mathcal{O},x)$ not because no operator was ever
capable of correcting it, but because the operator's precondition failed.
This is repair deficit of the energetic kind, and it is worth keeping
conceptually separate from the central-pattern-generator example of
\cref{ex:cpg}, which illustrates the structural side of the same
definition: redundancy among operators that remain executable, rather
than the failure of a single operator's executability.
\end{example}

\begin{openproblem}
\Cref{def:deficit-types} may not be exhaustive. A third candidate, not yet
formalized here, is \emph{kinetic} repair deficit: an operator $o$ with $o
\in \operatorname{Exec}(x)$ that nonetheless fails to correct $d$ because
its completion time exceeds the timescale over which $d$ compounds or
propagates. This would be a recurrence of the Expiatory Gap's rate-mismatch
logic ($D_{\mathrm{sys}} > D_{\max}$ in \cref{subsec:expiatory-gap})
relocated from the diagnostic channel to the repair operator itself, and
its existence raises a question this chapter does not yet resolve: whether
throughput failure is properly a phenomenon unique to the Signal $\to$
Diagnosis boundary, or a failure pattern (rate exceeding capacity) that can
recur at any stage of the chain, of which the Expiatory Gap as defined in
\cref{subsec:expiatory-gap} is only the first instance.
\end{openproblem}

Reachability as stated in \cref{def:reach} is binary: a deviation is
either reachable or it is not. This conceals a distinction that matters
in practice, since a deviation reachable by exactly one operator and a
deviation reachable by ten independent operators are both formally free
of repair deficit, yet occupy very different positions with respect to
future risk. The repair reserve refines \cref{def:reach} accordingly.

\begin{definition}[Repair reserve]
\label{def:reserve}
For a diagnosed deviation $d$ at coordinate $x$, the repair reserve is
\[
R(d,x) \;=\; \bigl|\{\, o \in \mathcal{O} : o \text{ executable at } x,\ o \text{ corrects } d \,\}\bigr|.
\]
Repair deficit as defined in \cref{def:reach} is then the special case
$R(d,x) = 0$, and $d \in \operatorname{Reach}(\mathcal{O},x)$ if and only
if $R(d,x) > 0$. Values of $R(d,x)$ greater than one indicate overlapping,
mutually substitutable corrective pathways rather than a single point of
failure.
\end{definition}

\begin{example}[Central pattern generators and repair reserve]
\label{ex:cpg}
Locomotor gait in most vertebrates is not produced by a single controller
but by a nested hierarchy of corrective operators acting at different
scales: cellular oscillators, local spinal circuits, limb-level
coordination, gait-level coordination, and cortical modulation. Each layer
corrects deviations at its own scale before they can propagate to the
next. Consequently $R(d,x)$ for a typical gait deviation is well above
one even when a high-level operator is lost: an animal that loses
cortical control can often still walk, because lower-level operators
remain executable and continue to correct the deviation. Repair deficit
for gait --- the condition $R(d,x) = 0$ --- typically requires the loss
of operators across several layers simultaneously, not the loss of any
single controller. The example illustrates that healthy biological
systems do not merely avoid repair deficit; they operate with
$R(d,x) \gg 1$ as a matter of course, and it is the erosion of that
reserve, rather than its first crossing of zero, that constitutes the
early warning signal of an approaching repair deficit.
\end{example}

\begin{remark}[Successful repair can still collapse fault identity]
\label{rem:repair-fiber-collapse}
$R(d,x) > 0$ certifies that some executable operator corrects $d$; it
says nothing about whether the act of repairing preserves the identity
of what was wrong. This concern is not idle: a repair operator is
itself a transition in the sense of Chapter 9's framework, and
repair operators built to accommodate anomalous evidence with minimal
disruption generically collapse fibers --- distinct prior faults
$f_1 \neq f_2$ routinely satisfy $R(f_1) = R(f_2)$, precisely because a
minimality criterion is more responsible for the repaired result than
the specific defect that triggered it. Admissibility can therefore be
restored while historical recoverability is lost in the same stroke:
$\text{repair} \not\Rightarrow \text{historical recoverability}$. This
is not a failure the four modes of this chapter currently register,
since a repair satisfying $R(d,x) > 0$ counts as successful by
\cref{def:reach,def:reserve} regardless of whether it preserves fault
identity. It bears directly on the residual-orbit machinery of
\cref{sec:two-kinds}: a repair word $w$ that reduces $\|\varepsilon\|$
to an acceptable tolerance has not thereby shown that the identity of
the original discrepancy survives in $H_t$, and \cref{def:residual-orbit}'s
apparatus tracks the residual's magnitude along the orbit, not whether
successive repairs are individually distinguishable from one another
after the fact. This is recorded here as an open phenomenon this
chapter's taxonomy does not yet have a formal home for, rather than
folded into one of the four existing modes by assumption.
\end{remark}

\begin{principle}[Repair--Recovery Separation]
\label{princ:repair-recovery-separation}
\[
\text{Repair} \;\not\Rightarrow\; \text{Historical Recoverability}.
\]
Restoring admissibility and preserving the identity of what was wrong
are independent achievements. An operator satisfying $R(d,x) > 0$
certifies the first; nothing in this chapter's taxonomy, on its own,
certifies the second. This is named separately from
\cref{rem:repair-fiber-collapse} because it recurs beyond this
chapter: Chapter 6's minimality-driven repair operators and Chapter 9's
fiber-collapsing repair transitions are independent instances of the
same separation, not three unrelated observations that happen to share
a conclusion.
\end{principle}

\subsection{Continuation Failure: composition failure}
\label{subsec:continuation-failure}

Continuation failure names the condition in which local repair succeeds
but the corrected structure does not compose with its surroundings into a
globally coherent whole.

Let $\mathcal{F}$ be a sheaf of local structure over a covered space
$(X, \{U_i\})$, and let each local repair yield a section $s_i \in
\mathcal{F}(U_i)$. An earlier version of this section wrote continuation
failure as the nonvanishing of a kernel from $\prod_i \mathcal{F}(U_i)$ to
$\Gamma(X,\mathcal{F})$, but that arrow is not generally canonical: global
sections restrict to local ones, not the reverse, so a map running the
other way cannot be the right object to take a kernel of. The correct
formulation distinguishes two failure modes the earlier notation
conflated.

\begin{definition}[Disagreement map]
\label{def:disagreement}
Define $\delta : \prod_i \mathcal{F}(U_i) \to \prod_{i,j} \mathcal{F}(U_i
\cap U_j)$ by $\delta((s_i)_i) = (s_i|_{U_i \cap U_j} - s_j|_{U_i \cap
U_j})_{i,j}$ for an abelian sheaf, or the corresponding pair of
restrictions in the set-valued case.
\end{definition}

\begin{claim}[Local Compatibility Criterion]
\label{thm:local-compatibility}
A family $(s_i)_i$ of local repairs is compatible on overlaps exactly
when $\delta((s_i)_i) = 0$. If $\mathcal{F}$ is a sheaf, every compatible
family glues uniquely to a global section $s \in \Gamma(X,\mathcal{F})$
with $s|_{U_i} = s_i$ for all $i$.
\end{claim}

\begin{proof}
$\delta((s_i)_i) = 0$ states precisely that $s_i$ and $s_j$ agree on
every overlap. Existence and uniqueness of the global section is the
sheaf gluing axiom.
\end{proof}

\begin{definition}[Compatibility failure and descent failure]
\label{def:two-failures}
A family of locally successful repairs exhibits \emph{compatibility
failure} when $\delta((s_i)_i) \neq 0$: the repairs disagree already on
some overlap. For a presheaf not known to satisfy the sheaf axiom, a
compatible family ($\delta((s_i)_i) = 0$) exhibits \emph{descent failure}
when no global section restricts to it. For a genuine sheaf, descent
failure cannot occur --- \cref{thm:local-compatibility} rules it out by
definition, which is why obtaining a nontrivial global obstruction
despite pairwise compatibility requires working with a presheaf, with
higher cocycle data, or with nonlinear constraints not captured by
pairwise equality alone.
\end{definition}

\begin{claim}[Cohomological Continuation Obstruction]
\label{thm:continuation-obstruction}
Suppose local repairs determine a \v{C}ech 1-cocycle $g = (g_{ij}) \in
\check{Z}^1(\mathcal{U},\mathcal{G})$ of transition corrections on
overlaps. A globally coherent repair exists whenever the class $[g] \in
\check{H}^1(\mathcal{U},\mathcal{G})$ vanishes. A nonzero class is an
obstruction to continuation.
\end{claim}

\begin{proof}
A global repair corresponds to local reparameterizations $h_i$ satisfying
$g_{ij} = h_j h_i^{-1}$ on each overlap, i.e., to $g$ being a coboundary.
Such $h_i$ exist exactly when the cohomology class of $g$ vanishes.
\end{proof}

Continuation failure in the sense this chapter needs is therefore best
stated as compatibility failure ($\delta((s_i)_i) \neq 0$) for the basic
case, with the cohomological obstruction of
\cref{thm:continuation-obstruction} available for the genuinely global
case where pairwise-compatible local repairs still fail to assemble.

\begin{example}[Cardiac synchronization as continuation failure]
\label{ex:pacemaker}
The heartbeat is not produced by a single clock winding down but by a
population of cellular oscillators in the sinoatrial node continually
re-synchronizing through gap-junction coupling. Each oscillator, taken
locally, is a valid section: it continues to fire rhythmically on its own
terms. A coherent heartbeat requires these local sections to glue into a
single global rhythm through sufficient electrical coupling between cells.
When coupling weakens, arrhythmia can result even though every individual
oscillator remains functional and every local repair, in isolation,
continues to succeed. The failure is not that signal generation stops or
that diagnosis degrades; it is compatibility failure in the sense of
\cref{def:two-failures}: adjacent oscillators' phases disagree across the
coupling that should reconcile them, $\delta((s_i)_i) \neq 0$, and local
success no longer implies global coherence.
\end{example}

\begin{openproblem}
The treatment of the ``existential expiatory gap'' in
sheaf-theoretic terms develops precisely this gluing
obstruction, but files it under the throughput-failure vocabulary of
\cref{subsec:expiatory-gap}. Determine whether that material is more
accurately relocated to this section, and whether its interpretation
changes once continuation failure is treated as a distinct failure mode
rather than a symptom of an unmanaged Expiatory Gap.
\end{openproblem}

\section{Relating the Four Modes}
\label{sec:relating}

\Cref{tab:summary} places the four failure modes side by side. The
central methodological claim of this chapter is that no single remedy
addresses all four: a system exhibiting diagnostic collapse will not be
repaired by measures designed for throughput failure, and a system
exhibiting continuation failure may show perfectly healthy local
diagnosis and repair at every individual site while still failing as a
whole.

\begin{center}
\begin{tabular}{llll}
\hline
Failure mode & Lost quantity & Formal signature & Remedy \\
\hline
Expiatory Gap & Throughput & $G = D_{\mathrm{sys}} - D_{\max}$ & Rate control \\
Diagnostic collapse & Information & $\I(F;E)\to 0$ & New diagnostic capacity \\
Repair deficit & Operators & $d \notin \operatorname{Reach}(\mathcal{O}, x)$ & New operators \\
Continuation failure & Composability & Compatibility failure, $\delta \neq 0$ & Global re-coordination \\
\hline
\end{tabular}
\captionof{table}{The four failure modes of the restorability boundary,
organized by which quantity is exhausted rather than by which stage of
the chain is affected. This framing generalizes beyond this chapter: for
any observed instance of a system failing to recover, the diagnostic
question is which of throughput, information, operators, or composability
is actually scarce, since the answer determines which of the four remedy
classes applies and rules out the other three.}
\label{tab:summary}
\end{center}

\begin{example}[Scale-dependent neural organization]
\label{ex:posani-scale}
Posani et al.\ (2026) provide an instructive case in which the apparent
organization of a system changes with descriptive scale. Analyzing
neural responses across the cortical hierarchy, they find that
individual neurons are rarely cleanly categorical: local response
profiles are heterogeneous and commonly mix several task-related
variables. At the population level, however, those distributed responses
can remain highly separable, so that relevant variables are recoverable
from the collective geometry even though they are not assigned to
sharply specialized local units.

This is not itself a measurement of restorability, and Chapter 5's
comparison already establishes that the paper's participation ratio
should not be identified with this book's reachability quantities. What
the example exposes instead is a scale dependence a theory of
restorability must permit for. At the single-unit scale, the absence of
categorical specialization may look like weak localization: there is no
unique component whose inspection or replacement corresponds to one
represented variable. At the population scale, the same system may
possess substantial diagnostic separability, because information is
distributed across decorrelated response directions rather than
concentrated in any one of them. A failure can therefore appear
differently under different covers or levels of description. Damage to
one locally mixed component need not destroy the population-level
distinction, since other components participate in representing it;
conversely, a perturbation aligned with a low-dimensional population
direction can erase a globally recoverable distinction while leaving
many individual units active and, by any local measure, undamaged.
\end{example}

\begin{remark}[The slogan this example licenses]
\label{rem:posani-slogan}
\cref{ex:posani-scale} supports a compact statement worth having on
hand: restorability depends on separability, not categoricity. A system
does not need to sort its states into clean, labeled categories to
support restoration --- Posani et al.'s own finding is that cortex
mostly does not. It needs enough geometric structure, at whatever scale
diagnosis is actually performed, that the distinctions repair depends on
remain recoverable. Diagnostic collapse (\cref{subsec:diagnostic-collapse})
is the failure of exactly this condition, stated earlier in this chapter
in information-theoretic terms; this example gives it a second,
geometric reading.
\end{remark}

\begin{remark}[A graded refinement of the architecture gate]
\label{rem:graded-gate}
\cref{ex:posani-scale} also sharpens the architecture gate of
\cref{def:gate}. That gate is binary: $\chi_\pi(x) \in \{0,1\}$,
according to whether $x \in \operatorname{Im}(\pi)$. Posani's result
shows the interesting structure can live entirely on the $\chi_\pi = 0$
side of that boundary and still admit further distinctions: a
represented state can be highly separable, weakly separable, or
effectively inseparable from its neighbors, depending on the local
geometry of the representation at the scale it is queried. An
architecture can contain a phenomenon --- $x \in \operatorname{Im}(\pi)$
--- while still lacking the geometry to isolate it from nearby
alternatives. This does not revise \cref{def:gate}; it observes that the
gate answers a coarser question than the one diagnosis usually needs
answered, and that a graded notion of diagnostic accessibility inside
$\operatorname{Im}(\pi)$ is a real further distinction this chapter has
not formalized.
\end{remark}

\begin{claim}[Scale relativity of restorability]
\label{claim:scale-restorability}
A system may lie inside the restorability boundary relative to one
descriptive scale and outside it relative to another. For fine and
coarse observation maps $q_f : X \to X_f$ and $q_c : X \to X_c$ with a
scale map $p : X_f \to X_c$ relating them, the corresponding effective
corrective ecologies $\mathcal{E}_f$ and $\mathcal{E}_c$ need not
satisfy $p(\Rboundary[_f]) = \Rboundary[_c]$, and restorability at one
scale need not imply restorability at the other.
\end{claim}

Coarse-graining can move a system relative to the restorability boundary
in either direction, and for different reasons in each. It can hide a
local failure by merging it with states that remain functionally
equivalent at the coarser scale --- the population-level separability of
\cref{ex:posani-scale}, which survives damage no single fine-scale unit
could survive by itself. Or it can destroy the very distinction a
diagnosis or repair needs, by collapsing exactly the fine-scale
difference the coarse scale was never built to track --- the reverse
failure, in which population-level structure conceals rather than
compensates for a fine-scale problem. Neither direction is the default;
which one occurs depends on which distinctions a given coarsening
removes, which is why \cref{claim:scale-restorability} states an
inequality rather than a direction. The operative restorability boundary
for a system is accordingly not a property of the system alone. It
depends on which scale diagnosis and repair are actually attempted at,
a dependency none of the four failure modes developed earlier in this
chapter make explicit on their own.

\section{Worked Examples Across the Four Modes}

The four failure modes are illustrated across this chapter with examples
chosen for concreteness rather than uniformity of domain: the stellar
reconstruction of \cref{ex:banik} demonstrates diagnostic collapse in a
statistical inference setting, the sodium-potassium pump discussed in the
introduction demonstrates that continuous repair can constitute the
maintained state itself rather than a response to its failure, the
central-pattern-generator hierarchy of \cref{ex:cpg} demonstrates that
repair deficit is generally a property of a coordinate within a
distributed operator ecology rather than of a single missing controller,
and the cardiac synchronization example of \cref{ex:pacemaker}
demonstrates that continuation failure can occur with every local
component fully functional. The scale-dependent example of
\cref{ex:posani-scale} adds a fifth lesson orthogonal to the other four:
restorability boundaries are not scale-invariant, and a system's
apparent position relative to one can reverse depending on the
descriptive resolution at which repair is attempted. Taken together,
these examples are intended to make the taxonomy feel like an empirical
regularity across domains rather than a formal device native only to
reconstruction theory.

\section{Open Problems for This Chapter}

Several matters remain unresolved in this skeleton and are flagged here
rather than resolved prematurely. The relationship between repair deficit
and the operator-ecology material has been sharpened through the
reachability formulation of \cref{def:reach}, but still needs full
integration with that essay's keystone-operator vocabulary rather than a
side reference. The sheaf-theoretic material currently housed in the
companion Expiatory Gap essay's \S20 needs a decision about relocation to
\cref{subsec:continuation-failure} now that continuation failure has its
own worked example and no longer needs to borrow that essay's treatment
by default. And while each failure mode now has at least one concrete
non-computational instance, a fully worked civilizational or institutional
example --- as opposed to a biological one --- for repair deficit and
continuation failure would strengthen the chapter's claim to generality
beyond the inferential and physiological domains it currently draws on.

\chapter{Continuation}
% ------------------------------------------------------------------

\section{Introduction: A Short Chapter, Predicted in Advance}

The Admissibility chapter closed by stating that a chapter on
continuation would not need new foundational machinery, only trajectories
formalized properly, its pointwise-reduction theorem restated, and a
comparison between geometric repair and operator ecologies. This chapter
does those three things and one more: it resolves an apparent conflict
between two sources this corpus draws on, which the Admissibility
chapter's brevity did not have room to address.

\section{Trajectories and Coordination Geometries}
\label{sec:trajectories}

\begin{definition}[Continuation trajectory]
\label{def:trajectory}
A continuation trajectory is a pair sequence $(x(t), \mathcal{G}(t))$,
where $x(t)$ is a state in the sense of earlier chapters and
$\mathcal{G}(t)$ is a coordination geometry supplying a CLIO projection
$\pi_{\mathcal{G}(t)}$ at each $t$, following the construction in
\emph{Coordination Geometry}.
\end{definition}

\begin{claim}[Pointwise reduction]
\label{thm:pointwise-recall}
For fixed $\mathcal{G}(t)$ at each instant, ``$x(t)$ is
continuation-admissible'' is equivalent to $S_{\mathrm{RSVP}}(\pi_{\mathcal{G}(t)}(x(t)))
\leq 0$, the entropy-monotonicity admissibility condition of the
Admissibility chapter, evaluated at that instant. This introduces no
structure beyond that chapter's static relation.
\end{claim}

\Cref{thm:pointwise-recall} is stated here exactly as it was in the
Admissibility chapter, for a reason that becomes important in
\cref{sec:checking-authoring}: it is a claim about checking admissibility
at a \emph{fixed} $\mathcal{G}(t)$, and says nothing yet about how
$\mathcal{G}(t)$ itself changes.

\section{Checking Versus Authoring}
\label{sec:checking-authoring}

\emph{Histories Become Operators} makes a claim that, read against
\cref{thm:pointwise-recall} without care, looks like a direct
contradiction: admissibility is not a passive condition fixed in advance
by a trace, but something ``the continuer partly authors in the act of
continuing.'' Reading is admissibility-\emph{constituting}, on this
account, not admissibility-\emph{revealing}. If \cref{thm:pointwise-recall}
already reduces continuation-admissibility to checking a static
inequality, what is left for a continuer to author?

\begin{claim}[Resolution: checking and authoring operate at different levels]
\label{claim:resolution}
\cref{thm:pointwise-recall} is a claim about the relation between $x(t)$
and $\mathcal{G}(t)$ at a single, frozen instant: given $\mathcal{G}(t)$,
whether $x(t)$ is admissible is a fact to be checked, not authored. It is
silent about how $\mathcal{G}(t+1)$ is obtained from $\mathcal{G}(t)$,
$x(t)$, and the accumulated history $H_t$. That transition --- geometric
repair, in \emph{Coordination Geometry}'s own terms, triggered whenever
$S_{\mathrm{RSVP}}(\pi_{\mathcal{G}(t)}(x(t))) > 0$ --- is exactly where
\emph{Histories Become Operators}' authorship claim applies. Checking is
pointwise and static; authoring is the reorganization of the geometry
doing the checking, and is neither pointwise nor static.
\end{claim}

The apparent contradiction dissolves once ``continuation-admissibility''
is recognized as ambiguous between two questions: whether the present
state satisfies the present geometry's condition (checking, settled by
\cref{thm:pointwise-recall}), and how the geometry itself came to be the
one now doing the checking (authoring, not addressed by that theorem at
all). Both sources are right about the question each was actually
answering.

\section{Geometric Repair as Level-3 Modification}
\label{sec:level-3}

Chapter 6 already gave a name to exactly this kind of transition.

\begin{claim}[Geometric repair is a Level-3 modification]
\label{claim:geometric-is-level3}
Recall Chapter 6's three-level adaptation hierarchy: Level~1 state update
$x_{t+1} = F(x_t)$, Level~2 parameter update $x_{t+1} = F_{\theta_t}(x_t)$,
and Level~3 rule update $F_{t+1} = G(F_t, x_t, E_t)$, in which the
mapping governing future behavior is itself the object modified.
Geometric repair, $\mathcal{G}(t) \to \mathcal{G}(t+1)$ triggered by an
admissibility violation, is a Level~3 modification: $\mathcal{G}(t)$ plays
the role of $F_t$, and the trigger condition plays the role of $E_t$.
\end{claim}

This is not a new observation about geometric repair so much as a
recognition that Chapter 6's hierarchy already had room for it. It also
sharpens what ``authoring'' means in \cref{claim:resolution}: authorship
is Level~3 activity, in a sense already made precise, not a vague
appeal to agency or interpretation.

\section{Geometric Repair and Operator-Ecology Regeneration}
\label{sec:comparison}

Chapter 6 defined a regenerative operator as one that increases an
ecology's future capacity to produce repairs, rather than merely
repairing a single structure. Geometric repair looks, on its face, like
exactly this kind of operation applied to coordination geometries instead
of operator ecologies. Whether it actually is requires the same
discipline this book has applied to every other resemblance of this
shape.

\begin{remark}[Family resemblance, not identity]
\label{rem:geometric-vs-regenerative}
Chapter 6's regenerative operator $g$ acts on an operator ecology
$\mathcal{E}$, with the payoff measured in $\Reach(g(\mathcal{E}))$ or
induced reachability volume. Geometric repair acts on a coordination
geometry $\mathcal{G}$, with the payoff measured by whether
$S_{\mathrm{RSVP}}(\pi_{\mathcal{G}'}(x)) \leq 0$ is restored. These fail
the definition-by-substitution test this book applies before treating
any two same-shaped constructions as identical: passing a type test ---
both are operators that act on the acting-structure rather than on the
state directly, which is a genuine structural parallel --- is not the
same as passing a substitution test, and no map has been constructed
showing a coordination geometry is a special case of an operator
ecology, or the reverse. They are cousins under Chapter 6's Level~3
category, not shown to be the same construction wearing different
notation.
\end{remark}

\subsection{The Induced Ecology, Type-Corrected}
\label{sec:induced-ecology}

An admissible reprojection $\pi_\mathcal{G} : X \to X'$ is a map between
spaces, in the sense of Chapter 1's projection and embedding operations.
It is not, on its own, an element of an operator repertoire in Chapter
6's sense: the operators composed in $\mathcal{R}_\mathcal{E}(x) =
\bigcup_n \mathcal{R}_\mathcal{E}^{(n)}(x)$ must be self-maps on one
fixed space, since each stage of a composition $o_k \circ \cdots \circ
o_1(x)$ has to land back inside the domain of the next operator. Taking
$\{\pi_\mathcal{G}\}$ itself as the repertoire of an induced ecology
therefore fails a type test before any proposition about it can be
stated.

What is genuinely Level~3 in \cref{def:levels}'s sense is not
$\pi_\mathcal{G}$ but $\mathcal{G}$: $\pi_\mathcal{G}$ is the fixed rule
checked against at an instant, and the transition $\mathcal{G}(t) \to
\mathcal{G}(t+1)$, triggered by an admissibility violation, is the
object being modified. The induced ecology is accordingly built one
level up, on the augmented space $Y = X \times \mathfrak{G}$ of
state--geometry pairs, with the geometry-repair transitions themselves
as its operators.

\begin{definition}[Induced coordination ecology]
\label{def:coordination-ecology}
For a coordination geometry $\mathcal{G}$ with candidate repair
transitions $\rho : \mathfrak{G} \to \mathfrak{G}$ available to it, the
\emph{induced ecology} is
\[
\mathcal{E}_\mathcal{G} \;=\; \bigl\{\, \rho : \mathfrak{G} \to
\mathfrak{G} \bigm| \rho \text{ a candidate geometry-repair transition
available given } \mathcal{G}\text{'s current constraint network} \,\bigr\},
\]
acting on $Y = X \times \mathfrak{G}$ by $\rho(x,\mathcal{G}) =
(x,\rho(\mathcal{G}))$: a repair operator reconfigures which
$\pi_\mathcal{G}$ is doing the checking and leaves $x$ untouched.
\end{definition}

Since $\mathcal{R}^{(0)}_{\mathcal{E}_\mathcal{G}}(x,\mathcal{G}) =
\{(x,\mathcal{G})\}$ by definition, the identity (no-repair) case is
already present in $\mathcal{E}_\mathcal{G}$ without needing to be
posited as a separate hypothesis.

\subsection{The Regeneration Proposition}

\begin{proposition}[Geometric repair induces an ecology]
\label{prop:geometric-regenerative}
Let $\mathcal{G}$ be a coordination geometry whose candidate repair
transitions are closed under composition. Then $\mathcal{E}_\mathcal{G}$
of \cref{def:coordination-ecology} is an operator ecology in the sense
of Chapter 6. A transition $\rho$ satisfying \cref{thm:pointwise-reduction}
--- restoring $S_{\mathrm{RSVP}}(\pi_{\mathcal{G}(t)}(x(t))) \leq 0$ at
the current instant --- is \emph{reachability-restoring} in Chapter 6's
sense. It is \emph{regenerative} in $\mathcal{E}_\mathcal{G}$ if and
only if it additionally enlarges the future repertoire itself,
$\Reach(\rho(\mathcal{E}_\mathcal{G})) \supsetneq \Reach(\mathcal{E}_\mathcal{G})$,
rather than merely restoring the present-instant admissibility verdict.
\end{proposition}

The ``if and only if'' in \cref{prop:geometric-regenerative} is the
open content, not a settled equivalence: nothing yet characterizes,
independently of the definition itself, which reachability-restoring
transitions satisfy it. A transition can restore
$S_{\mathrm{RSVP}} \leq 0$ at the present instant while leaving
$\Reach(\mathcal{E}_\mathcal{G})$ unchanged --- a repair without
regeneration, exactly the gap this section develops an invariant for
rather than assumes away.

\begin{remark}[The converse is not automatic]
\label{rem:converse-obstruction}
Given an operator ecology $\mathcal{E}$, constructing a coordination
geometry realizing it would require every operator to be expressible as
an admissible reprojection of a common space, with ecological
composition corresponding to geometric composition. Chapter 6's own
Wittgenstein example is a candidate obstruction: the commissioned
left-hand repertoire is discrete, symbolic, and built by commission
rather than by continuous reprojection of a shared state space, with no
evident $\pi_\mathcal{G}$ family generating it by composition. The
expected result is a correspondence between a restricted class of
coordination geometries and a geometrically realizable class of operator
ecologies, not a full equivalence.
\end{remark}

\subsection{The Reachability Gain Vector}
\label{sec:gain-vector}

Chapter 6's effective repertoire dimension $D_{\mathrm{eff}}(\mathcal{E},x)$
is built from local effect vectors $v_o(x) \in T_xX$, one per operator.
A geometry-repair transition $\rho$ does not act on $x$ at all, so it
has no natural home in $T_xX$; the fix is to stop representing an
operator's effect as a tangent vector and represent the \emph{reachable
set it changes} as a vector instead, which is well-typed for every
operator, state-level or geometry-level alike.

\begin{definition}[Reachability gain vector]
\label{def:gain-vector}
Fix a measure space $(X,\mu)$. For an operator $o$ and state $x$, let
$\chi_S \in L^2(X,\mu)$ denote the indicator function of a measurable
set $S \subseteq X$. The \emph{reachability gain vector} of $o$ at $x$ is
\[
v_o(x) \;=\; \frac{\chi_{\mathcal{R}_{\mathcal{E}\cup\{o\}}(x)} -
\chi_{\mathcal{R}_{\mathcal{E}\setminus\{o\}}(x)}}
{\bigl\| \chi_{\mathcal{R}_{\mathcal{E}\cup\{o\}}(x)} -
\chi_{\mathcal{R}_{\mathcal{E}\setminus\{o\}}(x)} \bigr\|}.
\]
For a geometry-repair transition $\rho$, $\mathcal{R}$ is replaced by
$\Gamma$: $v_\rho(x)$ is the normalized indicator of
$\Gamma_{\rho(d)}(x) \setminus \Gamma_d(x)$.
\end{definition}

\begin{proposition}[Consistency with reachability-restoration]
\label{prop:gain-consistency}
$v_\rho(x)$ is well defined and nonzero exactly when $\rho$ is
reachability-restoring in Chapter 6's sense.
\end{proposition}

\begin{proof}
$v_\rho(x)$ is defined by normalizing $\chi_{\Gamma_{\rho(d)}(x)
\setminus \Gamma_d(x)}$, which is nonzero exactly when
$\Gamma_{\rho(d)}(x) \setminus \Gamma_d(x)$ has positive measure, i.e.\
when $\mu(\Gamma_{\rho(d)}(x)) > \mu(\Gamma_d(x))$ under the monotonicity
assumption $\Gamma_d(x) \subseteq \Gamma_{\rho(d)}(x)$ adopted for this
construction. This is exactly Chapter 6's reachability-restoring
condition.
\end{proof}

\Cref{prop:gain-consistency} depends on treating repair as monotone:
$\Gamma_d(x) \subseteq \Gamma_{\rho(d)}(x)$, so the signed difference
$\chi_{\Gamma_{\rho(d)}(x)} - \chi_{\Gamma_d(x)}$ is nonnegative and the
gain region is unambiguous. Repairs that both gain and lose reachable
states are not covered by \cref{def:gain-vector} as stated; they are
addressed in \cref{sec:signed}.

\begin{remark}[Recovering Chapter 6's tangent-vector construction]
In the smooth setting Chapter 6's original $v_o(x) \in T_xX$ was built
for, applying $o$ for small time $\varepsilon$ moves $x$ along $v_o(x)$
to first order, so the newly reachable region under $o$ is a shrinking
sliver in direction $v_o(x)$; rescaled, its indicator converges weakly
to a point mass along $v_o(x)$ as $\varepsilon \to 0$. \cref{def:gain-vector}
is accordingly a generalization of Chapter 6's construction, not a
replacement for it --- the original recovers as a limiting case.
Rigorously establishing that limit is a short lemma not carried out
here.
\end{remark}

\begin{claim}[Bounds theorem transfers]
\label{claim:bounds-transfer}
Define $C_{\mathcal{E}}(x) = \sum_o w_o\, v_o(x) v_o(x)^\top$ as an
operator on $L^2(X,\mu)$, with $D_{\mathrm{eff}}(\mathcal{E},x) =
(\sum_i \lambda_i)^2 / \sum_i \lambda_i^2$ over its eigenvalues. Then
$1 \leq D_{\mathrm{eff}}(\mathcal{E},x) \leq \operatorname{rank}
C_\mathcal{E}(x)$, exactly as in the finite-dimensional case.
\end{claim}

\begin{proof}
The proof of Chapter 6's bounds theorem uses only Cauchy--Schwarz
applied to the finitely many nonzero eigenvalues of $C_\mathcal{E}(x)$;
it requires finite rank, not finite ambient dimension. Since
$\operatorname{rank} C_\mathcal{E}(x)$ is bounded by the (finite) number
of contributing operators, the argument goes through verbatim on
$L^2(X,\mu)$.
\end{proof}

\subsection{Signed Extension: Non-Monotone Repair}
\label{sec:signed}

Dropping the monotonicity assumption of \cref{sec:gain-vector}
altogether, a repair $\rho$ acting on damage $d$ at $x$ generally both
gains and loses reachable states:
\[
\mathrm{gain}_\rho = \Gamma_{\rho(d)}(x) \setminus \Gamma_d(x), \qquad
\mathrm{loss}_\rho = \Gamma_d(x) \setminus \Gamma_{\rho(d)}(x),
\]
the positive and negative parts of the signed difference $\delta\chi_\rho
= \chi_{\Gamma_{\rho(d)}(x)} - \chi_{\Gamma_d(x)}$, an ordinary Jordan
decomposition and nothing further.

\begin{definition}[Unnormalized gain and loss vectors]
$g_\rho(x) = \chi_{\mathrm{gain}_\rho}$, $l_\rho(x) =
\chi_{\mathrm{loss}_\rho}$, left unnormalized so that $\|g_\rho\|^2 =
\mu(\mathrm{gain}_\rho)$ and $\|l_\rho\|^2 = \mu(\mathrm{loss}_\rho)$
carry magnitude rather than being discarded by normalization.
\end{definition}

\begin{definition}[Signed reachability covariance]
\label{def:signed-covariance}
\[
C(x) \;=\; \sum_\rho w_\rho \bigl( g_\rho(x)g_\rho(x)^\top -
l_\rho(x)l_\rho(x)^\top \bigr).
\]
\end{definition}

$C(x)$ is Hermitian but not in general positive semidefinite. Its trace,
\[
\operatorname{tr} C(x) \;=\; \sum_\rho w_\rho \bigl(
\mu(\mathrm{gain}_\rho) - \mu(\mathrm{loss}_\rho) \bigr),
\]
is the weighted sum of every candidate repair's net reachability
change, aggregating what Chapter 6's reachability-restoring condition
checks one repair at a time.

\begin{definition}[Net effective dimension]
\label{def:deff-net}
$D_{\mathrm{eff}}^{\mathrm{net}}(\mathcal{E},x) = (\operatorname{tr}
C(x))^2 / \operatorname{tr}(C(x)^2)$.
\end{definition}

\begin{claim}[Extended bounds]
\label{claim:deff-net-bounds}
$0 \leq D_{\mathrm{eff}}^{\mathrm{net}}(\mathcal{E},x) \leq
\operatorname{rank} C(x)$. The upper bound holds unconditionally; the
lower bound of $1$ from \cref{claim:bounds-transfer} need not hold once
$C(x)$ is indefinite, and $D_{\mathrm{eff}}^{\mathrm{net}} = 0$ exactly
when $\operatorname{tr} C(x) = 0$.
\end{claim}

\begin{proof}
The Cauchy--Schwarz step $(\sum_i \lambda_i)^2 \leq r \sum_i
\lambda_i^2$, for $r = \operatorname{rank} C(x)$, holds for any real
$\lambda_i$ regardless of sign, giving the upper bound unchanged. The
lower bound $(\sum_i \lambda_i)^2 \geq \sum_i \lambda_i^2$ used in the
positive-semidefinite case relied on every cross term $2\lambda_i
\lambda_j$ being nonnegative; with mixed signs this can fail, and the
inequality can be violated down to $\operatorname{tr} C(x) = 0$, at
which point $D_{\mathrm{eff}}^{\mathrm{net}}$ is $0$ by definition of
the ratio at a zero numerator.
\end{proof}

Pure-gain repair ($l_\rho \equiv 0$ for every $\rho$) recovers
\cref{claim:bounds-transfer} exactly as the special case
$D_{\mathrm{eff}}^{\mathrm{net}} = D_{\mathrm{eff}}$; nothing in the
pure-gain theorem required revision, only extension around it.

\begin{definition}[Strongly regenerative]
\label{def:strongly-regenerative}
A geometry-repair transition $\rho$ is \emph{strongly regenerative} in
$\mathcal{E}_\mathcal{G}$ if adding it strictly increases
$D_{\mathrm{eff}}^{\mathrm{net}}(\mathcal{E}_\mathcal{G},x)$, as
distinct from \emph{weakly regenerative} (\cref{prop:geometric-regenerative}'s
bare condition $\Reach(\rho(\mathcal{E}_\mathcal{G})) \supsetneq
\Reach(\mathcal{E}_\mathcal{G})$).
\end{definition}

A transition that trades one reachable region for a differently shaped
but equal-measure region is weakly regenerative --- the repertoire
count increases --- while contributing zero net gain, and is not
strongly regenerative. This is the same divergence Chapter 6 already
names between $K(o)$ and $K_V(o;x)$: inventory impact counts lost
executable operators, while generative impact measures the future
region their loss disconnects, recurring here inside the regeneration
criterion itself rather than only in keystone-impact accounting.

\begin{openproblem}
\label{prob:weights-recur}
The weights $w_\rho$ in \cref{def:signed-covariance} are estimates of
the same kind as Appendix D's retention-cost weights and Chapter 9's
allocation-problem weights, not a new estimation gap specific to this
construction. No general estimation method is supplied here, consistent
with the corresponding open problems in those chapters.
\end{openproblem}

\begin{openproblem}
\label{prob:cross-terms}
$\operatorname{tr}(C(x)^2)$ mixes gain and loss vectors across distinct
operators $\rho, \rho'$ through cross terms $g_\rho \cdot g_{\rho'}$,
$l_\rho \cdot l_{\rho'}$, and $g_\rho \cdot l_{\rho'}$. By the shape of
the construction this should discount correlated repairs --- two
candidate reprojections gaining overlapping regions --- the same way
Chapter 6's original $D_{\mathrm{eff}}$ discounts collinear operator
effects, but this has not been separately verified here.
\end{openproblem}

\begin{openproblem}[Marginal criterion]
\label{prob:marginal}
\cref{def:strongly-regenerative} is stated as a whole-ecology
diagnostic. A clean per-operator criterion would instead compare
$D_{\mathrm{eff}}^{\mathrm{net}}(\mathcal{E}_\mathcal{G} \cup \{\rho\},x)$
against $D_{\mathrm{eff}}^{\mathrm{net}}(\mathcal{E}_\mathcal{G},x)$
directly as a marginal quantity attributable to $\rho$ alone, in the
manner of Chapter 6's keystone impact $K(o)$. This marginal form is not
developed here.
\end{openproblem}

\section*{What This Chapter Adds}

Continuation, as this chapter treats it, is neither a new foundational
layer beneath admissibility nor a mere restatement of it. It is the
setting in which two true but previously separate claims about
admissibility --- that it is checked, and that it is authored --- turn
out to describe different levels of the same structure rather than
competing theories. That distinction, and the identification of
authorship with Level~3 modification already defined in Chapter 6, is
this chapter's actual content. Whether geometric repair and regenerative
operators are the same construction remains open, and is flagged as such
rather than resolved by the convenience of the resemblance.

\chapter{Reconstruction}
% ------------------------------------------------------------------

\section{Introduction: The Question This Chapter Actually Asks}

It is tempting to treat reconstruction as this book's concluding
technical convenience: distinctions were made in Chapter 1, information
and entropy quantified their consequences, repair and continuation kept
structures viable, and reconstruction is simply the operation that
recovers a prior state from what remains. This chapter argues that
framing understates the actual difficulty by exactly the amount that
matters. Recovery is not always a matter of degree, cheap at one end and
expensive at the other. In the generic case, a state's successor does
not encode which of several possible predecessors actually produced it,
and no quantity of computational resource closes that gap, because the
gap is not a resource-bounded search problem in the first place.

\section{Two Kinds of Materialization}
\label{sec:two-kinds-recall}

\begin{definition}[Generating input]
\label{def:generating-input}
Given a structure $X$ produced by a process $P$, a generating input for
$X$ is a state $s$, independent of the system's own execution history,
from which $P$ can produce $X$ without reference to any intermediate
state of the system that produced the original instance of $X$.
\end{definition}

\begin{definition}[Fiber]
\label{def:fiber}
For a transition function $F : X \times U \to X$ and a state $y \in X$,
the fiber of $y$, restricted to admissible predecessor pairs $P
\subseteq X \times U$, is $\mathrm{Fib}_P(y) = \{(x,u) \in P : F(x,u) =
y\}$. $F$ is non-injective over $P$ at $y$ when $|\mathrm{Fib}_P(y)| > 1$.
\end{definition}

\begin{definition}[Reconstructive irreversibility]
\label{def:irreversibility}
A transition into $y$ is reconstructively irreversible over
$\mathrm{Fib}_P(y)$ when $F$ is non-injective there and no generating
input exists from which the predecessor's identity could otherwise be
recovered.
\end{definition}

The distinction this chapter needs is not a spectrum running from cheap
reconstruction to expensive reconstruction to impossible reconstruction.
It is categorical.

\begin{claim}[Non-reducibility]
\label{claim:non-reducibility}
No monotonic transformation of cost, however extreme, converts a
reconstruction-cost materialization into a reconstructive-irreversibility
materialization or the reverse. The two are distinguished by the
existence of a generating input and by whether a retained history
channel is injective on the relevant fiber, neither of which is a cost.
\end{claim}

A chaotic Hamiltonian system illustrates why this distinction is not
merely definitional bookkeeping: such a system can require reconstruction
precision growing exponentially in $t$ to recover an earlier state to any
fixed tolerance, making expensive reconstruction entirely rational, while
never once crossing into the regime \cref{def:irreversibility} describes
--- because a Hamiltonian flow's fibers are always singletons
(\cref{cor:hamiltonian-recall} below), so a generating input always
exists in principle, however hard it is to use. Expense and the existence
of a generating input are independent properties.

\begin{claim}[Hamiltonian dynamics never collapse fibers]
\label{cor:hamiltonian-recall}
For a Hamiltonian system with a flow map $\Phi_t$ smooth enough to
guarantee existence and uniqueness of solutions, $\Phi_t$ is a
diffeomorphism for each $t$: bijective, with smooth inverse. Every fiber
under $\Phi_t$ is therefore a singleton, and $\Phi_{-t}(y)$ is always a
generating input for the predecessor state.
\end{claim}

This is the sharpest available contrast case: a system in which
reconstruction is, by construction, never reconstructively irreversible,
however expensive it may be to actually carry out. Most systems this
book has been concerned with --- repair algebras, operator ecologies,
diagnostic ecologies under throughput pressure --- are not Hamiltonian in
this sense. Their transitions are generically non-injective, and it is
worth being precise about what that costs them.

\section{The Resource Independence Theorem}
\label{sec:resource-independence}

\begin{claim}[Resource independence]
\label{thm:resource-independence}
Let $\mathrm{Fib}_P(y)$ be subject to reconstructive irreversibility.
Then for any candidate reconstruction procedure and any allocation of
time, space, or energy to it, no such allocation enables recovery of the
predecessor's identity.
\end{claim}

\begin{proof}
By \cref{def:irreversibility}, reconstructive irreversibility holds
exactly when no generating input exists for the predecessor's identity.
A reconstruction procedure, however much time, space, or energy it is
allocated, is still a function from some input to an output; enlarging
its resource bound changes which functions are computable within that
bound, but does not change whether an input exists to compute from.
Since no state or artifact independent of the system's own execution
history determines the predecessor's identity, there is no input, of any
kind, from which any procedure --- regardless of allocated resources ---
could recover it. The claim does not depend on complexity-theoretic
bounds; it depends only on the absence of a domain element to compute
from, which no resource allocation supplies.
\end{proof}

Time, space, and energy are, within limits, fungible with one another:
time-space tradeoffs and energy-time tradeoffs are ordinary subjects of
complexity theory and physical computation respectively. Continuation is
not a fourth item on that list. \cref{thm:resource-independence} is the
precise sense in which it fails to be fungible with any of the three:
more time does not help, more space does not help, more energy does not
help, not because the problem is merely very hard, but because there is
no input for additional resources to be spent on.

\begin{claim}[Continuation is not a quantity of memory]
\label{prop:not-memory}
Possessing unbounded memory does not by itself guarantee that a system's
retained structure is injective over any given admissibility-relevant
fiber. Injectivity is a property of \emph{what} is stored --- of the
retention map applied at the moment of collapse --- not of \emph{how
much} storage is available to store it in.
\end{claim}

A system with arbitrarily large memory can still apply a retention
policy that discards exactly the information distinguishing two
predecessors --- recording that an edit occurred without recording what
the edit was --- regardless of unused capacity. This is why more storage
is not the fix for a design that retains the wrong thing: the fix is
retaining a different thing, which is a design decision, not a resource
allocation.

\section{The Distinction Budget}
\label{sec:distinction-budget}

\begin{definition}[Distinguishability preservation]
\label{def:distinguishability}
For a history channel update $G(h_t, x, u)$ and a fiber
$\mathrm{Fib}_P(y)$, the augmented dynamics preserve distinguishability
over the fiber at $h_t$ when $(x,u) \mapsto G(h_t,x,u)$ is injective on
$\mathrm{Fib}_P(y)$.
\end{definition}

\begin{definition}[Distinction budget]
\label{def:distinction-budget}
Given a system's trajectory up to time $t$, its distinction budget $B_t$
is the set of collapsed fibers over which predecessor identity remains
recoverable --- either because a generating input still exists, or
because the retained history channel is injective over that fiber.
\end{definition}

\begin{claim}[Monotonic shrinkage]
\label{prop:monotonic-shrinkage}
Absent a positive retention action at the moment a fiber collapses,
$B_t$ can only lose members as $t$ increases, never regain them.
\end{claim}

\begin{proof}
A fiber leaves $B_t$ exactly when neither a generating input nor an
injective retained channel is available for it. Both conditions concern
facts fixed at or before the moment of collapse, and neither can be
altered afterward: a generating input that has ceased to exist cannot be
un-destroyed, and a channel that failed to be injective at the moment of
retention cannot become injective retroactively, since its value was
already fixed at that moment.
\end{proof}

\begin{remark}[A structural resemblance, not a special case]
\label{rem:dpi-resemblance}
\cref{prop:monotonic-shrinkage} has the flavor of the data processing
inequality from information theory --- no processing of a signal can
increase the mutual information it carries about its source. The
resemblance is structural, not formal: this chapter's setting concerns
discrete distinguishability over admissible predecessors, not mutual
information over a channel, and no claim is made that one result is a
special case of the other. What the two share is the one-directional
shape: a process can only destroy distinguishing information already
present, never manufacture it after the fact.
\end{remark}

This gives reconstruction, as this book treats it, a genuinely different
character from repair. Chapter 6 and Chapter 7 both concern operators
that restore admissibility to a damaged state, and Chapter 6's
Anti-Alignment result showed retained history can be more or less
adequate to that task. Nothing in either chapter's machinery is
monotonic in the strong sense $B_t$ is: a damaged operator ecology can
regain reachable futures through regeneration, but a fiber that has left
$B_t$ is gone in a stronger sense --- no operator, however constructed,
restores it, because \cref{thm:resource-independence} rules out recovery
by resource expenditure and \cref{prop:monotonic-shrinkage} rules out
recovery by waiting.

\section{Historical Observability}
\label{sec:historical-observability}

\begin{definition}[Historical observability]
\label{def:historical-observability}
Given a class $Q$ of possible future admissibility judgments a system
might need to make, its historical observability with respect to $Q$ at
time $t$ is the proportion of $Q$ whose judgments remain decidable given
the current distinction budget $B_t$.
\end{definition}

Historical observability is not a fixed property of hardware or raw
storage capacity. By \cref{thm:resource-independence}, it cannot be
increased after the fact by any allocation of time, space, or energy once
a relevant fiber has left $B_t$; it can only be preserved, at the moment
of collapse, by retaining the right thing there. This is the sense in
which continuation is the resource that determines which of a system's
future admissibility judgments remain answerable at all, independent of
how much of the other three resources the system is given. A system can
have unlimited time, space, and energy and still be structurally unable
to answer a question whose answer depended on a distinction its own
dynamics discarded before anyone thought to ask it.

\section{Admissibility-Relevant Fibers and Selective Retention}
\label{sec:selective-retention}

Not every collapsed fiber is worth defending against
\cref{prop:monotonic-shrinkage}'s shrinkage. A system that tried to keep
every non-injective transition it ever executed inside $B_t$ would
exhaust its retention budget long before the retention bought it
anything.

\begin{definition}[Admissibility-relevant fiber]
\label{def:admissibility-relevant}
A fiber $\mathrm{Fib}_P(y)$ is admissibility-relevant if there exists
some later reachable state $x_j$ and continuation family $C$ such that
the admissibility of $x_j$ relative to $C$ differs depending on which
member of $\mathrm{Fib}_P(y)$ actually occurred.
\end{definition}

\begin{claim}[Selective Retention]
\label{thm:selective-retention}
A system that wishes to preserve exactly the admissibility judgments its
continuation families support, at minimum retention cost, must allocate
injective history channels precisely over admissibility-relevant fibers,
and need not allocate them elsewhere.
\end{claim}

\Cref{thm:selective-retention} has two failure modes, not one.
Under-retention leaves an admissibility-relevant fiber undefended, so
some future judgment becomes undecidable --- typically not visibly, since
systems under-provisioned this way tend to resolve the ambiguity silently
by defaulting to whichever fiber member is most recent or most common,
rather than by raising a visible error. Over-retention spends retention
budget on fibers no downstream judgment ever consults. Chapter 5's
reinterpretation of Anti-Alignment as a monotonicity lemma inside a
retention-cost optimization, rather than a standalone argument for
maximal retention, is this claim's consequence stated in that chapter's
own vocabulary: the optimum is a minimal admissibility-sufficient trace,
not the fullest trace affordable.

\section{The Allocation Problem and a Design Principle}
\label{sec:allocation}

\begin{definition}[Retention budget and candidate fibers]
\label{def:allocation-problem}
Given candidate fibers $F_1,\ldots,F_n$ with relevance weights $w_i > 0$
and retention costs $c_i > 0$, and a total retention budget $\beta$, the
allocation problem is to choose $S \subseteq \{1,\ldots,n\}$ maximizing
$\sum_{i \in S} w_i$ subject to $\sum_{i \in S} c_i \leq \beta$.
\end{definition}

This is an instance of the weighted knapsack problem, known to be
NP-hard in general. That is stated plainly rather than smoothed over:
\cref{def:allocation-problem} gives the allocation decision a precise
combinatorial shape it did not have before, but a precise shape is not
the same as a tractable one.

\begin{principle}[Admissibility-Aware Externalization]
\label{princ:admissibility-aware}
Given finite retention capacity, a system should allocate history
channels to admissibility-relevant fibers in order of relevance per unit
retention cost, $w_i / c_i$, rather than uniformly, rather than in order
of raw reconstruction difficulty, and rather than by implementation
convenience.
\end{principle}

This is stated as a design principle, not a theorem: its force depends
on the accuracy of weights and costs a designer supplies, which can be
estimated but not guaranteed correct in advance of deployment. A designer
following it imperfectly will produce exactly \cref{thm:selective-retention}'s
two failure modes --- under-retention where relevance was underestimated,
over-retention where it was overestimated --- both detectable after the
fact by audit even when neither was foreseeable before it.

\begin{openproblem}
\label{prob:allocation-open}
Four gaps remain open. The greedy heuristic's approximation gap for
retention-allocation problems specifically, as opposed to knapsack
problems in general, has not been established. The weights and costs
\cref{def:allocation-problem} requires are estimates in any real system,
and no general estimation method is supplied beyond the audit apparatus
gestured at above. Whether a system's space of possible fibers is fixed
in advance or expands as evidence arrives bears directly on whether
candidate fibers can even be enumerated at design time. And a broader
conjecture --- that any persistent system under finite retention
capacity, shaped by selective pressure toward correct future
functioning, will exhibit retention correlating with
admissibility-relevance rather than raw reconstruction cost, with
departures observable as the two failure signatures above --- has been
illustrated with motivating candidates (legal and financial audit
trails, version control's near-unconditional retention explained by
near-zero marginal cost rather than uniform relevance, and biological
memory consolidation's bias toward salient events) but not tested by
any systematic survey. It is stated in a form intended to be testable by
existing methods, not one requiring new theoretical machinery, but the
survey itself remains undone.
\end{openproblem}

\section*{What This Chapter Adds to the Book's Spine}

Reconstruction, in the sense this chapter has developed it, is not the
operation that undoes damage; that is repair's job, treated in Chapter 6
and Chapter 7. Reconstruction is the prior question of what remains
available to be acted on at all --- and this chapter's central result is
that the answer is not always \emph{everything, at a price}. Some
distinctions are gone in a stronger sense than expensive: no
allocation of time, space, or energy, however large, brings them back,
because the resource that would need to be spent was never the kind of
thing money buys. What a system retains at the moment a fiber collapses
is the only chance it gets.

\chapter*{Conclusion}
\addcontentsline{toc}{chapter}{Conclusion}

This volume set out to build a proof spine, not a topic list. It is
worth being precise about which of those two things the preceding eight
chapters actually delivered. A topic list would have eight independent
essays, each competent on its own terms, connected by nothing stronger
than shared vocabulary. What follows instead is a chain of
results that genuinely depend on one another: Chapter 2's information
measures are checked against Chapter 1's deficit apparatus by
substitution, not analogy; Chapter 5's entropy inventory reduces three
of five candidate quantities to one underlying partition structure and
leaves the other two --- repair entropy and the RSVP field --- with
different, explicitly stated epistemic status; Chapter 6's repair theory
supplies the generator-closure map Chapter 5's admissible future region
needed and never had; Chapter 7's four failure modes are shown
independent of the architecture gate rather than subordinate to it, on
the strength of explicit counterexamples rather than assumption; Chapter
7 resolves an apparent contradiction between two sources this corpus
draws on by showing they describe different levels of the same
structure; and Chapter 9 closes the spine with a result --- resource
independence --- that could not have been stated without Chapters 1
through 6 already having built the distinction/fiber vocabulary it
depends on.

Three findings recur across chapters rather than staying local to the
one that introduced them, and are worth naming together here rather than
only where each first appeared. The \textbf{Repair--Recovery Separation
Principle} (Chapter 7) --- that restoring admissibility and preserving
the historical identity of what was wrong are independent achievements
--- reappears in Chapter 6's minimality-driven operators and Chapter
10's fiber-collapsing repair transitions as three instances of one
phenomenon, not three unrelated observations that happen to share a
conclusion. The distinction between \textbf{checking and authoring}
admissibility (Chapter 8) turned out to dissolve an apparent
contradiction between two independently-developed accounts by showing
each was answering a different question at a different level, a pattern
the methodology used throughout --- type test, then substitution test, before
any two same-named constructions are treated as identical --- was built
specifically to catch. And the book's \textbf{single central open
problem}, whether the two working notions of admissibility developed in
this corpus induce the same continuation family, was sharpened rather
than closed: from a badly-typed comparison of two unrelated-looking
quantities into a well-typed question about two set-valued constructions,
strengthened by a formal argument for skepticism borrowed from an
unrelated result about constants of motion, and left open throughout on
the explicit grounds that a book already this dependent on the answer
has every incentive to resolve it prematurely, which is exactly why it
should not.

What is absent is also worth stating plainly. A
chapter originally planned under the title ``Constraint'' was found, on
direct inspection of its presumed source material, to have no technical
content distinct from Admissibility, and the gap was left rather than
filled with something that did not belong there. Whether Coordination
deserves dedicated treatment beyond what Chapters 5 and 7 already give it
remains undecided. And the relationship among the three
reachability-flavored constructions in this volume --- the executable repertoire, the
operator-induced reachable set, and the admissible future region --- is
stated as a generator chain in Chapter 6, but the pairwise relationships
among effective reachability, physical reachability, and admissibility
proper were each established under different and non-interchangeable
assumptions, a distinction easily lost on a careless reading.

The discipline that produced all of this was not, in the end, a
technique for proving things. It was a standing refusal to let two
constructions share a name, a formula, or a chapter without first asking
whether they were actually the same thing. Most of the time in this
volume, they were not.






\chapter*{Appendix A: Notation and Symbol Index}
\addcontentsline{toc}{chapter}{Appendix A: Notation and Symbol Index}

This index collects every symbol used recurrently in this volume, in
order of first appearance, each with a one-line description and a
pointer to its formal definition. It is meant to be consulted rather
than read straight through.

\section*{Distinction and Information}

\begin{center}
\begin{tabular}{lp{4.6in}}
\hline
Symbol & Meaning \\
\hline
$(X,\sim)$ & A distinguishability space: a set with an equivalence
relation of observational indistinguishability. \\
$\delta_T(x)$ & Coding deficit of $x$ under ontology $T$: excess
description length over the optimum, $L_T(x) - L^*(x)$. \\
$\Delta_T(x)$ & Distinguishability deficit of $x$ under $T$: inaccessible
distinction capacity, $D^*(x) - D_T(x)$. \\
$L_T(x)$, $L^*(x)$ & Shortest description of $x$ reachable within $T$,
and the shortest achievable by any ontology. \\
$D_T(x)$, $D^*(x)$ & Number of objects separable from $x$ within $T$,
and the maximum achievable by any ontology. \\
$\sigma, \phi, \rho, \tau$ & The four operations on distinguishability
structure: projection, embedding, revision, transport. \\
$H(\mathcal{P})$ & Shannon entropy of a partition, $-\sum_i p_i \log p_i$. \\
$S_\pi(m)$ & Representational entropy at $m$ under a representation
$\pi$: $\log|\pi^{-1}(m)|$, the log-size of the collapsed fiber. \\
$N_D$, $\Omega(P_i)$ & Number of partition cells (distinction count),
and multiplicity within a cell --- dual readings of one partition
structure. \\
\hline
\end{tabular}
\end{center}

\section*{Entropy and Reachability}

\begin{center}
\begin{tabular}{lp{4.6in}}
\hline
Symbol & Meaning \\
\hline
$S(x,t)$ & The RSVP entropy field: a monotonic quantity, $\dot{S} \geq
0$, entering the entropy-based admissibility relation. \\
$D(t)$ & Recoverable distinction capacity; the Second Law of this
volume states $\dot{D} \leq 0$. \\
$\overline{\mathcal{R}}_t(x)$ & Effective reachable set: the physically
reachable set from $x$, quotiented by the distinguishability relation at
time $t$. \\
$V_{\mathrm{eff}}(x,t)$ & Effective reachability volume, $\bar\mu_t(\overline{\mathcal{R}}_t(x))$. \\
$\Gamma_t(x)$ & Admissible future region: the effective reachable set,
read as the set of states an observer's current admissibility structure
licenses. \\
$\SR(F \mid T)$ & Repair entropy: uncertainty over fault locations
consistent with a symptom, given retained trace $T$. \\
$\mathcal{H}_{\mathrm{fault}}(F,T)$, $\Omega_F(s)$ & The live fault
hypotheses given $T$, and the full fault hypothesis space for symptom
$s$. \\
$\mathcal{P}_T$ & The partition of a hypothesis space induced by a
retained trace $T$. \\
\hline
\end{tabular}
\end{center}

\section*{Admissibility}

\begin{center}
\begin{tabular}{lp{4.6in}}
\hline
Symbol & Meaning \\
\hline
$\pi : \mathcal{M} \to \Sigma$ & A model's prediction map, with
$\operatorname{Im}(\pi)$ the observations it can in principle produce. \\
$\chi_\pi(x)$ & The architecture gate: $0$ if $x \in \operatorname{Im}(\pi)$
(inferential boundary), $1$ otherwise (ontological boundary). \\
$\delta_\pi^{\mathrm{arch}}(x)$ & Architectural representation deficit:
$\inf_{m} d_\Sigma(\pi(m), x)$, a graded quantity whose zero/nonzero
threshold reproduces $\chi_\pi$. \\
$C(x)$ & Continuation family: the set of trajectories from $x$
considered permitted. \\
$C_S(x)$, $C_\delta(x)$ & The continuation families induced by the
entropy-monotonicity and deficit-bounded admissibility constructions,
respectively --- the two objects at the center of this volume's central
open problem. \\
$\pi_{\mathcal{G}}$ & The CLIO projection supplied by a coordination
geometry $\mathcal{G}$. \\
\hline
\end{tabular}
\end{center}

\section*{Repair and Operator Ecologies}

\begin{center}
\begin{tabular}{lp{4.6in}}
\hline
Symbol & Meaning \\
\hline
$\mathcal{E}$ & An operator ecology: a constraint network determining
which operators remain executable. \\
$\Reach(\mathcal{E})$ & The reachable executable repertoire of
$\mathcal{E}$: a set of operators. \\
$\mathcal{R}_\mathcal{E}(x)$ & Operator-induced reachable set: the states
reachable from $x$ by composing operators in $\Reach(\mathcal{E})$. A
different, generated object from $\Reach(\mathcal{E})$ itself. \\
$K(o)$ & Keystone impact (inventory): $|\Reach(\mathcal{E})| -
|\Reach(\mathcal{E}\setminus\{o\})|$. \\
$K_V(o;x)$ & Generative keystone impact: the corresponding loss in
$\mu(\mathcal{R}_\mathcal{E}(x))$, which need not track $K(o)$. \\
$D_{\mathrm{eff}}(\mathcal{E},x)$ & Effective repertoire dimension: a
participation-ratio-style measure of decorrelated operator diversity at
$x$, bounded between $1$ and $\operatorname{rank}(C_\mathcal{E}(x))$. \\
\hline
\end{tabular}
\end{center}

\section*{Restorability and History}

\begin{center}
\begin{tabular}{lp{4.6in}}
\hline
Symbol & Meaning \\
\hline
$\Rboundary$, $\Rboundary[_t]$ & The restorability boundary, and its
history-indexed form $\Rboundary(X,\mathcal{E},\mathcal{R},H_t)$. \\
$H_t$ & Repair history: the remembered sequence $(\varepsilon_0, R_1,
\varepsilon_1, \ldots, R_n, \varepsilon_n)$ of residuals and attempted
repairs up to time $t$. \\
$\operatorname{Orb}_\mathcal{R}(\varepsilon_0)$ & Repair orbit: the set
of residuals reachable from $\varepsilon_0$ by any word in the repair
algebra $\mathcal{R}$. \\
$P(\varepsilon_0,\mathcal{R})$ & Persistent residual: $\inf_{w}
\|w(\varepsilon_0)\|$ over the full repair monoid. \\
$\widehat{P}_t$ & Empirical persistence bound: the same infimum, taken
only over words actually attempted by time $t$; an upper bound on $P$. \\
\hline
\end{tabular}
\end{center}

\section*{Reconstruction}

\begin{center}
\begin{tabular}{lp{4.6in}}
\hline
Symbol & Meaning \\
\hline
$F : X \times U \to X$ & A transition function; $\mathrm{Fib}_P(y)$
its fiber over $y$, restricted to admissible predecessor pairs $P$. \\
$B_t$ & Distinction budget: the set of collapsed fibers over which
predecessor identity remains recoverable at time $t$. Monotonically
non-increasing absent a retention decision at the moment of collapse. \\
\hline
\end{tabular}
\end{center}

\chapter*{Appendix B: Quotients, Fibers, and Projections}
\addcontentsline{toc}{chapter}{Appendix B: Quotients, Fibers, and Projections}

Every chapter of this volume depends, at some point, on the same basic
construction: a map that collapses some states together, and the
question of what survives that collapse. This appendix develops that
construction once, with enough generality to support its uses in
Chapters 1, 2, 6, and 10 without repetition.

\section*{Equivalence Relations and Quotients}

\begin{definition}[Quotient by an equivalence relation]
For a set $X$ and an equivalence relation $\sim$ on $X$, the quotient
$X/{\sim}$ is the set of equivalence classes $[x]_\sim = \{y \in X : y
\sim x\}$, with the quotient map $q : X \to X/{\sim}$ sending $x \mapsto
[x]_\sim$.
\end{definition}

\begin{definition}[Coarsening order]
For equivalence relations $\sim_1, \sim_2$ on $X$, $\sim_1$ is finer
than $\sim_2$ (equivalently, $\sim_2$ is coarser) when $x \sim_1 y
\implies x \sim_2 y$ for all $x,y$. This defines a partial order on the
set of equivalence relations on $X$, with the discrete relation (every
element its own class) as the finest element and the indiscrete
relation (one class) as the coarsest.
\end{definition}

A coarsening $\sim_1 \subseteq \sim_2$ induces a canonical surjection
$c_{21} : X/{\sim_1} \to X/{\sim_2}$ sending each $\sim_1$-class to the
$\sim_2$-class containing it, well defined because $\sim_1 \subseteq
\sim_2$ guarantees every $\sim_1$-class lies inside a single
$\sim_2$-class.

\section*{Fibers and Projections}

\begin{definition}[Projection and fiber]
A projection is a surjective map $\pi : X \to Y$. The fiber over $y \in
Y$ is $\pi^{-1}(y) = \{x \in X : \pi(x) = y\}$. A projection induces an
equivalence relation on $X$ by $x \sim_\pi x' \iff \pi(x) = \pi(x')$,
whose classes are exactly the fibers, so every projection is a quotient
map for the equivalence relation it induces, and conversely every
quotient map $X \to X/{\sim}$ is a projection with fibers equal to the
$\sim$-classes.
\end{definition}

\begin{claim}[Fiber monotonicity]
If $\sigma : X \to Y$ factors as $X \xrightarrow{q} X/{\sim_1}
\xrightarrow{c} Y$ for a coarsening $\sim_1 \subseteq \sim_\sigma$, then
each fiber of $\sigma$ is a union of $\sim_1$-classes, and $|\sigma^{-1}(y)|
\geq |q^{-1}([x]_{\sim_1})|$ for any $x$ with $\sigma(x) = y$. Coarser
factoring maps have fibers containing at least as much as finer ones.
\end{claim}

\begin{definition}[Restricted fiber]
For a transition function $F : X \times U \to X$ and a subset $P
\subseteq X \times U$ of admissible predecessor pairs, the restricted
fiber of $y$ is $\mathrm{Fib}_P(y) = \{(x,u) \in P : F(x,u) = y\}$. $F$
is injective over $P$ at $y$ when $|\mathrm{Fib}_P(y)| \leq 1$.
\end{definition}

This is the fiber notion used in Chapter 9: a restriction of the
ordinary fiber of $F$ to the pairs an admissibility structure has
already licensed, which is what allows a system to be non-injective in
general while remaining injective on the specific fibers that matter to
it.

\section*{Pullbacks}

\begin{definition}[Pullback]
For maps $f : X \to Z$ and $g : Y \to Z$, the pullback $X \times_Z Y =
\{(x,y) \in X \times Y : f(x) = g(y)\}$ is the set of pairs agreeing on
their image in $Z$, with projections to $X$ and $Y$ that commute with
$f$ and $g$.
\end{definition}

Pullbacks are the construction underlying the disagreement map of
Chapter 7: given local sections $s_i \in \mathcal{F}(U_i)$ over a cover
$\{U_i\}$, the overlap $U_i \cap U_j$ is where the pullback of the two
restriction maps is tested for agreement, and the sheaf gluing condition
is exactly the statement that a family agreeing on every pairwise
pullback determines a unique global section.

\section*{A Disciplined Notion of Projection}

The projections of this appendix are required only to be constant on
the equivalence classes they collapse. A stronger and useful notion
requires, in addition, that nothing \emph{beyond} those classes gets
collapsed, and that singular cases be given structured treatment rather
than excluded.

\begin{definition}[Constraint-leveraged inference operator]
\label{def:CLIO-appB}
Let $X$ be a state space, $\mathcal{A}$ a system of local admissibility
constraints on $X$ inducing an equivalence relation $\sim_\mathcal{A}$,
and $M$ a compressed invariant space. A map $\pi_\mathcal{A} : X \to M$
is a constraint-leveraged inference operator when it satisfies three
conditions: it is constant on $\sim_\mathcal{A}$-classes (an ordinary
quotient map); its fibers are \emph{exactly} the $\sim_\mathcal{A}$-classes,
$\pi_\mathcal{A}(x) = \pi_\mathcal{A}(x') \iff x \sim_\mathcal{A} x'$
(obstruction faithfulness --- no accidental collapsing beyond what the
equivalence relation itself demands); and it extends to a class of
controlled singular inputs $\hat{X} \supset X$, with singularities
treated as structured morphism types carrying definite invariant data
rather than as excluded configurations (singular extension).
\end{definition}

\begin{remark}[This sharpens, rather than replaces, the ordinary quotient map]
\label{rem:CLIO-vs-quotient}
Every quotient map $q : X \to X/{\sim}$ of this appendix already
satisfies obstruction faithfulness by construction --- $q(x) = q(x')$
iff $x \sim x'$ is definitional. What \cref{def:CLIO-appB} adds is the
third condition, singular extension, which the ordinary quotient map
construction is silent about: it says nothing about how to treat inputs
where the equivalence relation itself becomes singular or ill-defined.
This is exactly the situation the Admissibility chapter's discussion of
constrained Hamiltonian systems encounters --- a singular Legendre
transform producing genuinely non-injective structure --- treated there
only by analogy to physics. \cref{def:CLIO-appB} gives that situation a
formal name and a required treatment (structured morphism types, not
exclusion) rather than leaving it as an unformalized precedent.
\end{remark}

\section*{Induced Partitions}

\begin{claim}[Composability of quotients]
For a chain of coarsenings $\sim_1 \subseteq \sim_2 \subseteq \cdots
\subseteq \sim_n$, the induced surjections compose:
$c_{n1} = c_{n,n-1} \circ \cdots \circ c_{21}$, and the composite factors
$X/{\sim_1} \to X/{\sim_n}$ through every intermediate quotient.
\end{claim}

This is the fact underlying Chapter 7's effective reachability
contraction: distinguishability coarsening over time, $\sim_{t_1}
\subseteq \sim_{t_2} \subseteq \cdots$, produces a chain of quotients
through which any reachable set's image can only lose elements as it
passes through successive coarsenings, since a surjective image can
never contain more distinct elements than its domain.

\section*{Summary Table}

\begin{center}
\begin{tabular}{lp{4.3in}}
\hline
Construction & Where it is used \\
\hline
Quotient $X/{\sim}$ & The distinguishability space of Chapter 1;
partitions underlying Chapter 2's entropy measures. \\
Fiber $\pi^{-1}(y)$ & Representational entropy $S_\pi(m)$ (Chapter 2);
the architecture gate $\chi_\pi$ (Chapter 7). \\
Restricted fiber $\mathrm{Fib}_P(y)$ & Reconstructive irreversibility
and the distinction budget $B_t$ (Chapter 9). \\
Coarsening order & The Second Law and effective reachability
contraction (Chapter 7). \\
Pullback & The disagreement map and sheaf gluing condition for
continuation failure (Chapter 7). \\
\hline
\end{tabular}
\end{center}

\chapter*{Appendix C: Continuation and Reachability Structures}
\addcontentsline{toc}{chapter}{Appendix C: Continuation and Reachability Structures}

Three closely related but formally distinct constructions recur through
the second half of this volume: an executable repertoire, a reachable
set generated from it, and an admissible future region built from that.
A separate object, the continuation family, organizes how admissibility
itself is stated. This appendix develops both in one place, together
with the relationships among them, so that later chapters can use all
four without re-deriving their connections each time.

\section*{Trajectory Space and Continuation Families}

\begin{definition}[Trajectory space]
For a state space $X$, let $\mathcal{T}$ denote the space of
trajectories: sequences (or, in continuous time, paths) $\tau :
[0,\infty) \to X$ or $\tau : \mathbb{N} \to X$ with $\tau(0) = x$ for
some $x \in X$, the trajectory's basepoint. Write $\tau|_{\geq t}$ for
the suffix of $\tau$ from time $t$ onward, itself a trajectory with
basepoint $\tau(t)$.
\end{definition}

\begin{definition}[Continuation family]
\label{def:continuation-family-formal}
A continuation family is a map $C : X \to \mathcal{P}(\mathcal{T})$
assigning to each $x \in X$ a set $C(x)$ of trajectories with basepoint
$x$, considered permitted. A state $x_i$ occurring within a trajectory
$\tau$ is admissible relative to $\tau$ and $C$ when the suffix of
$\tau$ beginning at $x_i$ is a member of $C(x_i)$.
\end{definition}

Two continuation families organize the admissibility discussion of
Chapter 3: $C_S(x) = \{\tau : \dot{S}(\tau) \geq 0\}$, built from entropy
monotonicity, and $C_\delta(x) = \{\tau : \delta_T(\tau) \leq
\epsilon\}$, built from a distortion bound. Neither closure property
below is automatic for either.

\begin{definition}[Closure under restriction]
$C$ is closed under restriction if $\tau \in C(x)$ implies $\tau|_{\geq
t} \in C(\tau(t))$ for every $t$ in $\tau$'s domain.
\end{definition}

\begin{definition}[Closure under concatenation]
$C$ is closed under concatenation if, whenever $\tau_1 \in C(x)$ has
endpoint $y$ and $\tau_2 \in C(y)$, the concatenation $\tau_1 \cdot
\tau_2 \in C(x)$.
\end{definition}

\begin{remark}[Neither closure property is free]
$C_S$ is closed under restriction --- a suffix of a non-decreasing path
is itself non-decreasing --- but its closure under concatenation depends
on continuity of $S$ at the join point, a substantive assumption rather
than a consequence of $\dot{S} \geq 0$ on each piece. $C_\delta$ need not
be closed under either property in general: a suffix can have higher
deficit relative to its own basepoint if $T$ is state-dependent, and
concatenation can exceed a threshold $\epsilon$ if deficits are
sub-additive rather than bounded by the maximum of the pieces.
\end{remark}

\begin{openproblem}
Characterizing the conditions under which $C_S$ and $C_\delta$ are
closed under concatenation, and whether either can be modified to
guarantee it without changing its extension on single trajectories,
remains open.
\end{openproblem}

\section*{Morphisms and Equivalence}

\begin{definition}[Continuation morphism]
\label{def:continuation-morphism}
For continuation families $C_1$ on $X_1$ and $C_2$ on $X_2$, a
continuation morphism is a pair $(f,\phi)$ where $f : X_1 \to X_2$ and
$\phi$ sends each $\tau \in C_1(x)$ to a trajectory $\phi(\tau) \in
C_2(f(x))$, compatibly with restriction: $\phi(\tau|_{\geq t}) =
\phi(\tau)|_{\geq t}$.
\end{definition}

\begin{definition}[Continuation equivalence]
States $x, y \in X$ are continuation-equivalent, $x \sim_C y$, when
there exist mutually inverse continuation morphisms between $C(x)$ and
$C(y)$: that is, when $C(x) \cong C(y)$.
\end{definition}

\begin{remark}[Distinct from behavioral equivalence]
Chapter 6 develops a hierarchy of equivalence notions --- state,
structural, behavioral, regenerative --- comparing systems' capacity to
recreate future behavior after damage. $\sim_C$ compares permitted
continuation structure directly, and the two need not coincide: two
states can be continuation-equivalent while behaviorally distinct, if
their permitted trajectory sets are isomorphic without their actual
transition dynamics agreeing, or the reverse.
\end{remark}

\begin{openproblem}[The central open problem, in morphism form]
\label{prob:admissibility-morphism}
Does there exist a continuation morphism $(f,\phi) : C_\delta \to C_S$
with $f = \mathrm{id}_X$? If such a morphism exists and is invertible,
$C_\delta \cong C_S$ and the two admissibility notions coincide. If a
morphism exists in one direction only, one family refines the other. If
none exists in either direction, the two are formally incomparable ---
stronger evidence for genuine separation than an unproved equality
alone.
\end{openproblem}

\section*{Reachability Structures}

A second cluster of constructions --- an operator repertoire, a
generated reachable set, and an admissible future region --- appear
across Chapters 3, 5, and 6 and are related by generation and
restriction rather than by identity.

\begin{definition}[Executable repertoire]
For an operator ecology $\mathcal{E}$, $\Reach(\mathcal{E})$ is the set
of operators currently executable given $\mathcal{E}$'s constraint
network.
\end{definition}

\begin{definition}[Operator-induced reachable set]
For $x \in X$, $\mathcal{R}_\mathcal{E}(x) = \bigcup_{n \geq 0}
\mathcal{R}_\mathcal{E}^{(n)}(x)$, where $\mathcal{R}_\mathcal{E}^{(n)}(x)$
is the set of states reachable from $x$ by composing at most $n$
operators from $\Reach(\mathcal{E})$, every intermediate composition
executable.
\end{definition}

\begin{definition}[Admissible future region]
For a distinguishability quotient $q_t : X \to X/{\sim_t}$ and physical
reachable set $\mathcal{R}(x)$, the effective reachable set is
$\overline{\mathcal{R}}_t(x) = q_t(\mathcal{R}(x))$, and the admissible
future region is $\Gamma_t(x) = \overline{\mathcal{R}}_t(x)$.
\end{definition}

\begin{claim}[The three constructions form a generator chain]
\[
\Reach(\mathcal{E}) \;\longrightarrow\; \mathcal{R}_\mathcal{E}(x) \;\longrightarrow\; \Gamma_t(x).
\]
$\Reach(\mathcal{E})$ is a set of operators. $\mathcal{R}_\mathcal{E}(x)$
is the closure of that repertoire under composition: the states
reachable by chaining executable operators together. $\Gamma_t(x)$ is a
further restriction of a reachable set --- here,
$\overline{\mathcal{R}}_t(x)$ --- to the states an observer's current
admissibility structure licenses. Each stage is a genuine restriction or
generation of the one before it, not a synonym for it: $\Reach(\mathcal{E})$
takes operators as elements, $\mathcal{R}_\mathcal{E}(x)$ and $\Gamma_t(x)$
take states, so the first arrow fails a type test outright and cannot be
an identity; the second arrow is a quotient, not an inclusion.
\end{claim}

\begin{claim}[Repertoire--reachability monotonicity]
For extension-monotone ecologies --- $\mathcal{O}_{\mathcal{E}_1}
\subseteq \mathcal{O}_{\mathcal{E}_2}$ implies every composition
executable in $\mathcal{E}_1$ remains executable in $\mathcal{E}_2$ ---
$\Reach(\mathcal{E}_1) \subseteq \Reach(\mathcal{E}_2)$ implies
$\mathcal{R}_{\mathcal{E}_1}(x) \subseteq \mathcal{R}_{\mathcal{E}_2}(x)$
for every $x$, and consequently $V_{\mathcal{E}_1}(x) \leq
V_{\mathcal{E}_2}(x)$ whenever $\mu$ is monotone under inclusion.
\end{claim}

\begin{claim}[Effective reachability contraction]
For distinguishability coarsening $\sim_{t_1} \subseteq \sim_{t_2}$,
$\overline{\mathcal{R}}_{t_2}(x) = c_{21}(\overline{\mathcal{R}}_{t_1}(x))$
for the canonical surjection $c_{21}$ induced by the coarsening, and
$|\overline{\mathcal{R}}_{t_2}(x)| \leq |\overline{\mathcal{R}}_{t_1}(x)|$
in the finite case.
\end{claim}

\begin{remark}[Reading the diagram]
Growth in a repertoire propagates forward through both arrows: more
operators generate at least as large a reachable set, which (holding the
distinguishability quotient fixed) generates at least as large an
admissible future region. Entropy production acts on the second arrow
only, coarsening the quotient that turns a reachable set into an
admissible one, and is silent about the first: an ecology can gain
operators while an observer's distinguishability structure simultaneously
coarsens, with the two effects pulling the composite chain in opposite
directions.
\end{remark}

\chapter*{Appendix D: Persistence and Reconstruction}
\addcontentsline{toc}{chapter}{Appendix D: Persistence and Reconstruction}

Two bodies of machinery, developed separately in Chapters 6 and 10,
concern the same underlying question from different sides: given a
retained history, what can still be known, and what determines whether
retaining more of it would help. This appendix collects both, and
supplies one piece neither chapter states formally: a criterion for
exactly which fibers a retention policy needs to protect.

\section*{Residual Orbits and Persistence}

\begin{definition}[Residual orbit and persistent residual]
Let $\mathcal{R}$ be a repair algebra acting on a discrepancy space
$\mathcal{D}$. For a repair word $w = R_n \circ \cdots \circ R_1 \in
\langle \mathcal{R} \rangle$, write $\varepsilon_w = w(\varepsilon_0)$.
The repair orbit of $\varepsilon_0$ is $\operatorname{Orb}_\mathcal{R}(\varepsilon_0)
= \{w(\varepsilon_0) : w \in \langle\mathcal{R}\rangle\}$. For a
discrepancy norm $\|\cdot\|$, the persistent residual is
$P(\varepsilon_0,\mathcal{R}) = \inf_{w \in \langle\mathcal{R}\rangle}
\|w(\varepsilon_0)\|$.
\end{definition}

\begin{definition}[Empirical persistence bound]
For a history $H_t$ containing attempted repair words $w_1,\ldots,w_n$,
$\widehat{P}_t = \min_{1 \leq k \leq n} \|w_k(\varepsilon_0)\|$.
\end{definition}

\begin{claim}[Monotonicity of the empirical persistence bound]
As repair history grows, $\widehat{P}_{t+1} \leq \widehat{P}_t$, and
$P(\varepsilon_0,\mathcal{R}) \leq \widehat{P}_t$ for every finite
history.
\end{claim}

$H_t$ records only attempted words, so it can only approximate
$P(\varepsilon_0,\mathcal{R})$ from above; repeated failure does not by
itself prove structural persistence unless the attempted repair family
is sufficiently exhaustive.

\begin{claim}[Persistence identification under exhaustive search]
If the attempted repair words become dense in $\langle\mathcal{R}\rangle$
under a topology for which $w \mapsto \|w(\varepsilon_0)\|$ is
continuous, then $\widehat{P}_t \to P(\varepsilon_0,\mathcal{R})$.
\end{claim}

This distinguishes ``every repair tried so far has failed'' from ``every
repair available in the architecture must fail'' --- the second is much
stronger, and is what separates Mercury's perihelion (where the
Newtonian repair algebra was plausibly searched near-exhaustively) from
a stellar-age reconstruction's noise floor (where no comparable density
argument has been made).

\section*{Generating Inputs, Fibers, and Reconstructive Irreversibility}

\begin{definition}[Generating input and fiber]
Given a structure $X$ produced by a process $P$, a generating input for
$X$ is a state $s$, independent of the system's own execution history,
from which $P$ can produce $X$ without reference to any intermediate
state of the original instance. For a transition function $F : X \times
U \to X$ and admissible predecessor pairs $P \subseteq X \times U$, the
restricted fiber of $y$ is $\mathrm{Fib}_P(y) = \{(x,u) \in P : F(x,u) =
y\}$; $F$ is injective over $P$ at $y$ when $|\mathrm{Fib}_P(y)| \leq
1$.
\end{definition}

\begin{definition}[Reconstructive irreversibility]
A transition into $y$ is reconstructively irreversible over
$\mathrm{Fib}_P(y)$ when $F$ is non-injective there and no generating
input exists from which the predecessor's identity could otherwise be
recovered.
\end{definition}

\begin{claim}[Non-reducibility]
No monotonic transformation of cost converts a reconstruction-cost
materialization (a generating input exists; cost is the only issue) into
a reconstructive-irreversibility materialization (no generating input
exists at all), or the reverse. The two are distinguished by the
existence of a generating input, which is not a cost.
\end{claim}

\begin{claim}[Hamiltonian dynamics never collapse fibers]
For a Hamiltonian system with a flow map $\Phi_t$ smooth enough to
guarantee existence and uniqueness of solutions, $\Phi_t$ is a
diffeomorphism for each $t$, so every fiber under $\Phi_t$ is a
singleton and $\Phi_{-t}(y)$ is always a generating input.
\end{claim}

This is the sharpest available contrast case: reconstruction can be
arbitrarily expensive in a chaotic Hamiltonian system while never once
becoming reconstructively irreversible, since expense and the existence
of a generating input are independent properties.

\section*{Resource Independence}

\begin{theorem}[Resource independence]
Let $\mathrm{Fib}_P(y)$ be subject to reconstructive irreversibility.
Then for any candidate reconstruction procedure and any allocation of
time, space, or energy to it, no such allocation enables recovery of the
predecessor's identity.
\end{theorem}

The proof turns on the absence of a domain element to compute from, not
on any complexity bound: enlarging a procedure's resource allocation
changes which functions are computable within that bound, but does not
supply an input where none exists. Time, space, and energy are, within
limits, fungible with one another; resource independence is the precise
sense in which continuation fails to be fungible with any of the three.

\begin{claim}[Continuation is not a quantity of memory]
Injectivity of a retained history channel is a property of \emph{what}
is stored, not of \emph{how much} storage is available. Unbounded
memory does not by itself guarantee that a retained structure is
injective over any given fiber.
\end{claim}

\section*{The Distinction Budget}

\begin{definition}[Distinguishability preservation and distinction budget]
For a history channel update $G(h_t,x,u)$ and a fiber $\mathrm{Fib}_P(y)$,
the augmented dynamics preserve distinguishability over the fiber at
$h_t$ when $(x,u) \mapsto G(h_t,x,u)$ is injective on $\mathrm{Fib}_P(y)$.
Given a trajectory up to time $t$, the distinction budget $B_t$ is the
set of collapsed fibers over which predecessor identity remains
recoverable --- either a generating input still exists, or the retained
history channel is injective over that fiber.
\end{definition}

\begin{claim}[Monotonic shrinkage]
Absent a positive retention action at the moment a fiber collapses,
$B_t$ can only lose members as $t$ increases, never regain them.
\end{claim}

A fiber leaves $B_t$ exactly when neither a generating input nor an
injective retained channel is available for it, and both conditions are
fixed at or before the moment of collapse: nothing afterward can restore
either.

\begin{definition}[Historical observability]
Given a class $Q$ of possible future admissibility judgments, historical
observability with respect to $Q$ at time $t$ is the proportion of $Q$
whose judgments remain decidable given $B_t$.
\end{definition}

By resource independence, historical observability cannot be increased
after the fact by any allocation of time, space, or energy once a
relevant fiber has left $B_t$; it can only be preserved, at the moment
of collapse, by retaining the right thing there.

\section*{Selective Retention}

The preceding sections establish that some fibers are worth defending
against monotonic shrinkage and others are not, without yet saying
which. This section supplies the criterion.

\begin{definition}[Admissibility-relevant fiber]
\label{def:admissibility-relevant-appD}
A fiber $\mathrm{Fib}_P(y)$ is admissibility-relevant if there exists
some later state $x_j$, reachable from $y$, and some continuation family
$C$, such that the admissibility of $x_j$ relative to $C$ differs
depending on which member of $\mathrm{Fib}_P(y)$ actually occurred.
\end{definition}

\begin{example}
If two distinct reasoning traces --- one that correctly ruled out an
unsafe action, and one that merely ran out of budget before considering
it --- are compressed into the same summarized state, and a later step
re-proposes the same action, whether that proposal is admissible depends
on which predecessor occurred: the fiber is admissibility-relevant. If
two different low-level operations happen to produce identical log
entries, and nothing downstream ever depends on which specific operation
produced the entry, the corresponding fiber is collapsed but not
admissibility-relevant.
\end{example}

\begin{theorem}[Selective Retention]
\label{thm:selective-retention-appD}
A system that wishes to preserve exactly the admissibility judgments its
continuation families support, at minimum retention cost, must allocate
injective history channels precisely over admissibility-relevant fibers,
and need not allocate them elsewhere.
\end{theorem}

\begin{proof}[Sketch]
Necessity: if a fiber is admissibility-relevant, some future state's
admissibility verdict depends on which predecessor occurred. A
non-injective history channel over that fiber leaves the predecessor's
identity unrecoverable, so the verdict cannot be determined. Sufficiency
and minimality: by \cref{def:admissibility-relevant-appD}, no future
admissibility judgment depends on distinguishing predecessors within a
non-admissibility-relevant fiber, so injective retention there preserves
a distinction with no consumer, at nonzero cost, and can be dropped
without any judgment becoming unsupportable.
\end{proof}

\begin{corollary}[Two failure modes]
A retention architecture fails \cref{thm:selective-retention-appD} in
exactly two ways: under-retention, leaving an admissibility-relevant
fiber undefended so that some future judgment becomes undecidable; or
over-retention, maintaining injective channels over fibers that are not
admissibility-relevant, paying a cost no future judgment consumes.
\end{corollary}

Under-retention is invisible until the specific future state that needed
the distinction actually arrives and cannot be evaluated. Over-retention
is visible directly, as wasted storage or compute, which is why naive
economic accounts of retention notice and penalize it while missing
under-retention entirely.

\begin{remark}[Consequence for Chapter 5's retention-cost objective]
Chapter 5 poses the objective $J(T) = \SR(F \mid T) + \lambda K(T)$
without specifying what the retention cost functional $K(T)$ should
actually penalize. \cref{thm:selective-retention-appD} answers this: $K$
should be understood as counting injective retention maintained over
fibers that are not admissibility-relevant, and the minimal
admissibility-sufficient trace $T^*$ that objective seeks is exactly the
trace achieving injectivity over the admissibility-relevant fibers of
\cref{def:admissibility-relevant-appD} and nowhere else.
\end{remark}

\section*{Summary}

Reconstruction, across both chapters this appendix draws together, is
not the operation that undoes damage --- that is repair's province. It
is the prior question of what remains available to be acted on at all.
The persistent residual $P(\varepsilon_0,\mathcal{R})$ asks whether a
discrepancy would survive unlimited repair attempts; the distinction
budget $B_t$ asks whether a predecessor's identity would survive
unlimited resource expenditure. Both answers can be no, for reasons that
have nothing to do with how much repair or resource is available.
Selective Retention closes the gap between the two: it says which
fibers a system must protect from the second failure to guarantee it
never faces a version of the first that could have been prevented.

\chapter*{Appendix E: Worked Examples}
\addcontentsline{toc}{chapter}{Appendix E: Worked Examples}

The preceding chapters illustrate their formal claims with examples
drawn from astrophysics, neurophysiology, and cortical population
recordings. This appendix develops one extended example instead,
returned to across three escalating tiers, because a single running
example makes the relationships between chapters --- which nothing else
in this book states in one place --- visible in a way scattered
one-chapter illustrations cannot.

\section*{Setup: A Lane of Vehicles}

Consider a lane of vehicles indexed $i = 1, 2, \ldots$, each with state
$x_i \in X$ (position, velocity) evolving under local sensing rather
than centralized control. No vehicle has access to a global state; each
has access only to its own sensors and, where extended below, a small
number of communication channels to neighbors and infrastructure.

\section*{Tier 1: Stigmergic Local Control}

Each vehicle does not track its absolute position or speed as the
quantity governing safe following. It tracks the gap $g_i(t)$ to the
vehicle ahead and the closing rate $\dot{g}_i(t)$: whether that gap is
shrinking or growing. This is stigmergy in the classical sense ---
coordination through local, indirect signals in the shared environment,
the mechanism behind ant pheromone trails and Reynolds flocking rules
--- rather than through any shared global representation.

\begin{definition}[Pairwise admissibility]
For adjacent vehicles $i, i+1$ with joint state $(x_i, x_{i+1}) \in X
\times X$, a continuation family on the product space, $C : X \times X
\to \mathcal{P}(\mathcal{T})$, determines whether a joint trajectory
maintains $g_i(t) \geq g_{\min}$ for a safety threshold $g_{\min}$.
\end{definition}

This is an ordinary application of the continuation family of Appendix
C, not a generalization of it: admissibility here is a property of the
pair, but the pair itself is simply a state in the product space $X
\times X$, and \cref{def:continuation-family-formal} already covers
states in any space, product or otherwise. No new formal machinery is
needed for a fixed pair.

\begin{remark}[Where this framing breaks]
The product-space framing depends on each vehicle having exactly one
relevant neighbor. Merging lanes, variable following distance, or a
vehicle with two adjacent neighbors on either side all break the fixed
pair into a variable-size neighborhood, at which point $X \times X$ is
no longer the right domain and something closer to a local system ---
an assignment of continuation families varying coherently across
overlapping neighborhoods, in the spirit of the sheaf-theoretic
machinery Chapter 7 uses for continuation failure --- would be needed
instead. This is not developed here.
\end{remark}

The safety condition $g_i(t) \geq g_{\min}$ does not require any vehicle
to categorize its situation into discrete named states --- congested,
free-flowing, merging. It requires only that the local geometry of gap
and closing rate remain separable enough to support one admissibility
judgment. This is \cref{rem:posani-slogan}'s claim in a different
domain: restorability depends on separability, not categoricity. A
vehicle does not need a clean category for its situation any more than
a cortical neuron needs one for the variable it encodes.

Correction under this scheme is also distributed in exactly the sense
\cref{ex:cpg} develops: no single controller adjusts the whole lane, and
a local perturbation --- one vehicle braking --- is absorbed by
neighboring vehicles' local responses without any global recomputation,
the repair reserve $R(d,x)$ of \cref{def:reserve} realized as
neighboring vehicles' independent capacity to respond, not as a single
central controller's bandwidth.

\section*{Tier 2: Diagnostic Escalation and Geometric Repair}

Vehicles in this system have no steering wheel. Passengers hold tablets;
drones are stationed along the road; a central console displays a
ranked list of pending questions; towers with human pilots can assume
direct control during accidents or extreme ambiguity. Four distinct
mechanisms operate across this hierarchy, each realizing a different
piece of earlier chapters.

\subsection*{Ambiguous objects and diagnostic collapse}

When an onboard camera cannot classify an object with confidence, its
diagnostic channel is approaching $\I(F;E) \to 0$ locally
(\cref{subsec:diagnostic-collapse}). The system's response is not to
apply more processing to the same camera feed; it routes the question to
a passenger's tablet, a structurally different diagnostic channel.
Chapter 7 states the general remedy for diagnostic collapse as
requiring a qualitatively new channel, not more of the same one; this is
that remedy with a concrete instantiation rather than an abstract one.

\subsection*{Drone control as geometric repair}

A passenger who takes control of a local drone ahead does not correct
their own vehicle's trajectory. They modify the shared coordination
geometry other vehicles are responding to --- exactly geometric repair
in the sense of \cref{thm:pointwise-reduction}, and a Level~3
modification in \cref{def:levels}'s hierarchy: the object being changed
is the rule governing future behavior, not a state or a parameter within
a fixed rule.

\begin{remark}[Checking versus authoring, instantiated]
A camera's classification check against the current coordination
geometry is checking, in \cref{claim:resolution}'s sense: evaluated
against a geometry held fixed at that instant. A passenger's decision to
take control of a drone, reorganizing what the geometry permits going
forward, is authoring. The two happen at different levels of the same
system, exactly as Chapter 8 argues they must.
\end{remark}

\section*{Tier 3: The Salience Screen}

The central console does not surface every ambiguity uniformly. It
ranks pending questions by how much they matter, which is precisely
Admissibility-Aware Externalization (\cref{princ:admissibility-aware}):
allocate attention to admissibility-relevant fibers
(\cref{def:admissibility-relevant-appD}) in order of relevance per unit
cost of driver attention, not by raw ambiguity or uniform coverage. A
question about an object that no downstream decision depends on --- the
non-admissibility-relevant case of Appendix D's own example --- does not
belong on the screen even if the classifier's confidence is low; a
question that determines whether a merge is safe belongs there even if
the classifier is otherwise fairly confident.

\section*{Tier 4: Tower Pilots and the Architecture Gate}

Tower pilots are invoked far less often than passenger tablet queries,
and their intervention is qualitatively different from every tier below.
The condition that triggers it --- ``accidents or extremely ambiguous
situations'' --- is doing real work: it distinguishes an inferential
boundary, where more asking within the existing sensor-and-passenger
architecture would eventually resolve the ambiguity, from an ontological
boundary, where the architecture itself cannot represent what has
happened and a different control modality is required regardless of how
much more is asked of it. This is the architecture gate
$\chi_\pi$ of \cref{def:gate}, with the same two-sided character as
Mercury's perihelion versus Banik's stellar-age reconstruction
(\cref{ex:mercury}): most ambiguous situations are inferential,
resolved by tiers 1 through 3; a genuine accident can be ontological,
outside what the vehicle's own sensors and passengers' judgment can
represent at all.

\begin{remark}[A keystone at the top of the hierarchy]
Tower-pilot intervention is rare and high-leverage in exactly
\cref{rem:keystone-divergence}'s sense: its low invocation frequency
does not indicate low importance. Its unavailability would disconnect a
disproportionate share of the system's correctable state space ---
exactly the situations tiers 1 through 3 cannot resolve --- which is a
generative keystone impact $K_V$ far larger than its usage frequency
alone would suggest, the same asymmetry apprenticeship training carries
in the two-cities argument of Chapter 6.
\end{remark}

\section*{Summary Table}

\begin{center}
\begin{tabular}{lp{3.6in}}
\hline
Tier & Formal object instantiated \\
\hline
Gap and closing rate & Pairwise admissibility on a product-space
continuation family; separability over categoricity. \\
Passenger tablet query & A new diagnostic channel resolving diagnostic
collapse, not more capacity on the same channel. \\
Drone control & Geometric repair; a Level-3 modification; the authoring
side of the checking/authoring distinction. \\
Salience screen & Admissibility-Aware Externalization: ranking by
relevance per unit cost, not raw ambiguity. \\
Tower pilot & The architecture gate: inferential versus ontological
boundary, with keystone-scale generative impact despite rare
invocation. \\
\hline
\end{tabular}
\end{center}

Each tier of this single system realizes a construction introduced
independently, in a different chapter, for a different reason. That
convergence is not itself a proof of anything; it is offered as evidence
that the formal vocabulary developed across this volume describes a
recurring structural pattern rather than a set of unrelated devices each
suited only to the domain that motivated it.

\chapter*{References}
\addcontentsline{toc}{chapter}{References}

\begin{thebibliography}{99}

\bibitem{banik2026}
Banik, I., Kudakolawa Kaluarachchige, T., Cookson, S., \& Desmond, H.
(2026). The age of the Universe from a large sample of the oldest
Galactic stars. \textit{Monthly Notices of the Royal Astronomical
Society}, in press. arXiv:2607.00764.

\bibitem{posani2026}
Posani, L., Wang, S., Muscinelli, S.\ P., Paninski, L., \& Fusi, S.
(2026). Rarely categorical, highly separable representations along the
cortical hierarchy. \textit{Nature}. https://doi.org/10.1038/s41586-026-10668-4

\bibitem{flyxion-distinguishability}
Flyxion. \textit{Distinguishability Geometry}. Companion essay.

\bibitem{flyxion-admissibility-field}
Flyxion. \textit{The Admissibility Field}. Companion essay.

\bibitem{flyxion-constraint}
Flyxion. \textit{Constraint Before Content}. Companion monograph.

\bibitem{flyxion-coordination}
Flyxion. \textit{Coordination Geometry: From Synchrony to Structure}.
Companion essay.

\bibitem{flyxion-operator-ecology}
Flyxion. \textit{Operator Ecology}. Companion essay.

\bibitem{flyxion-history-before-function}
Flyxion. \textit{History Before Function: Compression as the Origin of
Operators}. Companion essay.

\bibitem{flyxion-recursive-continuation}
Flyxion. \textit{Recursive Continuation: Autopoiesis, Sympoiesis, and
the Compression of Self-Modification}. Companion essay.

\bibitem{flyxion-histories-become-operators}
Flyxion. \textit{Histories Become Operators: Biological Notes on
Retained Admissibility}. Companion essay.

\bibitem{flyxion-ecology-of-distinctions}
Flyxion. \textit{The Ecology of Distinctions}. Companion monograph.

\bibitem{flyxion-persistent-anomalies}
Flyxion. \textit{Persistent Anomalies and the Geometry of Ontology
Revision: Repair, Admissibility, and Non-Abelian Tears}. Companion
monograph.

\bibitem{flyxion-hydra}
Flyxion. \textit{What Survives Reconstruction}. Companion monograph.

\bibitem{flyxion-expiatory-gap}
Flyxion. \textit{The Expiatory Gap: Civilizational Incompressibility and
the Limits of Superintelligent Optimization}. Companion essay.

\bibitem{flyxion-generative-continuation}
Flyxion. \textit{Generative Continuation Geometry}. Companion essay.

\end{thebibliography}
\end{document}
